% Môn: Vật lý
% Lớp: 11
% Chương: Chương I. Dao động
% Bài: L11C1 Bài 5. Động năng. Thế năng. Sự chuyển hóa giữa động năng và thế năng trong dao động điều hòa
% Số câu: 320
% App đọc file này trực tiếp — sửa rồi Commit trên GitHub, bấm Đồng bộ GitHub .tex

% L11C1 Bài 5. Động năng. Thế năng. Sự chuyển hóa giữa động năng và thế năng trong dao động điều hòa

% ===== Câu 1 =====
% ID: L11C1B5-01-DS
% Mức: TH
\dangbt{Chu kì biến thiên của động năng, thế năng}
\begin{ex}
% ID: L11C1B5-01-DS
% Mức: TH
Một vật dao động điều hòa quanh vị trí cân bằng theo phương trình li độ $x = A\cos(\omega t + \varphi)$. Xét tính đúng/sai của các nhận định sau đây về sự biến thiên năng lượng của vật trong quá trình dao động:
\choiceTF
{Động năng của vật đạt giá trị cực đại khi vật đi qua một trong hai vị trí biên.}
{\True Thế năng của vật đạt giá trị cực đại khi vật ở vị trí biên dương hoặc biên âm.}
{\True Cơ năng của vật dao động điều hòa luôn được bảo toàn và không thay đổi theo thời gian nếu bỏ qua mọi lực cản.}
{\True Khi vật di chuyển từ vị trí biên về vị trí cân bằng, thế năng đàn hồi hoặc thế năng trọng trường giảm dần và chuyển hóa thành động năng.}
\loigiai{
\begin{itemchoice}
\item (Sai) Động năng của vật đạt giá trị cực đại khi vật đi qua một trong hai vị trí biên. (Vì): Mệnh đề này sai vì khi vật qua vị trí biên thì tốc độ của vật bằng $0$, do đó động năng của vật đạt giá trị tối thiểu bằng $0$
\item (Đúng) Thế năng của vật đạt giá trị cực đại khi vật ở vị trí biên dương hoặc biên âm. (Vì): Mệnh đề này đúng vì tại các vị trí biên $x = \pm A$, độ lệch của vật khỏi vị trí cân bằng đạt cực đại, kéo theo thế năng đạt giá trị cực đại
\item (Đúng) Cơ năng của vật dao động điều hòa luôn được bảo toàn và không thay đổi theo thời gian nếu bỏ qua mọi lực cản. (Vì): Mệnh đề này đúng vì khi hệ không có ma sát hay lực cản, cơ năng biến thiên tuần hoàn của động năng và thế năng luôn có tổng là hằng số bảo toàn
\item (Đúng) Khi vật di chuyển từ vị trí biên về vị trí cân bằng, thế năng đàn hồi hoặc thế năng trọng trường giảm dần và chuyển hóa thành động năng. (Vì): Mệnh đề này đúng vì khi vật đi từ biên về vị trí cân bằng, li độ giảm (thế năng giảm) và tốc độ tăng (động năng tăng), thế năng chuyển hóa dần thành động năng
\end{itemchoice}
}
\end{ex}

% ===== Câu 2 =====
% ID: L11C1B5-01-TL
% Mức: TH
\begin{bt}
% ID: L11C1B5-01-TL
% Mức: TH
Con lắc lò xo có $k=50$ N/m, biên độ $4\,\text{cm}$. Tính cơ năng; động năng và thế năng tại $x=2$ cm.
\end{bt}

% ===== Câu 3 =====
% ID: L11C1B5-01-TLN
% Mức: TH
\begin{ex}
% ID: L11C1B5-01-TLN
% Mức: TH
Con lắc lò xo có độ cứng $50\,\text{N}$ /m và biên độ $4\,\text{cm}$. Cơ năng theo đơn vị $J$ bằng bao nhiêu?
\shortans{0,04}
\end{ex}

% ===== Câu 4 =====
% ID: L11C1B5-01-TN
% Mức: NB
\begin{ex}
% ID: L11C1B5-01-TN
% Mức: NB
Một vật dao động điều hòa với chu kì $T$. Động năng và thế năng của vật biến thiên tuần hoàn theo thời gian với chu kì $T'$ bằng:
\choice
{$2T$}
{$T$}
{\True $\dfrac{T}{2}$}
{$\dfrac{T}{4}$}
\loigiai{
• Biểu thức thế năng theo thời gian:
 $W_{\text{t}} = \dfrac{1}{2}m\omega^2 A^2 \cos^2(\omega t + \varphi) = \dfrac{1}{4}m\omega^2 A^2 \left[1 + \cos(2\omega t + 2\varphi)\right].$
 
• Tần số góc biến thiên của năng lượng là $\omega' = 2\omega$.
 
• Chu kì biến thiên tuần hoàn của năng lượng:
 $T' = \dfrac{2\pi}{\omega'} = \dfrac{2\pi}{2\omega} = \dfrac{T}{2}.$
}
\end{ex}

% ===== Câu 5 =====
% ID: L11C1B5-02-DS
% Mức: TH
\begin{ex}
% ID: L11C1B5-02-DS
% Mức: TH
Một vật dao động điều hòa quanh vị trí cân bằng theo phương trình li độ $x = A\cos(\omega t + \varphi)$. Xét tính đúng/sai của các nhận định sau đây về sự biến thiên năng lượng của vật trong quá trình dao động:
\choiceTF
{Động năng của vật đạt giá trị cực đại khi vật đi qua một trong hai vị trí biên.}
{\True Thế năng của vật đạt giá trị cực đại khi vật ở vị trí biên dương hoặc biên âm.}
{\True Cơ năng của vật dao động điều hòa luôn được bảo toàn và không thay đổi theo thời gian nếu bỏ qua mọi lực cản.}
{\True Khi vật di chuyển từ vị trí biên về vị trí cân bằng, thế năng đàn hồi hoặc thế năng trọng trường giảm dần và chuyển hóa thành động năng.}
\loigiai{
\begin{itemchoice}
\item (Sai) Động năng của vật đạt giá trị cực đại khi vật đi qua một trong hai vị trí biên. (Vì): Mệnh đề này sai vì khi vật qua vị trí biên thì tốc độ của vật bằng $0$, do đó động năng của vật đạt giá trị tối thiểu bằng $0$
\item (Đúng) Thế năng của vật đạt giá trị cực đại khi vật ở vị trí biên dương hoặc biên âm. (Vì): Mệnh đề này đúng vì tại các vị trí biên $x = \pm A$, độ lệch của vật khỏi vị trí cân bằng đạt cực đại, kéo theo thế năng đạt giá trị cực đại
\item (Đúng) Cơ năng của vật dao động điều hòa luôn được bảo toàn và không thay đổi theo thời gian nếu bỏ qua mọi lực cản. (Vì): Mệnh đề này đúng vì khi hệ không có ma sát hay lực cản, cơ năng biến thiên tuần hoàn của động năng và thế năng luôn có tổng là hằng số bảo toàn
\item (Đúng) Khi vật di chuyển từ vị trí biên về vị trí cân bằng, thế năng đàn hồi hoặc thế năng trọng trường giảm dần và chuyển hóa thành động năng. (Vì): Mệnh đề này đúng vì khi vật đi từ biên về vị trí cân bằng, li độ giảm (thế năng giảm) và tốc độ tăng (động năng tăng), thế năng chuyển hóa dần thành động năng
\end{itemchoice}
}
\end{ex}

% ===== Câu 6 =====
% ID: L11C1B5-02-TL
% Mức: VD
\begin{bt}
% ID: L11C1B5-02-TL
% Mức: VD
Con lắc lò xo có $k=80$ N/m, biên độ $6\,\text{cm}$. Tính cơ năng; động năng và thế năng tại $x=3$ cm.
\end{bt}

% ===== Câu 7 =====
% ID: L11C1B5-02-TLN
% Mức: VD
\begin{ex}
% ID: L11C1B5-02-TLN
% Mức: VD
Con lắc lò xo có độ cứng $80\,\text{N}$ /m và biên độ $6\,\text{cm}$. Cơ năng theo đơn vị $J$ bằng bao nhiêu?
\shortans{0,144}
\end{ex}

% ===== Câu 8 =====
% ID: L11C1B5-02-TN
% Mức: NB
\begin{ex}
% ID: L11C1B5-02-TN
% Mức: NB
Một vật dao động điều hòa với chu kì $T$. Động năng và thế năng của vật biến thiên tuần hoàn theo thời gian với chu kì $T'$ bằng:
\choice
{$2T$}
{$T$}
{\True $\dfrac{T}{2}$}
{$\dfrac{T}{4}$}
\loigiai{
• Biểu thức thế năng theo thời gian:
 $W_{\text{t}} = \dfrac{1}{2}m\omega^2 A^2 \cos^2(\omega t + \varphi) = \dfrac{1}{4}m\omega^2 A^2 \left[1 + \cos(2\omega t + 2\varphi)\right].$
 
• Tần số góc biến thiên của năng lượng là $\omega' = 2\omega$.
 
• Chu kì biến thiên tuần hoàn của năng lượng:
 $T' = \dfrac{2\pi}{\omega'} = \dfrac{2\pi}{2\omega} = \dfrac{T}{2}.$
}
\end{ex}

% ===== Câu 9 =====
% ID: L11C1B5-03-DS
% Mức: TH
\begin{ex}
% ID: L11C1B5-03-DS
% Mức: TH
Một vật dao động điều hòa quanh vị trí cân bằng theo phương trình li độ $x = A\cos(\omega t + \varphi)$. Xét tính đúng/sai của các nhận định sau đây về sự biến thiên năng lượng của vật trong quá trình dao động:
\choiceTF
{Động năng của vật đạt giá trị cực đại khi vật đi qua một trong hai vị trí biên.}
{\True Thế năng của vật đạt giá trị cực đại khi vật ở vị trí biên dương hoặc biên âm.}
{\True Cơ năng của vật dao động điều hòa luôn được bảo toàn và không thay đổi theo thời gian nếu bỏ qua mọi lực cản.}
{\True Khi vật di chuyển từ vị trí biên về vị trí cân bằng, thế năng đàn hồi hoặc thế năng trọng trường giảm dần và chuyển hóa thành động năng.}
\loigiai{
\begin{itemchoice}
\item (Sai) Động năng của vật đạt giá trị cực đại khi vật đi qua một trong hai vị trí biên. (Vì): Mệnh đề này sai vì khi vật qua vị trí biên thì tốc độ của vật bằng $0$, do đó động năng của vật đạt giá trị tối thiểu bằng $0$
\item (Đúng) Thế năng của vật đạt giá trị cực đại khi vật ở vị trí biên dương hoặc biên âm. (Vì): Mệnh đề này đúng vì tại các vị trí biên $x = \pm A$, độ lệch của vật khỏi vị trí cân bằng đạt cực đại, kéo theo thế năng đạt giá trị cực đại
\item (Đúng) Cơ năng của vật dao động điều hòa luôn được bảo toàn và không thay đổi theo thời gian nếu bỏ qua mọi lực cản. (Vì): Mệnh đề này đúng vì khi hệ không có ma sát hay lực cản, cơ năng biến thiên tuần hoàn của động năng và thế năng luôn có tổng là hằng số bảo toàn
\item (Đúng) Khi vật di chuyển từ vị trí biên về vị trí cân bằng, thế năng đàn hồi hoặc thế năng trọng trường giảm dần và chuyển hóa thành động năng. (Vì): Mệnh đề này đúng vì khi vật đi từ biên về vị trí cân bằng, li độ giảm (thế năng giảm) và tốc độ tăng (động năng tăng), thế năng chuyển hóa dần thành động năng
\end{itemchoice}
}
\end{ex}

% ===== Câu 10 =====
% ID: L11C1B5-03-TL
% Mức: VD
\begin{bt}
% ID: L11C1B5-03-TL
% Mức: VD
Con lắc lò xo có $k=100$ N/m, biên độ $8\,\text{cm}$. Tính cơ năng; động năng và thế năng tại $x=4$ cm.
\end{bt}

% ===== Câu 11 =====
% ID: L11C1B5-03-TLN
% Mức: VD
\begin{ex}
% ID: L11C1B5-03-TLN
% Mức: VD
Con lắc lò xo có độ cứng $100\,\text{N}$ /m và biên độ $8\,\text{cm}$. Cơ năng theo đơn vị $J$ bằng bao nhiêu?
\shortans{0,32}
\end{ex}

% ===== Câu 12 =====
% ID: L11C1B5-03-TN
% Mức: NB
\begin{ex}
% ID: L11C1B5-03-TN
% Mức: NB
Một vật dao động điều hòa với chu kì $T$. Động năng và thế năng của vật biến thiên tuần hoàn theo thời gian với chu kì $T'$ bằng:
\choice
{$2T$}
{$T$}
{\True $\dfrac{T}{2}$}
{$\dfrac{T}{4}$}
\loigiai{
• Biểu thức thế năng theo thời gian:
 $W_{\text{t}} = \dfrac{1}{2}m\omega^2 A^2 \cos^2(\omega t + \varphi) = \dfrac{1}{4}m\omega^2 A^2 \left[1 + \cos(2\omega t + 2\varphi)\right].$
 
• Tần số góc biến thiên của năng lượng là $\omega' = 2\omega$.
 
• Chu kì biến thiên tuần hoàn của năng lượng:
 $T' = \dfrac{2\pi}{\omega'} = \dfrac{2\pi}{2\omega} = \dfrac{T}{2}.$
}
\end{ex}

% ===== Câu 13 =====
% ID: L11C1B5-04-DS
% Mức: TH
\begin{ex}
% ID: L11C1B5-04-DS
% Mức: TH
Động năng và thế năng của một vật dao động điều hòa có chu kì dao động của li độ là $T$, tần số dao động là $f$. Khảo sát tính chất tuần hoàn theo thời gian của các loại năng lượng này:
\choiceTF
{Động năng và thế năng biến thiên tuần hoàn theo thời gian với chu kì bằng chu kì dao động $T$ của li độ.}
{\True Động năng và thế năng của vật biến thiên tuần hoàn với tần số biến thiên gấp hai lần tần số dao động $f$ của li độ.}
{Tần số góc biến thiên của động năng và thế năng bằng đúng tần số góc $\omega$ của li độ vật.}
{\True Khoảng thời gian ngắn nhất giữa hai lần liên tiếp động năng bằng thế năng của hệ dao động là $T/4$.}
\loigiai{
\begin{itemchoice}
\item (Sai) Động năng và thế năng biến thiên tuần hoàn theo thời gian với chu kì bằng chu kì dao động $T$ của li độ. (Vì): Mệnh đề này sai vì động năng và thế năng biến thiên với chu kì bằng một nửa chu kì dao động của li độ: $T' = T/2$
\item (Đúng) Động năng và thế năng của vật biến thiên tuần hoàn với tần số biến thiên gấp hai lần tần số dao động $f$ của li độ. (Vì): Mệnh đề này đúng vì tần số biến thiên của năng lượng là $f' = 2f$, tức là gấp đôi tần số dao động của li độ vật
\item (Sai) Tần số góc biến thiên của động năng và thế năng bằng đúng tần số góc $\omega$ của li độ vật. (Vì): Mệnh đề này sai vì tần số góc biến thiên của động năng và thế năng bằng hai lần tần số góc của li độ: $\omega' = 2\omega$
\item (Đúng) Khoảng thời gian ngắn nhất giữa hai lần liên tiếp động năng bằng thế năng của hệ dao động là $T/4$.
\end{itemchoice}
}
\end{ex}

% ===== Câu 14 =====
% ID: L11C1B5-04-TL
% Mức: TH
\dangbt{Chưa có dạng}
\begin{ex}
% ID: L11C1B5-04-TL
% Mức: TH
Một vật có khối lượng $m = 0$, $4 \mathrm{kg}$, dao động điều hoà với chu kì $T = 0$,2 $\pi(s),$ biên độ bằng $10 \mathrm{cm}$. Tính cơ năng của dao động. $\nguon{ly11 kntt.pdf}$
\choiceTF

\loigiai{
Cơ năng của dao động: $W = \dfrac{m\omega^{2}A^{2}}{2} = \dfrac{m}{2}\left( \dfrac{2\pi}{T} \right)^{2}A^{2} = \dfrac{0,4 \cdot \left( \dfrac{2\pi}{0,2\pi} \right)^{2} \cdot (0,1)^{2}}{2} = 0,2\text{J}.$
}
\end{ex}

% ===== Câu 15 =====
% ID: L11C1B5-04-TLN
% Mức: TH
\dangbt{Nhận diện phát biểu về năng lượng dao động}
\begin{ex}
% ID: L11C1B5-04-TLN
% Mức: TH
Con lắc lò xo có độ cứng $50\,\text{N}$ /m và biên độ $4\,\text{cm}$. Cơ năng theo đơn vị $J$ bằng bao nhiêu?
\shortans{0,04}
\end{ex}

% ===== Câu 16 =====
% ID: L11C1B5-04-TN
% Mức: NB
\dangbt{Chu kì biến thiên của động năng, thế năng}
\begin{ex}
% ID: L11C1B5-04-TN
% Mức: NB
Một chất điểm dao động điều hòa với tần số $f = 5\ \mathrm{Hz}$. Thế năng và động năng của chất điểm biến thiên tuần hoàn theo thời gian với tần số bằng:
\choice
{$5\ \mathrm{Hz}$}
{\True $10\ \mathrm{Hz}$}
{$2,5\ \mathrm{Hz}$}
{$20\ \mathrm{Hz}$}
\loigiai{
Tần số biến thiên của động năng và thế năng gấp hai lần tần số dao động của li độ:
  \begin{aligned}
 f' &= 2f 
&= 2 $\cdot5 = 10\\mathrm{Hz}$.
 \end{aligned}
}
\end{ex}

% ===== Câu 17 =====
% ID: L11C1B5-05-DS
% Mức: TH
\begin{ex}
% ID: L11C1B5-05-DS
% Mức: TH
Động năng và thế năng của một vật dao động điều hòa có chu kì dao động của li độ là $T$, tần số dao động là $f$. Khảo sát tính chất tuần hoàn theo thời gian của các loại năng lượng này:
\choiceTF
{Động năng và thế năng biến thiên tuần hoàn theo thời gian với chu kì bằng chu kì dao động $T$ của li độ.}
{\True Động năng và thế năng của vật biến thiên tuần hoàn với tần số biến thiên gấp hai lần tần số dao động $f$ của li độ.}
{Tần số góc biến thiên của động năng và thế năng bằng đúng tần số góc $\omega$ của li độ vật.}
{\True Khoảng thời gian ngắn nhất giữa hai lần liên tiếp động năng bằng thế năng của hệ dao động là $T/4$.}
\loigiai{
\begin{itemchoice}
\item (Sai) Động năng và thế năng biến thiên tuần hoàn theo thời gian với chu kì bằng chu kì dao động $T$ của li độ. (Vì): Mệnh đề này sai vì động năng và thế năng biến thiên với chu kì bằng một nửa chu kì dao động của li độ: $T' = T/2$
\item (Đúng) Động năng và thế năng của vật biến thiên tuần hoàn với tần số biến thiên gấp hai lần tần số dao động $f$ của li độ. (Vì): Mệnh đề này đúng vì tần số biến thiên của năng lượng là $f' = 2f$, tức là gấp đôi tần số dao động của li độ vật
\item (Sai) Tần số góc biến thiên của động năng và thế năng bằng đúng tần số góc $\omega$ của li độ vật. (Vì): Mệnh đề này sai vì tần số góc biến thiên của động năng và thế năng bằng hai lần tần số góc của li độ: $\omega' = 2\omega$
\item (Đúng) Khoảng thời gian ngắn nhất giữa hai lần liên tiếp động năng bằng thế năng của hệ dao động là $T/4$.
\end{itemchoice}
}
\end{ex}

% ===== Câu 18 =====
% ID: L11C1B5-05-TL
% Mức: TH
\dangbt{Chưa có dạng}
\begin{ex}
% ID: L11C1B5-05-TL
% Mức: TH
Một chất điểm có khối lượng $100\,\mathrm{g}$ dao động điều hoà trên quỹ đạo là đoạn thẳng MN (dài hơn $8 \mathrm{cm}$). Tại điểm $P$ cách $M$ một khoảng $4 \mathrm{cm}$ và tại điểm $Q$ cách $N$ một khoảng $2 \mathrm{cm}$, chất điểm có động năng tương ứng là 32$\cdot$ 10^{-3} $J$ và 18$\cdot$ 10^{-3} J. Tính tốc độ trung binh khi vật đi từ $P$ đến Q. $\nguon{ly11kntt.pdf}$
\choiceTF

\loigiai{
Tốc độ tại P: $v_{P} = \sqrt{\dfrac{2W_{dP}}{m}} = 80\text{cm/s}$; tại Q: $v_{Q} = \sqrt{\dfrac{2W_{dQ}}{m}} = 60\text{cm/s}$.

Do v_(P) > v_(Q) nên li độ |x_(P)|<|x_(Q)|: [Một chất điểm có khối lượng $100\,\text{g}$ dao động điều hoà]

[Một chất điểm có khối lượng $100\,\text{g}$ dao động điều hoà]

Giải hệ ta được: $A$ = $10\,\mathrm{cm}$ và ω = $10\,\mathrm{rad/s}$.

[Một chất điểm có khối lượng $100\,\text{g}$ dao động điều hoà]

Quãng đường PQ = OP + OQ = (A-4) + (A-2) = $14\,\mathrm{cm}$.

Thời gian vật đi từ $P$ đến $Q$ là $\Delta$ t với: $\text{\Delta }t = \dfrac{\text{\Delta }\varphi}{\omega}$.

$\text{\Delta }\varphi = \pi - \left( {\varphi_{1} + \varphi_{2}} \right) = \dfrac{\pi}{2}$, với $\left. \cos\varphi_{1} = \dfrac{\text{OP}}{\text{A}};\cos\varphi_{2} = \dfrac{\text{OQ}}{\text{A}}\Rightarrow\text{\Delta t} = \dfrac{\text{T}}{4} = \dfrac{\pi}{20} \right.\Rightarrow$  Tốc độ trung bình khi vật đi từ $P$ đến Q: $\overline{v} = \dfrac{PQ}{\text{\Delta }t} = \dfrac{14}{\dfrac{\pi}{20}} \approx 89\text{cm/s}$.
}
\end{ex}

% ===== Câu 19 =====
% ID: L11C1B5-05-TLN
% Mức: VD
\dangbt{Nhận diện phát biểu về năng lượng dao động}
\begin{ex}
% ID: L11C1B5-05-TLN
% Mức: VD
Con lắc lò xo có độ cứng $80\,\text{N}$ /m và biên độ $6\,\text{cm}$. Cơ năng theo đơn vị $J$ bằng bao nhiêu?
\shortans{0,144}
\end{ex}

% ===== Câu 20 =====
% ID: L11C1B5-05-TN
% Mức: NB
\dangbt{Chu kì biến thiên của động năng, thế năng}
\begin{ex}
% ID: L11C1B5-05-TN
% Mức: NB
Một chất điểm dao động điều hòa với tần số $f = 5\ \mathrm{Hz}$. Thế năng và động năng của chất điểm biến thiên tuần hoàn theo thời gian với tần số bằng:
\choice
{$5\ \mathrm{Hz}$}
{\True $10\ \mathrm{Hz}$}
{$2,5\ \mathrm{Hz}$}
{$20\ \mathrm{Hz}$}
\loigiai{
Tần số biến thiên của động năng và thế năng gấp hai lần tần số dao động của li độ:
  \begin{aligned}
 f' &= 2f 
&= 2 $\cdot5 = 10\\mathrm{Hz}$.
 \end{aligned}
}
\end{ex}

% ===== Câu 21 =====
% ID: L11C1B5-06-DS
% Mức: TH
\begin{ex}
% ID: L11C1B5-06-DS
% Mức: TH
Động năng và thế năng của một vật dao động điều hòa có chu kì dao động của li độ là $T$, tần số dao động là $f$. Khảo sát tính chất tuần hoàn theo thời gian của các loại năng lượng này:
\choiceTF
{Động năng và thế năng biến thiên tuần hoàn theo thời gian với chu kì bằng chu kì dao động $T$ của li độ.}
{\True Động năng và thế năng của vật biến thiên tuần hoàn với tần số biến thiên gấp hai lần tần số dao động $f$ của li độ.}
{Tần số góc biến thiên của động năng và thế năng bằng đúng tần số góc $\omega$ của li độ vật.}
{\True Khoảng thời gian ngắn nhất giữa hai lần liên tiếp động năng bằng thế năng của hệ dao động là $T/4$.}
\loigiai{
\begin{itemchoice}
\item (Sai) Động năng và thế năng biến thiên tuần hoàn theo thời gian với chu kì bằng chu kì dao động $T$ của li độ. (Vì): Mệnh đề này sai vì động năng và thế năng biến thiên với chu kì bằng một nửa chu kì dao động của li độ: $T' = T/2$
\item (Đúng) Động năng và thế năng của vật biến thiên tuần hoàn với tần số biến thiên gấp hai lần tần số dao động $f$ của li độ. (Vì): Mệnh đề này đúng vì tần số biến thiên của năng lượng là $f' = 2f$, tức là gấp đôi tần số dao động của li độ vật
\item (Sai) Tần số góc biến thiên của động năng và thế năng bằng đúng tần số góc $\omega$ của li độ vật. (Vì): Mệnh đề này sai vì tần số góc biến thiên của động năng và thế năng bằng hai lần tần số góc của li độ: $\omega' = 2\omega$
\item (Đúng) Khoảng thời gian ngắn nhất giữa hai lần liên tiếp động năng bằng thế năng của hệ dao động là $T/4$.
\end{itemchoice}
}
\end{ex}

% ===== Câu 22 =====
% ID: L11C1B5-06-TL
% Mức: TH
\dangbt{Chưa có dạng}
\begin{ex}
% ID: L11C1B5-06-TL
% Mức: TH
Một con lắc lò xo treo thẳng đứng vào điểm $I$ cố định, quả cầu có khối lượng $100\,\mathrm{g}$. Con lắc dao động điều hoà theo phương trình $x = 4 cos_10\sqrt{5t}(\mathrm{cm})$ với t tính theo giây. Lấy $g = 10 m/s_2$. Tính độ lớn lực đàn hồi lớn nhất và nhỏ nhất do lò xo tác dụng lên điểm I. $\nguon{ly11 kntt.pdf}$
\choiceTF

\loigiai{
Độ dãn của lò xo khi vật ở vị trí cân bằng: $\text{\Delta }l_{0} = \dfrac{\text{g}}{\omega^{2}} = \dfrac{10}{500} = 2\text{cm}$.

Biên độ dao động $A$ = $4\,\mathrm{cm}$.

Do A> $\Delta$ l₀ nên F_(min) = 0 (lúc lò xo không biến dạng).

Độ cứng của lò xo: k = $\dfrac{mg}{\text{\Delta }l_{0}} = \dfrac{0,1 \cdot 10}{0,02}$ = $50\,\mathrm{N}$ /m.

Lực đàn hồi cực đại F_(max) = k $(\Delta l₀+A)$ = 50.0,06 = $3\,\mathrm{N}$.
}
\end{ex}

% ===== Câu 23 =====
% ID: L11C1B5-06-TLN
% Mức: VD
\dangbt{Nhận diện phát biểu về năng lượng dao động}
\begin{ex}
% ID: L11C1B5-06-TLN
% Mức: VD
Con lắc lò xo có độ cứng $100\,\text{N}$ /m và biên độ $8\,\text{cm}$. Cơ năng theo đơn vị $J$ bằng bao nhiêu?
\shortans{0,32}
\end{ex}

% ===== Câu 24 =====
% ID: L11C1B5-06-TN
% Mức: NB
\dangbt{Chu kì biến thiên của động năng, thế năng}
\begin{ex}
% ID: L11C1B5-06-TN
% Mức: NB
Một chất điểm dao động điều hòa với tần số $f = 5\ \mathrm{Hz}$. Thế năng và động năng của chất điểm biến thiên tuần hoàn theo thời gian với tần số bằng:
\choice
{$5\ \mathrm{Hz}$}
{\True $10\ \mathrm{Hz}$}
{$2,5\ \mathrm{Hz}$}
{$20\ \mathrm{Hz}$}
\loigiai{
Tần số biến thiên của động năng và thế năng gấp hai lần tần số dao động của li độ:
  \begin{aligned}
 f' &= 2f 
&= 2 $\cdot5 = 10\\mathrm{Hz}$.
 \end{aligned}
}
\end{ex}

% ===== Câu 25 =====
% ID: L11C1B5-07-DS
% Mức: TH
\begin{ex}
% ID: L11C1B5-07-DS
% Mức: TH
Một con lắc lò xo treo thẳng đứng dao động điều hòa quanh vị trí cân bằng. Trong quá trình dao động, ta xét sự biến đổi năng lượng của hệ:
\choiceTF
{\True Thế năng của hệ dao động này bao gồm thế năng đàn hồi của lò xo và thế năng trọng trường của quả cầu nặng.}
{\True Trong suốt quá trình dao động điều hòa, tổng động năng và thế năng của hệ luôn được bảo toàn nếu bỏ qua mọi lực cản.}
{Động năng của quả cầu nặng đạt giá trị cực đại khi vật đi qua vị trí lò xo không biến dạng.}
{\True Khi quả cầu đi thẳng từ vị trí thấp nhất lên vị trí cao nhất của quỹ đạo, thế năng trọng trường của quả cầu tăng liên tục.}
\loigiai{
\begin{itemchoice}
\item (Đúng) Thế năng của hệ dao động này bao gồm thế năng đàn hồi của lò xo và thế năng trọng trường của quả cầu nặng. (Vì): Mệnh đề này đúng vì con lắc lò xo thẳng đứng chịu tác dụng của cả lực đàn hồi và trọng lực nên thế năng của hệ là tổng thế năng đàn hồi và thế năng trọng trường
\item (Đúng) Trong suốt quá trình dao động điều hòa, tổng động năng và thế năng của hệ luôn được bảo toàn nếu bỏ qua mọi lực cản. (Vì): Mệnh đề này đúng vì khi bỏ qua các lực hao tán như ma sát, cơ năng của hệ luôn là một đại lượng bảo toàn
\item (Sai) Động năng của quả cầu nặng đạt giá trị cực đại khi vật đi qua vị trí lò xo không biến dạng. (Vì): Mệnh đề này sai vì động năng đạt cực đại khi tốc độ đạt cực đại, tức là tại vị trí cân bằng, không phải tại vị trí lò xo không biến dạng (trừ khi tại vị trí cân bằng lò xo không biến dạng, điều này không xảy ra với con lắc lò xo treo thẳng đứng
\item (Đúng) Khi quả cầu đi thẳng từ vị trí thấp nhất lên vị trí cao nhất của quỹ đạo, thế năng trọng trường của quả cầu tăng liên tục. (Vì): Mệnh đề này đúng vì khi đi từ vị trí thấp nhất lên cao nhất, độ cao tăng liên tục nên thế năng trọng trường $W_{\text{tr}} = mgh$ tăng liên tục
\end{itemchoice}
}
\end{ex}

% ===== Câu 26 =====
% ID: L11C1B5-07-TL
% Mức: TH
\dangbt{Chưa có dạng}
\begin{ex}
% ID: L11C1B5-07-TL
% Mức: TH
Một con lắc lò xo treo thẳng đứng. Biết rằng trong quá trình dao động, tỉ số giữa độ lớn lực 7 đàn hồi lớn nhất và nhỏ nhất là 3, biên độ dao động là $10 \mathrm{cm}$. Lấy $g = 10 m/s_2$. Tính tần số dao động của vật. $\nguon{ly11 kntt.pdf}$
\choiceTF

\loigiai{
$\left. \dfrac{F_{\text{max}}}{F_{\text{min}}} = \dfrac{k\left( {\text{\Delta }l_{0} + A} \right)}{k\left( {\text{\Delta }l_{0} - A} \right)} = \dfrac{7}{3}\Rightarrow 3\left( {\text{\Delta }l_{0} + A} \right) = 7\left( {\text{\Delta }l_{0} - A} \right) \right.\Rightarrow\Deltal₀= 2,5\,\mathrm{A}=25\,\mathrm{cm}=0,25\,\mathrm{m}$.

Với $\Delta$ l₀ là độ dãn của lò xo tại vị trí cân bằng.

$\left. \omega = \sqrt{\dfrac{g}{\text{\Delta }l_{0}}} = \sqrt{\dfrac{10}{0,25}} = 2\pi\left( {\text{rad}/\text{s}} \right)\Rightarrow f = \dfrac{\omega}{2\pi} = 1\text{Hz} \right.$.
}
\end{ex}

% ===== Câu 27 =====
% ID: L11C1B5-07-TLN
% Mức: VD
\dangbt{Năng lượng của con lắc đơn dao động điều hòa}
\begin{ex}
% ID: L11C1B5-07-TLN
% Mức: VD
Một con lắc đơn có vật nặng $0,2\,\text{kg}$, chiều dài $1\,\text{m}$, dao động điều hòa với biên độ góc 0,1 rad. Lấy g = $10\,\text{m}$ /s². Cơ năng của con lắc theo đơn vị $J$ bằng bao nhiêu?
\shortans{0,01}
\loigiai{
Với góc nhỏ, $W$ = mglα₀²/2 = 0,2·10·1·0,1²/2 = $0,01\,\text{J}$.
}
\end{ex}

% ===== Câu 28 =====
% ID: L11C1B5-07-TN
% Mức: TH
\dangbt{Chu kì biến thiên của động năng, thế năng}
\begin{ex}
% ID: L11C1B5-07-TN
% Mức: TH
Khoảng thời gian ngắn nhất giữa hai lần liên tiếp động năng của vật dao động điều hòa bằng thế năng của hệ là:
\choice
{$\dfrac{T}{2}$}
{\True $\dfrac{T}{4}$}
{$\dfrac{T}{8}$}
{$T$}
\loigiai{
• Động năng bằng thế năng tại các vị trí có li độ $x = \pm \dfrac{A}{\sqrt{2}}$.
 
• Biểu diễn trên đường tròn lượng giác, các vị trí này cách đều nhau một góc $90^\circ\left(\dfrac{\pi}{2}\ \mathrm{rad}\right)$.
 
• Khoảng thời gian ngắn nhất ứng với góc quét $90^\circ$, tức bằng một phần tư chu kì:
 $\Delta t = \dfrac{T}{4}.$
}
\end{ex}

% ===== Câu 29 =====
% ID: L11C1B5-08-DS
% Mức: TH
\begin{ex}
% ID: L11C1B5-08-DS
% Mức: TH
Một con lắc lò xo treo thẳng đứng dao động điều hòa quanh vị trí cân bằng. Trong quá trình dao động, ta xét sự biến đổi năng lượng của hệ:
\choiceTF
{\True Thế năng của hệ dao động này bao gồm thế năng đàn hồi của lò xo và thế năng trọng trường của quả cầu nặng.}
{\True Trong suốt quá trình dao động điều hòa, tổng động năng và thế năng của hệ luôn được bảo toàn nếu bỏ qua mọi lực cản.}
{Động năng của quả cầu nặng đạt giá trị cực đại khi vật đi qua vị trí lò xo không biến dạng.}
{\True Khi quả cầu đi thẳng từ vị trí thấp nhất lên vị trí cao nhất của quỹ đạo, thế năng trọng trường của quả cầu tăng liên tục.}
\loigiai{
\begin{itemchoice}
\item (Đúng) Thế năng của hệ dao động này bao gồm thế năng đàn hồi của lò xo và thế năng trọng trường của quả cầu nặng. (Vì): Mệnh đề này đúng vì con lắc lò xo thẳng đứng chịu tác dụng của cả lực đàn hồi và trọng lực nên thế năng của hệ là tổng thế năng đàn hồi và thế năng trọng trường
\item (Đúng) Trong suốt quá trình dao động điều hòa, tổng động năng và thế năng của hệ luôn được bảo toàn nếu bỏ qua mọi lực cản. (Vì): Mệnh đề này đúng vì khi bỏ qua các lực hao tán như ma sát, cơ năng của hệ luôn là một đại lượng bảo toàn
\item (Sai) Động năng của quả cầu nặng đạt giá trị cực đại khi vật đi qua vị trí lò xo không biến dạng. (Vì): Mệnh đề này sai vì động năng đạt cực đại khi tốc độ đạt cực đại, tức là tại vị trí cân bằng, không phải tại vị trí lò xo không biến dạng (trừ khi tại vị trí cân bằng lò xo không biến dạng, điều này không xảy ra với con lắc lò xo treo thẳng đứng
\item (Đúng) Khi quả cầu đi thẳng từ vị trí thấp nhất lên vị trí cao nhất của quỹ đạo, thế năng trọng trường của quả cầu tăng liên tục. (Vì): Mệnh đề này đúng vì khi đi từ vị trí thấp nhất lên cao nhất, độ cao tăng liên tục nên thế năng trọng trường $W_{\text{tr}} = mgh$ tăng liên tục
\end{itemchoice}
}
\end{ex}

% ===== Câu 30 =====
% ID: L11C1B5-08-TL
% Mức: TH
\dangbt{Chưa có dạng}
\begin{ex}
% ID: L11C1B5-08-TL
% Mức: TH
Một con lắc đơn dao động điều hoà với biên độ góc $\alphamax$. Lấy mốc cơ năng tại vị trí cân bằng. Tính li độ góc của con lắc khi nó ở vị trí có động năng bằng thế năng. $\nguon{ly11 kntt.pdf}$
\choiceTF

\loigiai{
Khi động năng bằng thế năng: W_(t) = W_(d), ta có:

$\left. 2\text{W}_{\text{t}} = \text{W}\Leftrightarrow 2\dfrac{\text{mgl}\alpha^{2}}{2} = \dfrac{\text{mgl}\alpha_{\text{max}}^{2}}{2}\Rightarrow\alpha = \pm \dfrac{\alpha_{\text{max}}\sqrt{2}}{2} \right.$
}
\end{ex}

% ===== Câu 31 =====
% ID: L11C1B5-08-TLN
% Mức: VD
\dangbt{Năng lượng của con lắc đơn dao động điều hòa}
\begin{ex}
% ID: L11C1B5-08-TLN
% Mức: VD
Một con lắc đơn có vật nặng $0,2\,\text{kg}$, chiều dài $1\,\text{m}$, dao động điều hòa với biên độ góc 0,1 rad. Lấy g = $10\,\text{m}$ /s². Cơ năng của con lắc theo đơn vị $J$ bằng bao nhiêu?
\shortans{30/12/1899}
\loigiai{
Với góc nhỏ, $W$ = mglα₀²/2 = 0,2·10·1·0,1²/2 = $0,01\,\text{J}$.
}
\end{ex}

% ===== Câu 32 =====
% ID: L11C1B5-08-TN
% Mức: TH
\dangbt{Chu kì biến thiên của động năng, thế năng}
\begin{ex}
% ID: L11C1B5-08-TN
% Mức: TH
Khoảng thời gian ngắn nhất giữa hai lần liên tiếp động năng của vật dao động điều hòa bằng thế năng của hệ là:
\choice
{$\dfrac{T}{2}$}
{\True $\dfrac{T}{4}$}
{$\dfrac{T}{8}$}
{$T$}
\loigiai{
• Động năng bằng thế năng tại các vị trí có li độ $x = \pm \dfrac{A}{\sqrt{2}}$.
 
• Biểu diễn trên đường tròn lượng giác, các vị trí này cách đều nhau một góc $90^\circ\left(\dfrac{\pi}{2}\ \mathrm{rad}\right)$.
 
• Khoảng thời gian ngắn nhất ứng với góc quét $90^\circ$, tức bằng một phần tư chu kì:
 $\Delta t = \dfrac{T}{4}.$
}
\end{ex}

% ===== Câu 33 =====
% ID: L11C1B5-09-DS
% Mức: TH
\begin{ex}
% ID: L11C1B5-09-DS
% Mức: TH
Một con lắc lò xo treo thẳng đứng dao động điều hòa quanh vị trí cân bằng. Trong quá trình dao động, ta xét sự biến đổi năng lượng của hệ:
\choiceTF
{\True Thế năng của hệ dao động này bao gồm thế năng đàn hồi của lò xo và thế năng trọng trường của quả cầu nặng.}
{\True Trong suốt quá trình dao động điều hòa, tổng động năng và thế năng của hệ luôn được bảo toàn nếu bỏ qua mọi lực cản.}
{Động năng của quả cầu nặng đạt giá trị cực đại khi vật đi qua vị trí lò xo không biến dạng.}
{\True Khi quả cầu đi thẳng từ vị trí thấp nhất lên vị trí cao nhất của quỹ đạo, thế năng trọng trường của quả cầu tăng liên tục.}
\loigiai{
\begin{itemchoice}
\item (Đúng) Thế năng của hệ dao động này bao gồm thế năng đàn hồi của lò xo và thế năng trọng trường của quả cầu nặng. (Vì): Mệnh đề này đúng vì con lắc lò xo thẳng đứng chịu tác dụng của cả lực đàn hồi và trọng lực nên thế năng của hệ là tổng thế năng đàn hồi và thế năng trọng trường
\item (Đúng) Trong suốt quá trình dao động điều hòa, tổng động năng và thế năng của hệ luôn được bảo toàn nếu bỏ qua mọi lực cản. (Vì): Mệnh đề này đúng vì khi bỏ qua các lực hao tán như ma sát, cơ năng của hệ luôn là một đại lượng bảo toàn
\item (Sai) Động năng của quả cầu nặng đạt giá trị cực đại khi vật đi qua vị trí lò xo không biến dạng. (Vì): Mệnh đề này sai vì động năng đạt cực đại khi tốc độ đạt cực đại, tức là tại vị trí cân bằng, không phải tại vị trí lò xo không biến dạng (trừ khi tại vị trí cân bằng lò xo không biến dạng, điều này không xảy ra với con lắc lò xo treo thẳng đứng
\item (Đúng) Khi quả cầu đi thẳng từ vị trí thấp nhất lên vị trí cao nhất của quỹ đạo, thế năng trọng trường của quả cầu tăng liên tục. (Vì): Mệnh đề này đúng vì khi đi từ vị trí thấp nhất lên cao nhất, độ cao tăng liên tục nên thế năng trọng trường $W_{\text{tr}} = mgh$ tăng liên tục
\end{itemchoice}
}
\end{ex}

% ===== Câu 34 =====
% ID: L11C1B5-09-TL
% Mức: TH
\dangbt{Chưa có dạng}
\begin{ex}
% ID: L11C1B5-09-TL
% Mức: TH
Một con lắc lò xo gồm một lò xo nhẹ có độ cứng k, được treo thẳng đứng vào một giá cố định và một vật có khối lượng $m = 100\$, $\mathrm{g}$. Khi vật ở vị trí cân bằng $O$, lò xo dãn $2,5 \mathrm{cm}$. Kéo vật dọc theo trục của lò xo xuống dưới cách vị trí cân bằng $O$ một đoạn $2 \mathrm{cm}$ rồi truyền cho nó vận tốc có độ lớn $40\sqrt{3} \mathrm{cm}$ /s theo phương thẳng đứng, hướng xuống dưới. Chọn trục toạ độ Ox theo phương thẳng đứng, gốc tại $O$, chiều dương hướng lên trên, gốc thời gian là lúc vật bắt đầu dao động. Lấy $g =$ 10 $\mathrm{m/s}$ 2 $. Biết chiều dài tự nhiên của của lò xo là$ 50 \mathrm{cm} $. a) Tính độ cứng của lò xo, viết phương trình dao động và tính cơ năng dao động của vật. b) Xác định li độ và vận tốc của vật khi thế năng dao động bằng 1/3 động năng. c) Tính thế năng dao động, động năng và vận tốc của vật tại vị trí có li độ$ x = 2\sqrt{2cm} $. d) Tính chiều dài, lực đàn hồi cực đại, cực tiểu của lò xo trong quá trình dao động.$ \nguon{ly11kntt.pdf}
\choiceTF

\loigiai{
a) Gọi $\Delta$ l₀ là độ dãn của lò xo tại vị trí cân bằng, ta có: $\Delta$ l₀= $2,5\,\mathrm{cm}$ = $0,025\,\mathrm{m}$.

Tại vị trí cân bằng: $\left. \text{k} \cdot \text{\Delta }l_{0} = \text{mg}\Rightarrow\text{k} = \dfrac{\text{mg}}{\text{\Delta }l_{0}} = \dfrac{0,1 \cdot 10}{0,025} = 40\text{N/m} \right.$.

$\omega = \sqrt{\dfrac{k}{m}} = \sqrt{\dfrac{40}{0,1}} = 20\text{rad/s}$.

Theo đề bài, khi t = 0 thì x = -2 và v = -40 $\sqrt{3}$ cm/s

$\left. \Rightarrow A = \sqrt{x^{2} + \dfrac{v^{2}}{\omega^{2}}} = \sqrt{( - 2)^{2} + \dfrac{\left( 40\sqrt{3} \right)^{2}}{(20)^{2}}} = 4\text{cm} \right.$.

Vậy tại thời điểm t = 0 thì x = - $2\,\mathrm{cm}$ = - $\dfrac{A}{2}$ và v<0, nên $\varphi = \dfrac{2\pi}{3}$, phương trình dao động là: x = 40cos $\left( {20t + \dfrac{2\pi}{3}} \right)$ (cm)

Cơ năng của dao động: $W$ = $\dfrac{1}{2}m\omega^{2}A^{2} = \dfrac{1}{2} \cdot 0,1(20)^{2}(0,04)^{2} = 0,032\text{J}$

b) Khi thế năng bằng $\dfrac{1}{3}$ động năng;

$. W_{t} = \dfrac{1}{3}W_{d}\Leftrightarrow W = 4W_{t}\Leftrightarrow\dfrac{1}{2}m\omega^{2}A^{2} = 4\dfrac{1}{2}m\omega^{2}x^{2} .. \Rightarrow x = \pm \dfrac{A}{2} = \pm 2\text{cm};W_{t} = \dfrac{W}{4}\Rightarrow W_{d} = \dfrac{$ 3 $\mathrm{W}$ }{4} = \dfrac{3.0,032}{4} = 0,024\text{J} \right.\left. \Rightarrow v = \pm \sqrt{\dfrac{2W_{d}}{m}} = \pm \sqrt{\dfrac{2 \cdot 0,024}{0,1}} = \pm \dfrac{2\sqrt{3}}{5}\text{m}/\text{s}. \right. $c) Khi x = 0,02$ \sqrt{2} $m, ta có thế năng:$ W_{t} = \dfrac{1}{2}m\omega^{2}x^{2} = \dfrac{1}{2} \cdot 0,1(20)^{2}\left( 0,02\sqrt{2} \right)^{2} = 0,016\text{J}. $Động năng W_(d)=$ W $- W_(t) = 0,032-0,016 = 0,$ 016\,\mathrm{J} $.

Vận tốc của vật khi đó:$v = \pm \sqrt{\dfrac{2W_{\text{đ}}}{m}} = \pm \sqrt{\dfrac{2 \cdot 0,016}{0,1}} = \pm \dfrac{2\sqrt{2}}{5}\text{m}/\text{s}$.

d) Khi lực đàn hồi của lò xo tác dụng lên giá treo đạt cực đại thì độ biến dạng của lò xo là:$\Delta$l_(max) =$\Delta l₀+A)\Rightarrow$Lực đàn hồi cực đại F_(max) = k$\Delta l₀+A)$= 40(2,5+4).10^{-2} = 2,$6\,\mathrm{N}$.

Lò xo có chiều dài cực đại: l_(max) =$(l₀+\Delta l₀+A)$= 50+2,5+4 =$56,5\,\mathrm{cm}$.

Vì A>$\Delta$l, nên trong quá trình dao động có lúc vật qua vị trí lò xo không biến dạng, khi đó:$\Delta$l_(min)=0\Rightarrow  Lực đàn hồi cực tiểu F_(min)= 0. Khi đó, lò xo có chiều dài tự nhiên: l₀ =$ 50\,\mathrm{cm}.
}
\end{ex}

% ===== Câu 35 =====
% ID: L11C1B5-09-TLN
% Mức: VD
\dangbt{Năng lượng của con lắc đơn dao động điều hòa}
\begin{ex}
% ID: L11C1B5-09-TLN
% Mức: VD
Một con lắc đơn có vật nặng $0,2\,\text{kg}$, chiều dài $1\,\text{m}$, dao động điều hòa với biên độ góc 0,1 rad. Lấy g = $10\,\text{m}$ /s². Cơ năng của con lắc theo đơn vị $J$ bằng bao nhiêu?
\shortans{30/12/1899}
\loigiai{
Với góc nhỏ, $W$ = mglα₀²/2 = 0,2·10·1·0,1²/2 = $0,01\,\text{J}$.
}
\end{ex}

% ===== Câu 36 =====
% ID: L11C1B5-09-TN
% Mức: TH
\dangbt{Chu kì biến thiên của động năng, thế năng}
\begin{ex}
% ID: L11C1B5-09-TN
% Mức: TH
Khoảng thời gian ngắn nhất giữa hai lần liên tiếp động năng của vật dao động điều hòa bằng thế năng của hệ là:
\choice
{$\dfrac{T}{2}$}
{\True $\dfrac{T}{4}$}
{$\dfrac{T}{8}$}
{$T$}
\loigiai{
• Động năng bằng thế năng tại các vị trí có li độ $x = \pm \dfrac{A}{\sqrt{2}}$.
 
• Biểu diễn trên đường tròn lượng giác, các vị trí này cách đều nhau một góc $90^\circ\left(\dfrac{\pi}{2}\ \mathrm{rad}\right)$.
 
• Khoảng thời gian ngắn nhất ứng với góc quét $90^\circ$, tức bằng một phần tư chu kì:
 $\Delta t = \dfrac{T}{4}.$
}
\end{ex}

% ===== Câu 37 =====
% ID: L11C1B5-10-DS
% Mức: VD
\begin{ex}
% ID: L11C1B5-10-DS
% Mức: VD
Đồ thị động năng theo thời gian $W_d(t)$ của một vật dao động điều hòa là một đường tuần hoàn có chu kì biến thiên của năng lượng đo được là $T' = 0,2$ s.
\choiceTF
{\True Chu kì dao động của li độ vật bằng $0,4$ s.}
{\True Tần số dao động của li độ vật bằng $2,5$ Hz.}
{\True Khoảng thời gian giữa hai lần liên tiếp động năng bằng thế năng của vật bằng $0,1$ s.}
{Trong một chu kì dao động toàn phần của li độ vật, thế năng của hệ đạt giá trị cực đại $4$ lần.}
\loigiai{
\begin{itemchoice}
\item (Đúng) Chu kì dao động của li độ vật bằng $0,4$ s. (Vì): Mệnh đề này đúng vì chu kì dao động $T$ của li độ gấp đôi chu kì biến thiên năng lượng $T'$: $T = 2T' = 2 \cdot 0,2 = 0,4$ s
\item (Đúng) Tần số dao động của li độ vật bằng $2,5$ Hz. (Vì): Mệnh đề này đúng vì tần số dao động li độ của vật là: $f = \dfrac{1}{T} = \dfrac{1}{0,4} = 2,5$ Hz
\item (Đúng) Khoảng thời gian giữa hai lần liên tiếp động năng bằng thế năng của vật bằng $0,1$ s. (Vì): Mệnh đề này đúng vì khoảng thời gian ngắn nhất giữa hai lần liên tiếp động năng bằng thế năng là $T/4 = 0,4/4 = 0,1$ s
\item (Sai) Trong một chu kì dao động toàn phần của li độ vật, thế năng của hệ đạt giá trị cực đại $4$ lần. (Vì): Mệnh đề này sai vì thế năng đạt giá trị cực đại khi vật ở vị trí biên, trong một chu kì vật đi qua biên hai lần (một lần biên dương, một lần biên âm) nên thế năng đạt cực đại 2 lần
\end{itemchoice}
}
\end{ex}

% ===== Câu 38 =====
% ID: L11C1B5-10-TL
% Mức: TH
\dangbt{Nhận diện phát biểu về năng lượng dao động}
\begin{bt}
% ID: L11C1B5-10-TL
% Mức: TH
Con lắc lò xo có $k=50$ N/m, biên độ $4\,\text{cm}$. Tính cơ năng; động năng và thế năng tại $x=2$ cm.
\end{bt}

% ===== Câu 39 =====
% ID: L11C1B5-10-TLN
% Mức: VD
\dangbt{Năng lượng của con lắc đơn dao động điều hòa}
\begin{ex}
% ID: L11C1B5-10-TLN
% Mức: VD
Một con lắc đơn có vật nặng $0,4\,\text{kg}$, chiều dài $1\,\text{m}$, dao động điều hòa với biên độ góc 0,1 rad. Lấy g = $10\,\text{m}$ /s². Cơ năng của con lắc theo đơn vị $J$ bằng bao nhiêu?
\shortans{30/12/1899}
\loigiai{
$W$ = mglα₀²/2 = 0,4·10·1·0,1²/2 = $0,02\,\text{J}$.
}
\end{ex}

% ===== Câu 40 =====
% ID: L11C1B5-10-TN
% Mức: TH
\dangbt{Chu kì biến thiên của động năng, thế năng}
\begin{ex}
% ID: L11C1B5-10-TN
% Mức: TH
Đồ thị động năng theo thời gian $W_{\text{đ}}(t)$ của một vật dao động điều hòa có chu kì $T$ là một đường tuần hoàn. Khoảng thời gian giữa hai lần liên tiếp động năng bằng $0$ (đáy của đồ thị động năng) bằng:
\choice
{$T$}
{\True $\dfrac{T}{2}$}
{$\dfrac{T}{4}$}
{$2T$}
\loigiai{
• Động năng bằng $0$ khi vận tốc của vật bằng $0$, tức là khi vật đang đi qua các vị trí biên.
 
• Trong một chu kì dao động đầy đủ, vật đi qua vị trí biên hai lần (biên âm và biên dương).
 
• Khoảng thời gian giữa hai lần liên tiếp vật đạt tới biên bằng nửa chu kì:
 $\Delta t = \dfrac{T}{2}.$
}
\end{ex}

% ===== Câu 41 =====
% ID: L11C1B5-100-TN
% Mức: VD
\dangbt{Xác định độ cứng lò xo từ năng lượng ban đầu}
\begin{ex}
% ID: L11C1B5-100-TN
% Mức: VD
Một con lắc lò xo dao động điều hòa. Đồ thị biểu diễn thế năng $W_t$ theo li độ $x$ được cho như hình vẽ. Từ đồ thị, hãy xác định độ cứng $k$ của lò xo và biên độ dao động $A$.
\choice
{\True $k = 100$ N/m; $A = 6$ cm}
{$k = 200$ N/m; $A = 4$ cm}
{$k = 150$ N/m; $A = 5$ cm}
{$k = 100$ N/m; $A = 4$ cm}
\loigiai{
Từ đồ thị, ta đọc được: thế năng cực đại $W_{t,max} = 0,18J$ đạt tại $x = \pm 6$ cm, và tại $x = 0$ thì $W_t = 0$. Công thức thế năng: $W_t = \frac{1}{2}kx^2$. Tại biên $x = A = 6\,\mathrm{cm}= 0,06$ m, $W_{t,max} = \frac{1}{2}kA^2$. Suy ra: $k = \frac{2W_{t,max}}{A^2} = \frac{2 \times 0,18}{(0,06)^2} = \frac{0,36}{0,0036} = 100$ N/m. Vậy $k = 100$ N/m và $A = 6$ cm. Đáp án đúng là A.
}
\end{ex}

% ===== Câu 42 =====
% ID: L11C1B5-101-TN
% Mức: VD
\begin{ex}
% ID: L11C1B5-101-TN
% Mức: VD
Một con lắc lò xo dao động điều hòa. Đồ thị biểu diễn thế năng $W_t$ theo thời gian $t$ được cho như hình vẽ. Biết khối lượng vật $m = 400$ g. Tính biên độ dao động $A$ và tần số góc $\omega$ của con lắc.
\choice
{$A = 5$ cm; $\omega = 10\pi$ rad/s}
{\True $A = 4$ cm; $\omega = 5\pi$ rad/s}
{$A = 5$ cm; $\omega = 5\pi$ rad/s}
{$A = 4$ cm; $\omega = 10\pi$ rad/s}
\loigiai{
Từ đồ thị, ta đọc được: thế năng dao động tuần hoàn với giá trị cực đại $W_{t,max} = 0,16J$ và chu kỳ của $W_t$ là $T_{Wt} = 0,2$ s. Vì thế năng dao động với chu kỳ bằng nửa chu kỳ dao động nên chu kỳ dao động thực là $T = 2T_{Wt} = 0,4$ s. Tần số góc: $\omega = \frac{2\pi}{T} = \frac{2\pi}{0,4} = 5\pi$ rad/s. Độ cứng lò xo: $k = m\omega^2 = 0,4 \times (5\pi)^2 = 0,4 \times 25\pi^2 \approx 98,7$ N/m. Biên độ từ $W_{t,max} = \frac{1}{2}kA^2$: $A = \sqrt{\frac{2W_{t,max}}{k}} = \sqrt{\frac{2 \times 0,16}{0,4 \times 25\pi^2}} = \sqrt{\frac{0,32}{10\pi^2}} \approx \sqrt{0,00324} \approx 0,04\,\mathrm{m}= 4$ cm. Vậy $A = 4$ cm và $\omega = 5\pi$ rad/s. Đáp án đúng là B.
}
\end{ex}

% ===== Câu 43 =====
% ID: L11C1B5-102-TN
% Mức: VD
\begin{ex}
% ID: L11C1B5-102-TN
% Mức: VD
Một con lắc lò xo dao động điều hòa có cơ năng 0,08 $J$ và biên độ $4\,\text{cm}$. Độ cứng của lò xo bằng bao nhiêu?
\choice
{$50\,\text{N}$ /m}
{$80\,\text{N}$ /m}
{\True $100\,\text{N}$ /m}
{$120\,\text{N}$ /m}
\loigiai{
$W$ = kA²/2 nên k = $2\,\text{W}$ /A² = 2·0,08/(0,04)² = $100\,\text{N}$ /m.
}
\end{ex}

% ===== Câu 44 =====
% ID: L11C1B5-103-TN
% Mức: VD
\begin{ex}
% ID: L11C1B5-103-TN
% Mức: VD
Một con lắc lò xo dao động điều hòa có cơ năng 0,08 $J$ và biên độ $4\,\text{cm}$. Độ cứng của lò xo bằng bao nhiêu?
\choice
{$50\,\text{N}$ /m}
{$80\,\text{N}$ /m}
{\True $100\,\text{N}$ /m}
{$120\,\text{N}$ /m}
\loigiai{
$W$ = kA²/2 nên k = $2\,\text{W}$ /A² = 2·0,08/(0,04)² = $100\,\text{N}$ /m.
}
\end{ex}

% ===== Câu 45 =====
% ID: L11C1B5-104-TN
% Mức: VD
\begin{ex}
% ID: L11C1B5-104-TN
% Mức: VD
Một con lắc lò xo dao động điều hòa có cơ năng 0,08 $J$ và biên độ $4\,\text{cm}$. Độ cứng của lò xo bằng bao nhiêu?
\choice
{$50\,\text{N}$ /m}
{$80\,\text{N}$ /m}
{\True $100\,\text{N}$ /m}
{$120\,\text{N}$ /m}
\loigiai{
$W$ = kA²/2 nên k = $2\,\text{W}$ /A² = 2·0,08/(0,04)² = $100\,\text{N}$ /m.
}
\end{ex}

% ===== Câu 46 =====
% ID: L11C1B5-11-DS
% Mức: VD
\dangbt{Chu kì biến thiên của động năng, thế năng}
\begin{ex}
% ID: L11C1B5-11-DS
% Mức: VD
Đồ thị động năng theo thời gian $W_d(t)$ của một vật dao động điều hòa là một đường tuần hoàn có chu kì biến thiên của năng lượng đo được là $T' = 0,2$ s.
\choiceTF
{\True Chu kì dao động của li độ vật bằng $0,4$ s.}
{\True Tần số dao động của li độ vật bằng $2,5$ Hz.}
{\True Khoảng thời gian giữa hai lần liên tiếp động năng bằng thế năng của vật bằng $0,1$ s.}
{Trong một chu kì dao động toàn phần của li độ vật, thế năng của hệ đạt giá trị cực đại $4$ lần.}
\loigiai{
\begin{itemchoice}
\item (Đúng) Chu kì dao động của li độ vật bằng $0,4$ s. (Vì): Mệnh đề này đúng vì chu kì dao động $T$ của li độ gấp đôi chu kì biến thiên năng lượng $T'$: $T = 2T' = 2 \cdot 0,2 = 0,4$ s
\item (Đúng) Tần số dao động của li độ vật bằng $2,5$ Hz. (Vì): Mệnh đề này đúng vì tần số dao động li độ của vật là: $f = \dfrac{1}{T} = \dfrac{1}{0,4} = 2,5$ Hz
\item (Đúng) Khoảng thời gian giữa hai lần liên tiếp động năng bằng thế năng của vật bằng $0,1$ s. (Vì): Mệnh đề này đúng vì khoảng thời gian ngắn nhất giữa hai lần liên tiếp động năng bằng thế năng là $T/4 = 0,4/4 = 0,1$ s
\item (Sai) Trong một chu kì dao động toàn phần của li độ vật, thế năng của hệ đạt giá trị cực đại $4$ lần. (Vì): Mệnh đề này sai vì thế năng đạt giá trị cực đại khi vật ở vị trí biên, trong một chu kì vật đi qua biên hai lần (một lần biên dương, một lần biên âm) nên thế năng đạt cực đại 2 lần
\end{itemchoice}
}
\end{ex}

% ===== Câu 47 =====
% ID: L11C1B5-11-TL
% Mức: VD
\dangbt{Nhận diện phát biểu về năng lượng dao động}
\begin{bt}
% ID: L11C1B5-11-TL
% Mức: VD
Con lắc lò xo có $k=80$ N/m, biên độ $6\,\text{cm}$. Tính cơ năng; động năng và thế năng tại $x=3$ cm.
\end{bt}

% ===== Câu 48 =====
% ID: L11C1B5-11-TLN
% Mức: VD
\dangbt{Năng lượng của con lắc đơn dao động điều hòa}
\begin{ex}
% ID: L11C1B5-11-TLN
% Mức: VD
Một con lắc đơn có vật nặng $0,4\,\text{kg}$, chiều dài $1\,\text{m}$, dao động điều hòa với biên độ góc 0,1 rad. Lấy g = $10\,\text{m}$ /s². Cơ năng của con lắc theo đơn vị $J$ bằng bao nhiêu?
\shortans{0,02}
\loigiai{
$W$ = mglα₀²/2 = 0,4·10·1·0,1²/2 = $0,02\,\text{J}$.
}
\end{ex}

% ===== Câu 49 =====
% ID: L11C1B5-11-TN
% Mức: TH
\dangbt{Chu kì biến thiên của động năng, thế năng}
\begin{ex}
% ID: L11C1B5-11-TN
% Mức: TH
Đồ thị động năng theo thời gian $W_{\text{đ}}(t)$ của một vật dao động điều hòa có chu kì $T$ là một đường tuần hoàn. Khoảng thời gian giữa hai lần liên tiếp động năng bằng $0$ (đáy của đồ thị động năng) bằng:
\choice
{$T$}
{\True $\dfrac{T}{2}$}
{$\dfrac{T}{4}$}
{$2T$}
\loigiai{
• Động năng bằng $0$ khi vận tốc của vật bằng $0$, tức là khi vật đang đi qua các vị trí biên.
 
• Trong một chu kì dao động đầy đủ, vật đi qua vị trí biên hai lần (biên âm và biên dương).
 
• Khoảng thời gian giữa hai lần liên tiếp vật đạt tới biên bằng nửa chu kì:
 $\Delta t = \dfrac{T}{2}.$
}
\end{ex}

% ===== Câu 50 =====
% ID: L11C1B5-12-DS
% Mức: VD
\begin{ex}
% ID: L11C1B5-12-DS
% Mức: VD
Đồ thị động năng theo thời gian $W_d(t)$ của một vật dao động điều hòa là một đường tuần hoàn có chu kì biến thiên của năng lượng đo được là $T' = 0,2$ s.
\choiceTF
{\True Chu kì dao động của li độ vật bằng $0,4$ s.}
{\True Tần số dao động của li độ vật bằng $2,5$ Hz.}
{\True Khoảng thời gian giữa hai lần liên tiếp động năng bằng thế năng của vật bằng $0,1$ s.}
{Trong một chu kì dao động toàn phần của li độ vật, thế năng của hệ đạt giá trị cực đại $4$ lần.}
\loigiai{
\begin{itemchoice}
\item (Đúng) Chu kì dao động của li độ vật bằng $0,4$ s. (Vì): Mệnh đề này đúng vì chu kì dao động $T$ của li độ gấp đôi chu kì biến thiên năng lượng $T'$: $T = 2T' = 2 \cdot 0,2 = 0,4$ s
\item (Đúng) Tần số dao động của li độ vật bằng $2,5$ Hz. (Vì): Mệnh đề này đúng vì tần số dao động li độ của vật là: $f = \dfrac{1}{T} = \dfrac{1}{0,4} = 2,5$ Hz
\item (Đúng) Khoảng thời gian giữa hai lần liên tiếp động năng bằng thế năng của vật bằng $0,1$ s. (Vì): Mệnh đề này đúng vì khoảng thời gian ngắn nhất giữa hai lần liên tiếp động năng bằng thế năng là $T/4 = 0,4/4 = 0,1$ s
\item (Sai) Trong một chu kì dao động toàn phần của li độ vật, thế năng của hệ đạt giá trị cực đại $4$ lần. (Vì): Mệnh đề này sai vì thế năng đạt giá trị cực đại khi vật ở vị trí biên, trong một chu kì vật đi qua biên hai lần (một lần biên dương, một lần biên âm) nên thế năng đạt cực đại 2 lần
\end{itemchoice}
}
\end{ex}

% ===== Câu 51 =====
% ID: L11C1B5-12-TL
% Mức: VD
\dangbt{Nhận diện phát biểu về năng lượng dao động}
\begin{bt}
% ID: L11C1B5-12-TL
% Mức: VD
Con lắc lò xo có $k=100$ N/m, biên độ $8\,\text{cm}$. Tính cơ năng; động năng và thế năng tại $x=4$ cm.
\end{bt}

% ===== Câu 52 =====
% ID: L11C1B5-12-TLN
% Mức: VD
\dangbt{Năng lượng của con lắc đơn dao động điều hòa}
\begin{ex}
% ID: L11C1B5-12-TLN
% Mức: VD
Một con lắc đơn có vật nặng $0,4\,\text{kg}$, chiều dài $1\,\text{m}$, dao động điều hòa với biên độ góc 0,1 rad. Lấy g = $10\,\text{m}$ /s². Cơ năng của con lắc theo đơn vị $J$ bằng bao nhiêu?
\shortans{0,02}
\loigiai{
$W$ = mglα₀²/2 = 0,4·10·1·0,1²/2 = $0,02\,\text{J}$.
}
\end{ex}

% ===== Câu 53 =====
% ID: L11C1B5-12-TN
% Mức: TH
\dangbt{Chu kì biến thiên của động năng, thế năng}
\begin{ex}
% ID: L11C1B5-12-TN
% Mức: TH
Đồ thị động năng theo thời gian $W_{\text{đ}}(t)$ của một vật dao động điều hòa có chu kì $T$ là một đường tuần hoàn. Khoảng thời gian giữa hai lần liên tiếp động năng bằng $0$ (đáy của đồ thị động năng) bằng:
\choice
{$T$}
{\True $\dfrac{T}{2}$}
{$\dfrac{T}{4}$}
{$2T$}
\loigiai{
• Động năng bằng $0$ khi vận tốc của vật bằng $0$, tức là khi vật đang đi qua các vị trí biên.
 
• Trong một chu kì dao động đầy đủ, vật đi qua vị trí biên hai lần (biên âm và biên dương).
 
• Khoảng thời gian giữa hai lần liên tiếp vật đạt tới biên bằng nửa chu kì:
 $\Delta t = \dfrac{T}{2}.$
}
\end{ex}

% ===== Câu 54 =====
% ID: L11C1B5-13-DS
% Mức: VD
\begin{ex}
% ID: L11C1B5-13-DS
% Mức: VD
Một vật có khối lượng $m = 400$ g thực hiện dao động điều hòa. Đồ thị thế năng theo thời gian $W_t(t)$ được biểu diễn bằng một đường sóng hình sin nâng trục có giá trị cực đại bằng $0,08$ J. Lấy $\pi^2 = 10$.
\choiceTF
{\True Cơ năng dao động của vật bằng $0,08$ J.}
{\True Khoảng thời gian từ lúc thế năng bằng 0 đến lúc thế năng cực đại lần đầu tiên tương ứng một phần tư chu kì dao động $T$ của vật.}
{\True Nếu chu kì biến thiên của thế năng là $0,1$ s thì chu kì dao động của vật là $0,2$ s.}
{Biên độ dao động của vật có giá trị bằng $10$ cm.}
\loigiai{
\begin{itemchoice}
\item (Đúng) Cơ năng dao động của vật bằng $0,08$ J. (Vì): Mệnh đề này đúng vì giá trị cực đại của thế năng chính bằng cơ năng của hệ dao động: $W = 0,08$ J
\item (Đúng) Khoảng thời gian từ lúc thế năng bằng 0 đến lúc thế năng cực đại lần đầu tiên tương ứng một phần tư chu kì dao động $T$ của vật. (Vì): Mệnh đề này đúng vì thế năng bằng 0 ở vị trí cân bằng và cực đại ở vị trí biên, thời gian ngắn nhất đi từ vị trí cân bằng ra biên bằng một phần tư chu kì dao động $T/4$
\item (Đúng) Nếu chu kì biến thiên của thế năng là $0,1$ s thì chu kì dao động của vật là $0,2$ s. (Vì): Mệnh đề này đúng vì chu kì dao động li độ gấp đôi chu kì thế năng: $T = 2T' = 2 \cdot 0,1 = 0,2$ s
\item (Sai) Biên độ dao động của vật có giá trị bằng $10$ cm. (Vì): Mệnh đề này sai vì từ chu kì $T = 0,2$ s ta có tần số góc $\omega = 10\pi$ rad/s. Cơ năng $W = \dfrac{1}{2}m\omega^2 A^2 \Rightarrow 0,08 = 0,5 \cdot 0,4 \cdot (10\pi)^2 \cdot A^2 \Rightarrow 0,08 = 200 A^2 \Rightarrow A^2 = 0,0004 \Rightarrow A = 0,02\,\mathrm{m}= 2$ cm
\end{itemchoice}
}
\end{ex}

% ===== Câu 55 =====
% ID: L11C1B5-13-TL
% Mức: TH
\dangbt{Năng lượng của con lắc đơn dao động điều hòa}
\begin{bt}
% ID: L11C1B5-13-TL
% Mức: TH
Con lắc đơn có $m=0.2$ kg, $l=1$ m, biên độ góc 0.1 rad, lấy $g=10$ m/s $^2$. Tính cơ năng trong gần đúng góc nhỏ.
\end{bt}

% ===== Câu 56 =====
% ID: L11C1B5-13-TLN
% Mức: TH
\dangbt{Phân tích đồ thị năng lượng}
\begin{ex}
% ID: L11C1B5-13-TLN
% Mức: TH
Con lắc lò xo có độ cứng $50\,\text{N}$ /m và biên độ $4\,\text{cm}$. Cơ năng theo đơn vị $J$ bằng bao nhiêu?
\shortans{0,04}
\end{ex}

% ===== Câu 57 =====
% ID: L11C1B5-13-TN
% Mức: VD
\dangbt{Chu kì biến thiên của động năng, thế năng}
\begin{ex}
% ID: L11C1B5-13-TN
% Mức: VD
Một chất điểm dao động điều hòa. Biết khoảng thời gian ngắn nhất giữa hai lần liên tiếp động năng bằng thế năng là $0,1\ \mathrm{s}$. Chu kì dao động của chất điểm là:
\choice
{$0,1\ \mathrm{s}$}
{$0,2\ \mathrm{s}$}
{\True $0,4\ \mathrm{s}$}
{$0,8\ \mathrm{s}$}
\loigiai{
Khoảng thời gian ngắn nhất giữa hai lần liên tiếp động năng bằng thế năng là $\Delta t = \dfrac{T}{4}$, do đó
  \begin{aligned}
 $\dfrac{T}{4}$ &= 0,1 
$\RightarrowT$ &= 4 \cdot 0,1 = 0,4 $\\mathrm{s}$.
 \end{aligned}
}
\end{ex}

% ===== Câu 58 =====
% ID: L11C1B5-14-DS
% Mức: VD
\begin{ex}
% ID: L11C1B5-14-DS
% Mức: VD
Một vật có khối lượng $m = 400$ g thực hiện dao động điều hòa. Đồ thị thế năng theo thời gian $W_t(t)$ được biểu diễn bằng một đường sóng hình sin nâng trục có giá trị cực đại bằng $0,08$ J. Lấy $\pi^2 = 10$.
\choiceTF
{\True Cơ năng dao động của vật bằng $0,08$ J.}
{\True Khoảng thời gian từ lúc thế năng bằng 0 đến lúc thế năng cực đại lần đầu tiên tương ứng một phần tư chu kì dao động $T$ của vật.}
{\True Nếu chu kì biến thiên của thế năng là $0,1$ s thì chu kì dao động của vật là $0,2$ s.}
{Biên độ dao động của vật có giá trị bằng $10$ cm.}
\loigiai{
\begin{itemchoice}
\item (Đúng) Cơ năng dao động của vật bằng $0,08$ J. (Vì): Mệnh đề này đúng vì giá trị cực đại của thế năng chính bằng cơ năng của hệ dao động: $W = 0,08$ J
\item (Đúng) Khoảng thời gian từ lúc thế năng bằng 0 đến lúc thế năng cực đại lần đầu tiên tương ứng một phần tư chu kì dao động $T$ của vật. (Vì): Mệnh đề này đúng vì thế năng bằng 0 ở vị trí cân bằng và cực đại ở vị trí biên, thời gian ngắn nhất đi từ vị trí cân bằng ra biên bằng một phần tư chu kì dao động $T/4$
\item (Đúng) Nếu chu kì biến thiên của thế năng là $0,1$ s thì chu kì dao động của vật là $0,2$ s. (Vì): Mệnh đề này đúng vì chu kì dao động li độ gấp đôi chu kì thế năng: $T = 2T' = 2 \cdot 0,1 = 0,2$ s
\item (Sai) Biên độ dao động của vật có giá trị bằng $10$ cm. (Vì): Mệnh đề này sai vì từ chu kì $T = 0,2$ s ta có tần số góc $\omega = 10\pi$ rad/s. Cơ năng $W = \dfrac{1}{2}m\omega^2 A^2 \Rightarrow 0,08 = 0,5 \cdot 0,4 \cdot (10\pi)^2 \cdot A^2 \Rightarrow 0,08 = 200 A^2 \Rightarrow A^2 = 0,0004 \Rightarrow A = 0,02\,\mathrm{m}= 2$ cm
\end{itemchoice}
}
\end{ex}

% ===== Câu 59 =====
% ID: L11C1B5-14-TL
% Mức: VD
\dangbt{Năng lượng của con lắc đơn dao động điều hòa}
\begin{bt}
% ID: L11C1B5-14-TL
% Mức: VD
Con lắc đơn có $m=0.25$ kg, $l=0.9$ m, biên độ góc 0.2 rad, lấy $g=10$ m/s $^2$. Tính cơ năng trong gần đúng góc nhỏ.
\end{bt}

% ===== Câu 60 =====
% ID: L11C1B5-14-TLN
% Mức: VD
\dangbt{Phân tích đồ thị năng lượng}
\begin{ex}
% ID: L11C1B5-14-TLN
% Mức: VD
Con lắc lò xo có độ cứng $80\,\text{N}$ /m và biên độ $6\,\text{cm}$. Cơ năng theo đơn vị $J$ bằng bao nhiêu?
\shortans{0,144}
\end{ex}

% ===== Câu 61 =====
% ID: L11C1B5-14-TN
% Mức: VD
\dangbt{Chu kì biến thiên của động năng, thế năng}
\begin{ex}
% ID: L11C1B5-14-TN
% Mức: VD
Một chất điểm dao động điều hòa. Biết khoảng thời gian ngắn nhất giữa hai lần liên tiếp động năng bằng thế năng là $0,1\ \mathrm{s}$. Chu kì dao động của chất điểm là:
\choice
{$0,1\ \mathrm{s}$}
{$0,2\ \mathrm{s}$}
{\True $0,4\ \mathrm{s}$}
{$0,8\ \mathrm{s}$}
\loigiai{
Khoảng thời gian ngắn nhất giữa hai lần liên tiếp động năng bằng thế năng là $\Delta t = \dfrac{T}{4}$, do đó
  \begin{aligned}
 $\dfrac{T}{4}$ &= 0,1 
$\RightarrowT$ &= 4 \cdot 0,1 = 0,4 $\\mathrm{s}$.
 \end{aligned}
}
\end{ex}

% ===== Câu 62 =====
% ID: L11C1B5-15-DS
% Mức: VD
\begin{ex}
% ID: L11C1B5-15-DS
% Mức: VD
Một vật có khối lượng $m = 400$ g thực hiện dao động điều hòa. Đồ thị thế năng theo thời gian $W_t(t)$ được biểu diễn bằng một đường sóng hình sin nâng trục có giá trị cực đại bằng $0,08$ J. Lấy $\pi^2 = 10$.
\choiceTF
{\True Cơ năng dao động của vật bằng $0,08$ J.}
{\True Khoảng thời gian từ lúc thế năng bằng 0 đến lúc thế năng cực đại lần đầu tiên tương ứng một phần tư chu kì dao động $T$ của vật.}
{\True Nếu chu kì biến thiên của thế năng là $0,1$ s thì chu kì dao động của vật là $0,2$ s.}
{Biên độ dao động của vật có giá trị bằng $10$ cm.}
\loigiai{
\begin{itemchoice}
\item (Đúng) Cơ năng dao động của vật bằng $0,08$ J. (Vì): Mệnh đề này đúng vì giá trị cực đại của thế năng chính bằng cơ năng của hệ dao động: $W = 0,08$ J
\item (Đúng) Khoảng thời gian từ lúc thế năng bằng 0 đến lúc thế năng cực đại lần đầu tiên tương ứng một phần tư chu kì dao động $T$ của vật. (Vì): Mệnh đề này đúng vì thế năng bằng 0 ở vị trí cân bằng và cực đại ở vị trí biên, thời gian ngắn nhất đi từ vị trí cân bằng ra biên bằng một phần tư chu kì dao động $T/4$
\item (Đúng) Nếu chu kì biến thiên của thế năng là $0,1$ s thì chu kì dao động của vật là $0,2$ s. (Vì): Mệnh đề này đúng vì chu kì dao động li độ gấp đôi chu kì thế năng: $T = 2T' = 2 \cdot 0,1 = 0,2$ s
\item (Sai) Biên độ dao động của vật có giá trị bằng $10$ cm. (Vì): Mệnh đề này sai vì từ chu kì $T = 0,2$ s ta có tần số góc $\omega = 10\pi$ rad/s. Cơ năng $W = \dfrac{1}{2}m\omega^2 A^2 \Rightarrow 0,08 = 0,5 \cdot 0,4 \cdot (10\pi)^2 \cdot A^2 \Rightarrow 0,08 = 200 A^2 \Rightarrow A^2 = 0,0004 \Rightarrow A = 0,02\,\mathrm{m}= 2$ cm
\end{itemchoice}
}
\end{ex}

% ===== Câu 63 =====
% ID: L11C1B5-15-TL
% Mức: VD
\dangbt{Năng lượng của con lắc đơn dao động điều hòa}
\begin{bt}
% ID: L11C1B5-15-TL
% Mức: VD
Con lắc đơn có $m=0.5$ kg, $l=1.6$ m, biên độ góc 0.1 rad, lấy $g=10$ m/s $^2$. Tính cơ năng trong gần đúng góc nhỏ.
\end{bt}

% ===== Câu 64 =====
% ID: L11C1B5-15-TLN
% Mức: VD
\dangbt{Phân tích đồ thị năng lượng}
\begin{ex}
% ID: L11C1B5-15-TLN
% Mức: VD
Con lắc lò xo có độ cứng $100\,\text{N}$ /m và biên độ $8\,\text{cm}$. Cơ năng theo đơn vị $J$ bằng bao nhiêu?
\shortans{0,32}
\end{ex}

% ===== Câu 65 =====
% ID: L11C1B5-15-TN
% Mức: VD
\dangbt{Chu kì biến thiên của động năng, thế năng}
\begin{ex}
% ID: L11C1B5-15-TN
% Mức: VD
Một chất điểm dao động điều hòa. Biết khoảng thời gian ngắn nhất giữa hai lần liên tiếp động năng bằng thế năng là $0,1\ \mathrm{s}$. Chu kì dao động của chất điểm là:
\choice
{$0,1\ \mathrm{s}$}
{$0,2\ \mathrm{s}$}
{\True $0,4\ \mathrm{s}$}
{$0,8\ \mathrm{s}$}
\loigiai{
Khoảng thời gian ngắn nhất giữa hai lần liên tiếp động năng bằng thế năng là $\Delta t = \dfrac{T}{4}$, do đó
  \begin{aligned}
 $\dfrac{T}{4}$ &= 0,1 
$\RightarrowT$ &= 4 \cdot 0,1 = 0,4 $\\mathrm{s}$.
 \end{aligned}
}
\end{ex}

% ===== Câu 66 =====
% ID: L11C1B5-16-DS
% Mức: TH
\dangbt{Nhận diện phát biểu về năng lượng dao động}
\begin{ex}
% ID: L11C1B5-16-DS
% Mức: TH
Một vật dao động điều hòa có chu kì $T$, tại t = 0 vật qua vị trí cân bằng. Hãy xác định tính đúng sai của các phát biểu sau.
\choiceTF
{\True Tại t = 0, động năng đạt cực đại.}
{\True Tại t = 0, thế năng bằng 0.}
{\True Sau thời gian T/4, vật đến một vị trí biên và thế năng đạt cực đại.}
{Cơ năng biến thiên tuần hoàn với chu kì T/2.}
\loigiai{
\begin{itemchoice}
\item (Đúng) Tại t = 0, động năng đạt cực đại.
\item (Đúng) Tại t = 0, thế năng bằng 0.
\item (Đúng) Sau thời gian T/4, vật đến một vị trí biên và thế năng đạt cực đại.
\item (Sai) Cơ năng biến thiên tuần hoàn với chu kì T/2.
\end{itemchoice}
}
\end{ex}

% ===== Câu 67 =====
% ID: L11C1B5-16-TL
% Mức: TH
\dangbt{Phân tích đồ thị năng lượng}
\begin{bt}
% ID: L11C1B5-16-TL
% Mức: TH
Con lắc lò xo có $k=50$ N/m, biên độ $4\,\text{cm}$. Tính cơ năng; động năng và thế năng tại $x=2$ cm.
\end{bt}

% ===== Câu 68 =====
% ID: L11C1B5-16-TLN
% Mức: TH
\dangbt{Sự chuyển hóa năng lượng trong dao động điều hòa}
\begin{ex}
% ID: L11C1B5-16-TLN
% Mức: TH
Con lắc lò xo có độ cứng $50\,\text{N}$ /m và biên độ $4\,\text{cm}$. Cơ năng theo đơn vị $J$ bằng bao nhiêu?
\shortans{0,04}
\end{ex}

% ===== Câu 69 =====
% ID: L11C1B5-16-TN
% Mức: NB
\dangbt{Chưa có dạng}
\begin{ex}
% ID: L11C1B5-16-TN
% Mức: NB
Đại lượng nào sau đây tăng gấp đôi khi biên độ của dao động điều hoà của con lắc lò xo tăng gấp đôi? $\nguon{ly11 kntt.pdf}$
\choice
{Cơ năng của con lắc.}
{Động năng của con lắc.}
{\True Vận tốc cực đại.}
{Thế năng của con lắc.}
\loigiai{
Chọn C.

Biên độ tăng gấp đôi thì cơ năng tăng 4 lần, động năng và thế năng phụ thuộc vào vị trí và vận tốc của vật. Chỉ có vận tốc cực đại v_(max) = Aω tăng gấp đôi.
}
\end{ex}

% ===== Câu 70 =====
% ID: L11C1B5-17-DS
% Mức: TH
\dangbt{Nhận diện phát biểu về năng lượng dao động}
\begin{ex}
% ID: L11C1B5-17-DS
% Mức: TH
Một vật dao động điều hòa có chu kì $T$, tại t = 0 vật qua vị trí cân bằng. Hãy xác định tính đúng sai của các phát biểu sau.
\choiceTF
{\True Tại t = 0, động năng đạt cực đại.}
{\True Tại t = 0, thế năng bằng 0.}
{\True Sau thời gian T/4, vật đến một vị trí biên và thế năng đạt cực đại.}
{Cơ năng biến thiên tuần hoàn với chu kì T/2.}
\loigiai{
\begin{itemchoice}
\item (Đúng) Tại t = 0, động năng đạt cực đại.
\item (Đúng) Tại t = 0, thế năng bằng 0.
\item (Đúng) Sau thời gian T/4, vật đến một vị trí biên và thế năng đạt cực đại.
\item (Sai) Cơ năng biến thiên tuần hoàn với chu kì T/2.
\end{itemchoice}
}
\end{ex}

% ===== Câu 71 =====
% ID: L11C1B5-17-TL
% Mức: VD
\dangbt{Phân tích đồ thị năng lượng}
\begin{bt}
% ID: L11C1B5-17-TL
% Mức: VD
Con lắc lò xo có $k=80$ N/m, biên độ $6\,\text{cm}$. Tính cơ năng; động năng và thế năng tại $x=3$ cm.
\end{bt}

% ===== Câu 72 =====
% ID: L11C1B5-17-TLN
% Mức: VD
\dangbt{Sự chuyển hóa năng lượng trong dao động điều hòa}
\begin{ex}
% ID: L11C1B5-17-TLN
% Mức: VD
Con lắc lò xo có độ cứng $80\,\text{N}$ /m và biên độ $6\,\text{cm}$. Cơ năng theo đơn vị $J$ bằng bao nhiêu?
\shortans{0,144}
\end{ex}

% ===== Câu 73 =====
% ID: L11C1B5-17-TN
% Mức: NB
\dangbt{Chưa có dạng}
\begin{ex}
% ID: L11C1B5-17-TN
% Mức: NB
Trong dao động điều hoà thì tập hợp ba đại lượng nào sau đây không thay đổi theo thời gian? $\nguon{ly11 kntt.pdf}$
\choice
{Lực kéo về; vận tốc; năng lượng toàn phần.}
{Biên độ; tần số góc; gia tốc.}
{Động năng; tần số; lực kéo về.}
{\True Biên độ; tần số góc; năng lượng toàn phần.}
\loigiai{
Chọn D.

Biên độ, tần số góc, năng lượng toàn phần là ba đại lượng không đổi trong dao động điều hoà.
}
\end{ex}

% ===== Câu 74 =====
% ID: L11C1B5-18-DS
% Mức: TH
\dangbt{Nhận diện phát biểu về năng lượng dao động}
\begin{ex}
% ID: L11C1B5-18-DS
% Mức: TH
Một vật dao động điều hòa có chu kì $T$, tại t = 0 vật qua vị trí cân bằng. Hãy xác định tính đúng sai của các phát biểu sau.
\choiceTF
{\True Tại t = 0, động năng đạt cực đại.}
{\True Tại t = 0, thế năng bằng 0.}
{\True Sau thời gian T/4, vật đến một vị trí biên và thế năng đạt cực đại.}
{Cơ năng biến thiên tuần hoàn với chu kì T/2.}
\loigiai{
\begin{itemchoice}
\item (Đúng) Tại t = 0, động năng đạt cực đại.
\item (Đúng) Tại t = 0, thế năng bằng 0.
\item (Đúng) Sau thời gian T/4, vật đến một vị trí biên và thế năng đạt cực đại.
\item (Sai) Cơ năng biến thiên tuần hoàn với chu kì T/2.
\end{itemchoice}
}
\end{ex}

% ===== Câu 75 =====
% ID: L11C1B5-18-TL
% Mức: VD
\dangbt{Phân tích đồ thị năng lượng}
\begin{bt}
% ID: L11C1B5-18-TL
% Mức: VD
Con lắc lò xo có $k=100$ N/m, biên độ $8\,\text{cm}$. Tính cơ năng; động năng và thế năng tại $x=4$ cm.
\end{bt}

% ===== Câu 76 =====
% ID: L11C1B5-18-TLN
% Mức: VD
\dangbt{Sự chuyển hóa năng lượng trong dao động điều hòa}
\begin{ex}
% ID: L11C1B5-18-TLN
% Mức: VD
Con lắc lò xo có độ cứng $100\,\text{N}$ /m và biên độ $8\,\text{cm}$. Cơ năng theo đơn vị $J$ bằng bao nhiêu?
\shortans{0,32}
\end{ex}

% ===== Câu 77 =====
% ID: L11C1B5-18-TN
% Mức: NB
\dangbt{Chưa có dạng}
\begin{ex}
% ID: L11C1B5-18-TN
% Mức: NB
Một chất điểm dao động điều hoà. Biết khoảng thời gian giữa năm lần liên tiếp động năng của chất điểm bằng thế năng của hệ là $0,4\,\mathrm{s}$. Tần số của dao động của chất điểm là $\nguon{ly11 kntt.pdf}$
\choice
{\True $2,5 \mathrm{Hz}$.}
{$3,125 \mathrm{Hz}$.}
{$5 \mathrm{Hz}$.}
{$6,25 \mathrm{Hz}$.}
\loigiai{
Chọn A.

Thời gian giữa năm lần liên tiếp động năng bằng thế năng là:

4. $\dfrac{T}{4}= 0,4\RightarrowT=0,4\,\mathrm{s}\Rightarrowf =\dfrac{1}{0,4}=2,5\,\mathrm{Hz}$
}
\end{ex}

% ===== Câu 78 =====
% ID: L11C1B5-19-DS
% Mức: TH
\dangbt{Nhận diện phát biểu về năng lượng dao động}
\begin{ex}
% ID: L11C1B5-19-DS
% Mức: TH
Chọn Đúng hoặc Sai cho từng ý:
\choiceTF

\loigiai{
A. Đúng — A. $W=\dfrac{1}{2}m\omega^2A^2$ không phụ thuộc thời gian.
B. Sai — B. Không phụ thuộc pha ban đầu.
C. Sai — C. $W_d$ và $W_t$ biến thiên theo thời gian.
D. Đúng — D. $W\propto A^2$.
}
\end{ex}

% ===== Câu 79 =====
% ID: L11C1B5-19-TLN
% Mức: TH
\dangbt{Tính cơ năng của con lắc lò xo từ biên độ dao động}
\begin{ex}
% ID: L11C1B5-19-TLN
% Mức: TH
Một con lắc lò xo có độ cứng $50\,\text{N}$ /m và biên độ $20\,\text{cm}$. Cơ năng của con lắc theo đơn vị $J$ bằng bao nhiêu?
\shortans{1}
\loigiai{
$W = kA²/2 = 50·0,2²/2 = 1 J.$
}
\end{ex}

% ===== Câu 80 =====
% ID: L11C1B5-19-TN
% Mức: NB
\dangbt{Chưa có dạng}
\begin{ex}
% ID: L11C1B5-19-TN
% Mức: NB
Một chất điểm có khối lượng m, dao động điều hoà với biên độ $A$, tần số góc $\omega$. Động năng cực đại của chất điểm là m $\omega$ 2 A_2 $\omega$ 2 A_2 mA $\omega$ _2 m $\omega$ A_2 $\nguon{ly11 kntt.pdf}$
\choice
{\True .}
{$2\,\text{m}$.}
{.}
{. 2 2 2}
\loigiai{
Chọn A.

Động năng cực đại bằng cơ năng: $W_{d\max} = W = \dfrac{1}{2}m\omega^{2}A^{2}$
}
\end{ex}

% ===== Câu 81 =====
% ID: L11C1B5-20-DS
% Mức: VD
\dangbt{Nhận diện phát biểu về năng lượng dao động}
\begin{ex}
% ID: L11C1B5-20-DS
% Mức: VD
Một vật dao động điều hòa có biên độ A. Tại vị trí x = A/2, hãy xác định tính đúng sai của các phát biểu sau.
\choiceTF
{\True Thế năng bằng một phần tư cơ năng.}
{\True Động năng bằng ba phần tư cơ năng.}
{\True Tỉ số động năng và thế năng bằng 3.}
{\True Độ lớn vận tốc bằng √3/2 lần vận tốc cực đại.}
\loigiai{
\begin{itemchoice}
\item (Đúng) Thế năng bằng một phần tư cơ năng.
\item (Đúng) Động năng bằng ba phần tư cơ năng.
\item (Đúng) Tỉ số động năng và thế năng bằng 3.
\item (Đúng) Độ lớn vận tốc bằng √3/2 lần vận tốc cực đại.
\end{itemchoice}
}
\end{ex}

% ===== Câu 82 =====
% ID: L11C1B5-20-TL
% Mức: VD
\dangbt{Sự chuyển hóa năng lượng trong dao động điều hòa}
\begin{bt}
% ID: L11C1B5-20-TL
% Mức: VD
Con lắc lò xo có $k=80$ N/m, biên độ $6\,\text{cm}$. Tính cơ năng; động năng và thế năng tại $x=3$ cm.
\end{bt}

% ===== Câu 83 =====
% ID: L11C1B5-20-TLN
% Mức: TH
\dangbt{Tính cơ năng của con lắc lò xo từ biên độ dao động}
\begin{ex}
% ID: L11C1B5-20-TLN
% Mức: TH
Một con lắc lò xo có độ cứng $50\,\text{N}$ /m và biên độ $20\,\text{cm}$. Cơ năng của con lắc theo đơn vị $J$ bằng bao nhiêu?
\shortans{1}
\loigiai{
$W = kA²/2 = 50·0,2²/2 = 1 J.$
}
\end{ex}

% ===== Câu 84 =====
% ID: L11C1B5-20-TN
% Mức: TH
\dangbt{Chưa có dạng}
\begin{ex}
% ID: L11C1B5-20-TN
% Mức: TH
Phương trình dao động điều hoà của một chất điểm dao động là: 2 $\pix = Acos(\omega t + ) (\mathrm{cm}) 3 Biểu thức động năng của nó biến thiên theo thời gian là mA_2 \omega_2 \pi mA_2 \omega_2 4\pi\nguon{ly11 kntt.pdf}$
\choice
{$W_0 = [1 + \cos (2\omega t + 3 )]$.}
{$Wd = [1 - \cos (2\omega t + 3 )]. 4 4 mA_2\omega_2 4\pimA_2\omega_2$ \pi}
{\True $W_0 = [1 + \cos (2\omega t + 3 )]$.}
{$Wd = [1 - \cos (2\omega t + 3 )]. 4 4$}
\loigiai{
Chọn C.

$W_{d} = \dfrac{1}{2}mv^{2} = \dfrac{1}{2}m\left( {x'} \right)^{2} = \dfrac{1}{2}m\left( {A\omega\cos\left( {\omega t + \dfrac{2\pi}{3} + \dfrac{\pi}{2}} \right)} \right)^{2}= \dfrac{1}{2}mA^{2}\omega^{2}\cos^{2}\left( {\omega t + \dfrac{7\pi}{6}} \right) = \dfrac{mA^{2}\omega^{2}}{4}\left( {1 + \cos\left( {2\omega t + \dfrac{7\pi}{3}} \right)} \right)$
}
\end{ex}

% ===== Câu 85 =====
% ID: L11C1B5-21-DS
% Mức: VD
\dangbt{Nhận diện phát biểu về năng lượng dao động}
\begin{ex}
% ID: L11C1B5-21-DS
% Mức: VD
Một vật dao động điều hòa có biên độ A. Tại vị trí x = A/2, hãy xác định tính đúng sai của các phát biểu sau.
\choiceTF
{\True Thế năng bằng một phần tư cơ năng.}
{\True Động năng bằng ba phần tư cơ năng.}
{\True Tỉ số động năng và thế năng bằng 3.}
{\True Độ lớn vận tốc bằng √3/2 lần vận tốc cực đại.}
\loigiai{
\begin{itemchoice}
\item (Đúng) Thế năng bằng một phần tư cơ năng.
\item (Đúng) Động năng bằng ba phần tư cơ năng.
\item (Đúng) Tỉ số động năng và thế năng bằng 3.
\item (Đúng) Độ lớn vận tốc bằng √3/2 lần vận tốc cực đại.
\end{itemchoice}
}
\end{ex}

% ===== Câu 86 =====
% ID: L11C1B5-21-TL
% Mức: VD
\dangbt{Sự chuyển hóa năng lượng trong dao động điều hòa}
\begin{bt}
% ID: L11C1B5-21-TL
% Mức: VD
Con lắc lò xo có $k=100$ N/m, biên độ $8\,\text{cm}$. Tính cơ năng; động năng và thế năng tại $x=4$ cm.
\end{bt}

% ===== Câu 87 =====
% ID: L11C1B5-21-TLN
% Mức: TH
\dangbt{Tính cơ năng của con lắc lò xo từ biên độ dao động}
\begin{ex}
% ID: L11C1B5-21-TLN
% Mức: TH
Một con lắc lò xo có độ cứng $50\,\text{N}$ /m và biên độ $20\,\text{cm}$. Cơ năng của con lắc theo đơn vị $J$ bằng bao nhiêu?
\shortans{1}
\loigiai{
$W = kA²/2 = 50·0,2²/2 = 1 J.$
}
\end{ex}

% ===== Câu 88 =====
% ID: L11C1B5-21-TN
% Mức: TH
\dangbt{Nhận diện phát biểu về năng lượng dao động}
\begin{ex}
% ID: L11C1B5-21-TN
% Mức: TH
Khi bỏ qua mọi lực cản, đại lượng nào sau đây không đổi trong suốt quá trình dao động điều hòa?
\choice
{\True Cơ năng}
{Động năng}
{Thế năng}
{Vận tốc}
\loigiai{
Khi không có lực cản, cơ năng được bảo toàn; động năng, thế năng và vận tốc đều biến thiên.
}
\end{ex}

% ===== Câu 89 =====
% ID: L11C1B5-22-DS
% Mức: VD
\begin{ex}
% ID: L11C1B5-22-DS
% Mức: VD
Một vật dao động điều hòa có biên độ A. Tại vị trí x = A/2, hãy xác định tính đúng sai của các phát biểu sau.
\choiceTF
{\True Thế năng bằng một phần tư cơ năng.}
{\True Động năng bằng ba phần tư cơ năng.}
{\True Tỉ số động năng và thế năng bằng 3.}
{\True Độ lớn vận tốc bằng √3/2 lần vận tốc cực đại.}
\loigiai{
\begin{itemchoice}
\item (Đúng) Thế năng bằng một phần tư cơ năng.
\item (Đúng) Động năng bằng ba phần tư cơ năng.
\item (Đúng) Tỉ số động năng và thế năng bằng 3.
\item (Đúng) Độ lớn vận tốc bằng √3/2 lần vận tốc cực đại.
\end{itemchoice}
}
\end{ex}

% ===== Câu 90 =====
% ID: L11C1B5-22-TL
% Mức: TH
\dangbt{Tính cơ năng của con lắc lò xo từ biên độ dao động}
\begin{bt}
% ID: L11C1B5-22-TL
% Mức: TH
Con lắc lò xo có $k=50$ N/m, biên độ $4\,\text{cm}$. Tính cơ năng; động năng và thế năng tại $x=2$ cm.
\end{bt}

% ===== Câu 91 =====
% ID: L11C1B5-22-TLN
% Mức: TH
\begin{ex}
% ID: L11C1B5-22-TLN
% Mức: TH
Một con lắc lò xo có độ cứng $80\,\text{N}$ /m và biên độ $10\,\text{cm}$. Cơ năng của con lắc theo đơn vị $J$ bằng bao nhiêu?
\shortans{0,4}
\loigiai{
$W = kA²/2 = 80·0,1²/2 = 0,4 J.$
}
\end{ex}

% ===== Câu 92 =====
% ID: L11C1B5-22-TN
% Mức: TH
\dangbt{Nhận diện phát biểu về năng lượng dao động}
\begin{ex}
% ID: L11C1B5-22-TN
% Mức: TH
Khi bỏ qua mọi lực cản, đại lượng nào sau đây không đổi trong suốt quá trình dao động điều hòa?
\choice
{\True Cơ năng}
{Động năng}
{Thế năng}
{Vận tốc}
\loigiai{
Khi không có lực cản, cơ năng được bảo toàn; động năng, thế năng và vận tốc đều biến thiên.
}
\end{ex}

% ===== Câu 93 =====
% ID: L11C1B5-23-DS
% Mức: TH
\dangbt{Năng lượng của con lắc đơn dao động điều hòa}
\begin{ex}
% ID: L11C1B5-23-DS
% Mức: TH
Xét các đặc điểm về mặt năng lượng của con lắc lò xo và con lắc đơn dao động điều hòa trong các điều kiện tiêu chuẩn:
\choiceTF
{\True Cơ năng của con lắc lò xo dao động điều hòa tỉ lệ thuận với bình phương biên độ dao động.}
{\True Khi góc lệch của con lắc đơn nhỏ ($\alpha_0 \le 10^\circ$), thế năng của con lắc đơn được tính xấp xỉ bằng công thức $W_t \approx \dfrac{1}{2}mgl\alpha^2$ (với góc lệch $\alpha$ tính bằng radian).}
{Tần số biến thiên của thế năng con lắc lò xo tỉ lệ nghịch với khối lượng của vật nặng treo vào lò xo.}
{\True Cơ năng của con lắc đơn dao động điều hòa tỉ lệ thuận với khối lượng của quả cầu treo.}
\loigiai{
\begin{itemchoice}
\item (Đúng) Cơ năng của con lắc lò xo dao động điều hòa tỉ lệ thuận với bình phương biên độ dao động. (Vì): Mệnh đề này đúng vì cơ năng con lắc lò xo được xác định theo hệ thức: $W = \dfrac{1}{2}kA^2$
\item (Đúng) Khi góc lệch của con lắc đơn nhỏ ($\alpha_0 \le 10^\circ$), thế năng của con lắc đơn được tính xấp xỉ bằng công thức $W_t \approx \dfrac{1}{2}mgl\alpha^2$ (với góc lệch $\alpha$ tính bằng radian). (Vì): Mệnh đề này đúng vì thế năng trọng trường $W_t = mgl(1 - \cos\alpha)$, với $\alpha$ nhỏ ta có xấp xỉ $1 - \cos\alpha \approx \dfrac{\alpha^2}{2}$, suy ra $W_t \approx \dfrac{1}{2}mgl\alpha^2$
\item (Sai) Tần số biến thiên của thế năng con lắc lò xo tỉ lệ nghịch với khối lượng của vật nặng treo vào lò xo. (Vì): Mệnh đề này sai vì tần số biến thiên của thế năng là $f' = 2f = \dfrac{1}{\pi}\sqrt{\dfrac{k}{m}}$, tỉ lệ nghịch với căn bậc hai của khối lượng vật nặng chứ không phải tỉ lệ nghịch đơn thuần
\item (Đúng) Cơ năng của con lắc đơn dao động điều hòa tỉ lệ thuận với khối lượng của quả cầu treo.
\end{itemchoice}
}
\end{ex}

% ===== Câu 94 =====
% ID: L11C1B5-23-TL
% Mức: VD
\dangbt{Tính cơ năng của con lắc lò xo từ biên độ dao động}
\begin{bt}
% ID: L11C1B5-23-TL
% Mức: VD
Con lắc lò xo có $k=80$ N/m, biên độ $6\,\text{cm}$. Tính cơ năng; động năng và thế năng tại $x=3$ cm.
\end{bt}

% ===== Câu 95 =====
% ID: L11C1B5-23-TLN
% Mức: TH
\begin{ex}
% ID: L11C1B5-23-TLN
% Mức: TH
Một con lắc lò xo có độ cứng $80\,\text{N}$ /m và biên độ $10\,\text{cm}$. Cơ năng của con lắc theo đơn vị $J$ bằng bao nhiêu?
\shortans{0,4}
\loigiai{
$W = kA²/2 = 80·0,1²/2 = 0,4 J.$
}
\end{ex}

% ===== Câu 96 =====
% ID: L11C1B5-23-TN
% Mức: TH
\dangbt{Nhận diện phát biểu về năng lượng dao động}
\begin{ex}
% ID: L11C1B5-23-TN
% Mức: TH
Khi bỏ qua mọi lực cản, đại lượng nào sau đây không đổi trong suốt quá trình dao động điều hòa?
\choice
{\True Cơ năng}
{Động năng}
{Thế năng}
{Vận tốc}
\loigiai{
Khi không có lực cản, cơ năng được bảo toàn; động năng, thế năng và vận tốc đều biến thiên.
}
\end{ex}

% ===== Câu 97 =====
% ID: L11C1B5-24-DS
% Mức: TH
\dangbt{Năng lượng của con lắc đơn dao động điều hòa}
\begin{ex}
% ID: L11C1B5-24-DS
% Mức: TH
Xét các đặc điểm về mặt năng lượng của con lắc lò xo và con lắc đơn dao động điều hòa trong các điều kiện tiêu chuẩn:
\choiceTF
{\True Cơ năng của con lắc lò xo dao động điều hòa tỉ lệ thuận với bình phương biên độ dao động.}
{\True Khi góc lệch của con lắc đơn nhỏ ($\alpha_0 \le 10^\circ$), thế năng của con lắc đơn được tính xấp xỉ bằng công thức $W_t \approx \dfrac{1}{2}mgl\alpha^2$ (với góc lệch $\alpha$ tính bằng radian).}
{Tần số biến thiên của thế năng con lắc lò xo tỉ lệ nghịch với khối lượng của vật nặng treo vào lò xo.}
{\True Cơ năng của con lắc đơn dao động điều hòa tỉ lệ thuận với khối lượng của quả cầu treo.}
\loigiai{
\begin{itemchoice}
\item (Đúng) Cơ năng của con lắc lò xo dao động điều hòa tỉ lệ thuận với bình phương biên độ dao động. (Vì): Mệnh đề này đúng vì cơ năng con lắc lò xo được xác định theo hệ thức: $W = \dfrac{1}{2}kA^2$
\item (Đúng) Khi góc lệch của con lắc đơn nhỏ ($\alpha_0 \le 10^\circ$), thế năng của con lắc đơn được tính xấp xỉ bằng công thức $W_t \approx \dfrac{1}{2}mgl\alpha^2$ (với góc lệch $\alpha$ tính bằng radian). (Vì): Mệnh đề này đúng vì thế năng trọng trường $W_t = mgl(1 - \cos\alpha)$, với $\alpha$ nhỏ ta có xấp xỉ $1 - \cos\alpha \approx \dfrac{\alpha^2}{2}$, suy ra $W_t \approx \dfrac{1}{2}mgl\alpha^2$
\item (Sai) Tần số biến thiên của thế năng con lắc lò xo tỉ lệ nghịch với khối lượng của vật nặng treo vào lò xo. (Vì): Mệnh đề này sai vì tần số biến thiên của thế năng là $f' = 2f = \dfrac{1}{\pi}\sqrt{\dfrac{k}{m}}$, tỉ lệ nghịch với căn bậc hai của khối lượng vật nặng chứ không phải tỉ lệ nghịch đơn thuần
\item (Đúng) Cơ năng của con lắc đơn dao động điều hòa tỉ lệ thuận với khối lượng của quả cầu treo.
\end{itemchoice}
}
\end{ex}

% ===== Câu 98 =====
% ID: L11C1B5-24-TL
% Mức: VD
\dangbt{Tính cơ năng của con lắc lò xo từ biên độ dao động}
\begin{bt}
% ID: L11C1B5-24-TL
% Mức: VD
Con lắc lò xo có $k=100$ N/m, biên độ $8\,\text{cm}$. Tính cơ năng; động năng và thế năng tại $x=4$ cm.
\end{bt}

% ===== Câu 99 =====
% ID: L11C1B5-24-TLN
% Mức: TH
\begin{ex}
% ID: L11C1B5-24-TLN
% Mức: TH
Một con lắc lò xo có độ cứng $80\,\text{N}$ /m và biên độ $10\,\text{cm}$. Cơ năng của con lắc theo đơn vị $J$ bằng bao nhiêu?
\shortans{0,4}
\loigiai{
$W = kA²/2 = 80·0,1²/2 = 0,4 J.$
}
\end{ex}

% ===== Câu 100 =====
% ID: L11C1B5-24-TN
% Mức: TH
\dangbt{Nhận diện phát biểu về năng lượng dao động}
\begin{ex}
% ID: L11C1B5-24-TN
% Mức: TH
Xét cơ năng của một dao động điều hòa thì:
\choice
{Động năng biến thiên tuần hoàn theo thời gian với tần số bằng tần số của dao động}
{Thế năng tỉ lệ thuận với li độ}
{Cơ năng tỉ lệ với biên độ}
{\True Tổng động năng và thế năng là một số không đổi}
\loigiai{
Cơ năng của một dao động điều hòa được định nghĩa là tổng động năng và thế năng của vật. Trong một hệ dao động điều hòa lý tưởng (không có ma sát hay lực cản), cơ năng được bảo toàn, tức là một số không đổi. Công thức cơ năng là $W = \frac{1}{2}kA^2 = \frac{1}{2}m\omega^2A^2$.
  - Phương án $A$ đúng: Cơ năng là tổng động năng và thế năng, và nó được bảo toàn trong dao động điều hòa.
  - Phương án $B$ sai: Thế năng $W_t = \frac{1}{2}kx^2$, tỉ lệ thuận với bình phương li độ $x^2$, không phải li độ $x$.
  - Phương án $C$ sai: Động năng $W_\text{đ} = \frac{1}{2}mv^2 = \frac{1}{2}m(\omega A \cos(\omega t + \phi))^2 = \frac{1}{4}m\omega^2 A^2 (1+\cos(2(\omega t + \phi)))$. Động năng biến thiên tuần hoàn với tần số $2\omega$, gấp đôi tần số dao động $\omega$.
  - Phương án $D$ sai: Cơ năng $W = \frac{1}{2}kA^2$, tỉ lệ thuận với bình phương biên độ $A^2$, không phải biên độ $A$.
}
\end{ex}

% ===== Câu 101 =====
% ID: L11C1B5-25-DS
% Mức: TH
\dangbt{Năng lượng của con lắc đơn dao động điều hòa}
\begin{ex}
% ID: L11C1B5-25-DS
% Mức: TH
Xét các đặc điểm về mặt năng lượng của con lắc lò xo và con lắc đơn dao động điều hòa trong các điều kiện tiêu chuẩn:
\choiceTF
{\True Cơ năng của con lắc lò xo dao động điều hòa tỉ lệ thuận với bình phương biên độ dao động.}
{\True Khi góc lệch của con lắc đơn nhỏ ($\alpha_0 \le 10^\circ$), thế năng của con lắc đơn được tính xấp xỉ bằng công thức $W_t \approx \dfrac{1}{2}mgl\alpha^2$ (với góc lệch $\alpha$ tính bằng radian).}
{Tần số biến thiên của thế năng con lắc lò xo tỉ lệ nghịch với khối lượng của vật nặng treo vào lò xo.}
{\True Cơ năng của con lắc đơn dao động điều hòa tỉ lệ thuận với khối lượng của quả cầu treo.}
\loigiai{
\begin{itemchoice}
\item (Đúng) Cơ năng của con lắc lò xo dao động điều hòa tỉ lệ thuận với bình phương biên độ dao động. (Vì): Mệnh đề này đúng vì cơ năng con lắc lò xo được xác định theo hệ thức: $W = \dfrac{1}{2}kA^2$
\item (Đúng) Khi góc lệch của con lắc đơn nhỏ ($\alpha_0 \le 10^\circ$), thế năng của con lắc đơn được tính xấp xỉ bằng công thức $W_t \approx \dfrac{1}{2}mgl\alpha^2$ (với góc lệch $\alpha$ tính bằng radian). (Vì): Mệnh đề này đúng vì thế năng trọng trường $W_t = mgl(1 - \cos\alpha)$, với $\alpha$ nhỏ ta có xấp xỉ $1 - \cos\alpha \approx \dfrac{\alpha^2}{2}$, suy ra $W_t \approx \dfrac{1}{2}mgl\alpha^2$
\item (Sai) Tần số biến thiên của thế năng con lắc lò xo tỉ lệ nghịch với khối lượng của vật nặng treo vào lò xo. (Vì): Mệnh đề này sai vì tần số biến thiên của thế năng là $f' = 2f = \dfrac{1}{\pi}\sqrt{\dfrac{k}{m}}$, tỉ lệ nghịch với căn bậc hai của khối lượng vật nặng chứ không phải tỉ lệ nghịch đơn thuần
\item (Đúng) Cơ năng của con lắc đơn dao động điều hòa tỉ lệ thuận với khối lượng của quả cầu treo.
\end{itemchoice}
}
\end{ex}

% ===== Câu 102 =====
% ID: L11C1B5-25-TL
% Mức: TH
\dangbt{Tính cơ năng của con lắc đơn}
\begin{bt}
% ID: L11C1B5-25-TL
% Mức: TH
Con lắc đơn có $m=0.2$ kg, $l=1$ m, biên độ góc 0.1 rad, lấy $g=10$ m/s $^2$. Tính cơ năng trong gần đúng góc nhỏ.
\end{bt}

% ===== Câu 103 =====
% ID: L11C1B5-25-TLN
% Mức: VD
\begin{ex}
% ID: L11C1B5-25-TLN
% Mức: VD
Một con lắc đơn có vật nặng $0,4\,\text{kg}$, khi qua vị trí cân bằng đạt tốc độ $0,5\,\text{m}$ /s. Cơ năng của con lắc theo đơn vị $J$ bằng bao nhiêu?
\shortans{0,05}
\loigiai{
Tại vị trí cân bằng, $W$ = mvmax²/2 = 0,4·0,5²/2 = $0,05\,\text{J}$.
}
\end{ex}

% ===== Câu 104 =====
% ID: L11C1B5-25-TN
% Mức: TH
\dangbt{Nhận diện phát biểu về năng lượng dao động}
\begin{ex}
% ID: L11C1B5-25-TN
% Mức: TH
Một dao động điều hòa có chu kỳ $T$ và tần số f. Chọn phát biểu sai:
\choice
{Tổng động năng và thế năng là một số không đổi}
{Động năng của vật biến thiên tuần hoàn với tần số $f'=2f$}
{\True Cơ năng của vật biến thiên tuần hoàn với tần số $f'=2f$}
{Thế năng biến thiên tuần hoàn với chu kỳ $T'=T/2$}
\loigiai{
Trong dao động điều hòa, cơ năng của vật được bảo toàn (là một hằng số) nếu bỏ qua ma sát. Do đó, cơ năng không biến thiên tuần hoàn. Động năng và thế năng của vật biến thiên tuần hoàn với tần số $f' = 2f$ và chu kỳ $T' = T/2$. Tổng động năng và thế năng chính là cơ năng và là một số không đổi. Vậy, phát biểu "Cơ năng của vật biến thiên tuần hoàn với tần số $f' = 2f$ " là sai.
}
\end{ex}

% ===== Câu 105 =====
% ID: L11C1B5-26-DS
% Mức: TH
\dangbt{Năng lượng của con lắc đơn dao động điều hòa}
\begin{ex}
% ID: L11C1B5-26-DS
% Mức: TH
[SGK CTST]
 Chọn Đúng hoặc Sai cho các phát biểu sau về năng lượng trong dao động điều hòa:
\choiceTF
{\True Khi vật ở vị trí cân bằng thì động năng cực đại}
{\True Khi vật ở vị trí biên thì thế năng cực đại}
{Khi vật ở vị trí cân bằng thì thế năng cực đại}
{Khi vật ở vị trí biên thì động năng cực đại}
\loigiai{
\begin{itemchoice}
\item (Đúng) Khi vật ở vị trí cân bằng thì động năng cực đại. (Vì): Ở vị trí cân bằng ($x=0$), tốc độ cực đại $\Rightarrow$ động năng cực đại, thế năng bằng 0
\item (Đúng) Khi vật ở vị trí biên thì thế năng cực đại. (Vì): Ở vị trí biên ($x=\pm A$), tốc độ bằng 0, thế năng cực đại
\item (Sai) Khi vật ở vị trí cân bằng thì thế năng cực đại. (Vì): Thế năng tại cân bằng là 0, không phải cực đại
\item (Sai) Khi vật ở vị trí biên thì động năng cực đại. (Vì): Ở biên, vận tốc bằng 0 nên động năng bằng 0
\end{itemchoice}
}
\end{ex}

% ===== Câu 106 =====
% ID: L11C1B5-26-TL
% Mức: VD
\dangbt{Tính cơ năng của con lắc đơn}
\begin{bt}
% ID: L11C1B5-26-TL
% Mức: VD
Con lắc đơn có $m=0.25$ kg, $l=0.9$ m, biên độ góc 0.2 rad, lấy $g=10$ m/s $^2$. Tính cơ năng trong gần đúng góc nhỏ.
\end{bt}

% ===== Câu 107 =====
% ID: L11C1B5-26-TLN
% Mức: VD
\begin{ex}
% ID: L11C1B5-26-TLN
% Mức: VD
Một con lắc đơn có vật nặng $0,4\,\text{kg}$, khi qua vị trí cân bằng đạt tốc độ $0,5\,\text{m}$ /s. Cơ năng của con lắc theo đơn vị $J$ bằng bao nhiêu?
\shortans{0,05}
\loigiai{
Tại vị trí cân bằng, $W$ = mvmax²/2 = 0,4·0,5²/2 = $0,05\,\text{J}$.
}
\end{ex}

% ===== Câu 108 =====
% ID: L11C1B5-26-TN
% Mức: NB
\dangbt{Năng lượng của con lắc đơn dao động điều hòa}
\begin{ex}
% ID: L11C1B5-26-TN
% Mức: NB
Một con lắc đơn gồm vật nhỏ khối lượng $m$, treo tại nơi có gia tốc trọng trường $g$. Khi con lắc dao động điều hòa với biên độ góc $\alpha_0$ (tính bằng radian), mốc thế năng chọn tại vị trí cân bằng, cơ năng của con lắc được xác định bằng công thức:
\choice
{$W = mgl(1 + \cos\alpha_0)$}
{\True $W = \dfrac{1}{2}mgl\alpha_0^2$}
{$W = \dfrac{1}{2}mg\alpha_0^2$}
{$W = mgl\alpha_0^2$}
\loigiai{
• Cơ năng của con lắc đơn bằng thế năng cực đại ở vị trí biên:
 $W = mgl(1 - \cos\alpha_0).$
 
• Với dao động điều hòa, biên độ góc $\alpha_0$ nhỏ $\alpha_0 \le 10^\circ)$, áp dụng công thức xấp xỉ:
 $1 - \cos\alpha_0 \approx \dfrac{\alpha_0^2}{2}.$
 
• Thay vào công thức ban đầu:
 $W = \dfrac{1}{2}mgl\alpha_0^2.$
}
\end{ex}

% ===== Câu 109 =====
% ID: L11C1B5-27-DS
% Mức: VD
\begin{ex}
% ID: L11C1B5-27-DS
% Mức: VD
Một con lắc đơn gồm quả cầu nhỏ khối lượng $m = 100$ g, treo vào một sợi dây mảnh có chiều dài $l = 1$ m tại nơi có gia tốc trọng trường $g = 10$ m/s $^2$. Kéo con lắc lệch khỏi phương thẳng đứng góc $\alpha_0 = 0,1$ rad rồi thả nhẹ cho dao động điều hòa. Chọn mốc thế năng tại vị trí cân bằng.
\choiceTF
{\True Tần số góc dao động điều hòa của con lắc đơn xấp xỉ bằng $\pi$ rad/s.}
{\True Cơ năng dao động của con lắc đơn có giá trị bằng $0,005$ J.}
{Tại vị trí có li độ góc $\alpha = 0,05$ rad, thế năng trọng trường của con lắc đơn bằng $0,0025$ J.}
{\True Động năng của vật nặng tại vị trí có li độ góc $\alpha = 0,05$ rad bằng $0,00375$ J.}
\loigiai{
\begin{itemchoice}
\item (Đúng) Tần số góc dao động điều hòa của con lắc đơn xấp xỉ bằng $\pi$ rad/s. (Vì): Mệnh đề này đúng vì tần số góc dao động của con lắc đơn là $\omega = \sqrt{\dfrac{g}{l}} = \sqrt{\dfrac{10}{1}} = \sqrt{10} \approx \pi$ rad/s
\item (Đúng) Cơ năng dao động của con lắc đơn có giá trị bằng $0,005$ J. (Vì): Mệnh đề này đúng vì cơ năng của con lắc là $W = \dfrac{1}{2}mgl\alpha_0^2 = 0,5 \cdot 0,1 \cdot 10 \cdot 1 \cdot 0,1^2 = 0,005$ J
\item (Sai) Tại vị trí có li độ góc $\alpha = 0,05$ rad, thế năng trọng trường của con lắc đơn bằng $0,0025$ J. (Vì): Mệnh đề này sai vì thế năng tại góc lệch $\alpha = 0,05$ rad là $W_t = \dfrac{1}{2}mgl\alpha^2 = 0,5 \cdot 0,1 \cdot 10 \cdot 1 \cdot 0,05^2 = 0,00125$ J
\item (Đúng) Động năng của vật nặng tại vị trí có li độ góc $\alpha = 0,05$ rad bằng $0,00375$ J. (Vì): Mệnh đề này đúng vì động năng của quả cầu bằng cơ năng trừ đi thế năng tương ứng: $W_d = W - W_t = 0,005 - 0,00125 = 0,00375$ J
\end{itemchoice}
}
\end{ex}

% ===== Câu 110 =====
% ID: L11C1B5-27-TL
% Mức: VD
\dangbt{Tính cơ năng của con lắc đơn}
\begin{bt}
% ID: L11C1B5-27-TL
% Mức: VD
Con lắc đơn có $m=0.5$ kg, $l=1.6$ m, biên độ góc 0.1 rad, lấy $g=10$ m/s $^2$. Tính cơ năng trong gần đúng góc nhỏ.
\end{bt}

% ===== Câu 111 =====
% ID: L11C1B5-27-TLN
% Mức: VD
\begin{ex}
% ID: L11C1B5-27-TLN
% Mức: VD
Một con lắc đơn có vật nặng $0,4\,\text{kg}$, khi qua vị trí cân bằng đạt tốc độ $0,5\,\text{m}$ /s. Cơ năng của con lắc theo đơn vị $J$ bằng bao nhiêu?
\shortans{0,05}
\loigiai{
Tại vị trí cân bằng, $W$ = mvmax²/2 = 0,4·0,5²/2 = $0,05\,\text{J}$.
}
\end{ex}

% ===== Câu 112 =====
% ID: L11C1B5-27-TN
% Mức: NB
\dangbt{Năng lượng của con lắc đơn dao động điều hòa}
\begin{ex}
% ID: L11C1B5-27-TN
% Mức: NB
Một con lắc đơn gồm vật nhỏ khối lượng $m$, treo tại nơi có gia tốc trọng trường $g$. Khi con lắc dao động điều hòa với biên độ góc $\alpha_0$ (tính bằng radian), mốc thế năng chọn tại vị trí cân bằng, cơ năng của con lắc được xác định bằng công thức:
\choice
{$W = mgl(1 + \cos\alpha_0)$}
{\True $W = \dfrac{1}{2}mgl\alpha_0^2$}
{$W = \dfrac{1}{2}mg\alpha_0^2$}
{$W = mgl\alpha_0^2$}
\loigiai{
• Cơ năng của con lắc đơn bằng thế năng cực đại ở vị trí biên:
 $W = mgl(1 - \cos\alpha_0).$
 
• Với dao động điều hòa, biên độ góc $\alpha_0$ nhỏ $\alpha_0 \le 10^\circ)$, áp dụng công thức xấp xỉ:
 $1 - \cos\alpha_0 \approx \dfrac{\alpha_0^2}{2}.$
 
• Thay vào công thức ban đầu:
 $W = \dfrac{1}{2}mgl\alpha_0^2.$
}
\end{ex}

% ===== Câu 113 =====
% ID: L11C1B5-28-DS
% Mức: VD
\begin{ex}
% ID: L11C1B5-28-DS
% Mức: VD
Một con lắc đơn gồm quả cầu nhỏ khối lượng $m = 100$ g, treo vào một sợi dây mảnh có chiều dài $l = 1$ m tại nơi có gia tốc trọng trường $g = 10$ m/s $^2$. Kéo con lắc lệch khỏi phương thẳng đứng góc $\alpha_0 = 0,1$ rad rồi thả nhẹ cho dao động điều hòa. Chọn mốc thế năng tại vị trí cân bằng.
\choiceTF
{\True Tần số góc dao động điều hòa của con lắc đơn xấp xỉ bằng $\pi$ rad/s.}
{\True Cơ năng dao động của con lắc đơn có giá trị bằng $0,005$ J.}
{Tại vị trí có li độ góc $\alpha = 0,05$ rad, thế năng trọng trường của con lắc đơn bằng $0,0025$ J.}
{\True Động năng của vật nặng tại vị trí có li độ góc $\alpha = 0,05$ rad bằng $0,00375$ J.}
\loigiai{
\begin{itemchoice}
\item (Đúng) Tần số góc dao động điều hòa của con lắc đơn xấp xỉ bằng $\pi$ rad/s. (Vì): Mệnh đề này đúng vì tần số góc dao động của con lắc đơn là $\omega = \sqrt{\dfrac{g}{l}} = \sqrt{\dfrac{10}{1}} = \sqrt{10} \approx \pi$ rad/s
\item (Đúng) Cơ năng dao động của con lắc đơn có giá trị bằng $0,005$ J. (Vì): Mệnh đề này đúng vì cơ năng của con lắc là $W = \dfrac{1}{2}mgl\alpha_0^2 = 0,5 \cdot 0,1 \cdot 10 \cdot 1 \cdot 0,1^2 = 0,005$ J
\item (Sai) Tại vị trí có li độ góc $\alpha = 0,05$ rad, thế năng trọng trường của con lắc đơn bằng $0,0025$ J. (Vì): Mệnh đề này sai vì thế năng tại góc lệch $\alpha = 0,05$ rad là $W_t = \dfrac{1}{2}mgl\alpha^2 = 0,5 \cdot 0,1 \cdot 10 \cdot 1 \cdot 0,05^2 = 0,00125$ J
\item (Đúng) Động năng của vật nặng tại vị trí có li độ góc $\alpha = 0,05$ rad bằng $0,00375$ J. (Vì): Mệnh đề này đúng vì động năng của quả cầu bằng cơ năng trừ đi thế năng tương ứng: $W_d = W - W_t = 0,005 - 0,00125 = 0,00375$ J
\end{itemchoice}
}
\end{ex}

% ===== Câu 114 =====
% ID: L11C1B5-28-TL
% Mức: TH
\dangbt{Tính động năng, thế năng và cơ năng tại một li độ hoặc thời điểm xác định}
\begin{bt}
% ID: L11C1B5-28-TL
% Mức: TH
Con lắc lò xo có $k=50$ N/m, biên độ $4\,\text{cm}$. Tính cơ năng; động năng và thế năng tại $x=2$ cm.
\end{bt}

% ===== Câu 115 =====
% ID: L11C1B5-28-TLN
% Mức: VD
\dangbt{Tính cơ năng của con lắc đơn}
\begin{ex}
% ID: L11C1B5-28-TLN
% Mức: VD
Một con lắc đơn có vật nặng $0,5\,\text{kg}$, khi qua vị trí cân bằng đạt tốc độ $0,4\,\text{m}$ /s. Cơ năng của con lắc theo đơn vị $J$ bằng bao nhiêu?
\shortans{0,04}
\loigiai{
$W = mvmax²/2 = 0,5·0,4²/2 = 0,04 J.$
}
\end{ex}

% ===== Câu 116 =====
% ID: L11C1B5-28-TN
% Mức: NB
\dangbt{Năng lượng của con lắc đơn dao động điều hòa}
\begin{ex}
% ID: L11C1B5-28-TN
% Mức: NB
Một con lắc đơn gồm vật nhỏ khối lượng $m$, treo tại nơi có gia tốc trọng trường $g$. Khi con lắc dao động điều hòa với biên độ góc $\alpha_0$ (tính bằng radian), mốc thế năng chọn tại vị trí cân bằng, cơ năng của con lắc được xác định bằng công thức:
\choice
{$W = mgl(1 + \cos\alpha_0)$}
{\True $W = \dfrac{1}{2}mgl\alpha_0^2$}
{$W = \dfrac{1}{2}mg\alpha_0^2$}
{$W = mgl\alpha_0^2$}
\loigiai{
• Cơ năng của con lắc đơn bằng thế năng cực đại ở vị trí biên:
 $W = mgl(1 - \cos\alpha_0).$
 
• Với dao động điều hòa, biên độ góc $\alpha_0$ nhỏ $\alpha_0 \le 10^\circ)$, áp dụng công thức xấp xỉ:
 $1 - \cos\alpha_0 \approx \dfrac{\alpha_0^2}{2}.$
 
• Thay vào công thức ban đầu:
 $W = \dfrac{1}{2}mgl\alpha_0^2.$
}
\end{ex}

% ===== Câu 117 =====
% ID: L11C1B5-29-DS
% Mức: VD
\begin{ex}
% ID: L11C1B5-29-DS
% Mức: VD
Một con lắc đơn gồm quả cầu nhỏ khối lượng $m = 100$ g, treo vào một sợi dây mảnh có chiều dài $l = 1$ m tại nơi có gia tốc trọng trường $g = 10$ m/s $^2$. Kéo con lắc lệch khỏi phương thẳng đứng góc $\alpha_0 = 0,1$ rad rồi thả nhẹ cho dao động điều hòa. Chọn mốc thế năng tại vị trí cân bằng.
\choiceTF
{\True Tần số góc dao động điều hòa của con lắc đơn xấp xỉ bằng $\pi$ rad/s.}
{\True Cơ năng dao động của con lắc đơn có giá trị bằng $0,005$ J.}
{Tại vị trí có li độ góc $\alpha = 0,05$ rad, thế năng trọng trường của con lắc đơn bằng $0,0025$ J.}
{\True Động năng của vật nặng tại vị trí có li độ góc $\alpha = 0,05$ rad bằng $0,00375$ J.}
\loigiai{
\begin{itemchoice}
\item (Đúng) Tần số góc dao động điều hòa của con lắc đơn xấp xỉ bằng $\pi$ rad/s. (Vì): Mệnh đề này đúng vì tần số góc dao động của con lắc đơn là $\omega = \sqrt{\dfrac{g}{l}} = \sqrt{\dfrac{10}{1}} = \sqrt{10} \approx \pi$ rad/s
\item (Đúng) Cơ năng dao động của con lắc đơn có giá trị bằng $0,005$ J. (Vì): Mệnh đề này đúng vì cơ năng của con lắc là $W = \dfrac{1}{2}mgl\alpha_0^2 = 0,5 \cdot 0,1 \cdot 10 \cdot 1 \cdot 0,1^2 = 0,005$ J
\item (Sai) Tại vị trí có li độ góc $\alpha = 0,05$ rad, thế năng trọng trường của con lắc đơn bằng $0,0025$ J. (Vì): Mệnh đề này sai vì thế năng tại góc lệch $\alpha = 0,05$ rad là $W_t = \dfrac{1}{2}mgl\alpha^2 = 0,5 \cdot 0,1 \cdot 10 \cdot 1 \cdot 0,05^2 = 0,00125$ J
\item (Đúng) Động năng của vật nặng tại vị trí có li độ góc $\alpha = 0,05$ rad bằng $0,00375$ J. (Vì): Mệnh đề này đúng vì động năng của quả cầu bằng cơ năng trừ đi thế năng tương ứng: $W_d = W - W_t = 0,005 - 0,00125 = 0,00375$ J
\end{itemchoice}
}
\end{ex}

% ===== Câu 118 =====
% ID: L11C1B5-29-TL
% Mức: VD
\dangbt{Tính động năng, thế năng và cơ năng tại một li độ hoặc thời điểm xác định}
\begin{bt}
% ID: L11C1B5-29-TL
% Mức: VD
Con lắc lò xo có $k=80$ N/m, biên độ $6\,\text{cm}$. Tính cơ năng; động năng và thế năng tại $x=3$ cm.
\end{bt}

% ===== Câu 119 =====
% ID: L11C1B5-29-TLN
% Mức: VD
\dangbt{Tính cơ năng của con lắc đơn}
\begin{ex}
% ID: L11C1B5-29-TLN
% Mức: VD
Một con lắc đơn có vật nặng $0,5\,\text{kg}$, khi qua vị trí cân bằng đạt tốc độ $0,4\,\text{m}$ /s. Cơ năng của con lắc theo đơn vị $J$ bằng bao nhiêu?
\shortans{0,04}
\loigiai{
$W = mvmax²/2 = 0,5·0,4²/2 = 0,04 J.$
}
\end{ex}

% ===== Câu 120 =====
% ID: L11C1B5-29-TN
% Mức: NB
\dangbt{Năng lượng của con lắc đơn dao động điều hòa}
\begin{ex}
% ID: L11C1B5-29-TN
% Mức: NB
Một con lắc đơn dao động với biên độ góc nhỏ. Chọn mốc thế năng ở vị trí cân bằng. Công thức tính thế năng của con lắc ở ly độ góc a là
\choice
{${{W}_{t}}=mgl\left( 1+\cos \alpha \right)$}
{\True $({W)_{t}}=\dfrac{1}{2}mgl.{{\alpha$ }^{2}}}
{$({W)_{t}}=2.mgl.{\cos}^{2}\dfrac{\alpha }{2}$}
{$({W)_{t}}=mgl.\sin\alpha$}
\loigiai{
Thế năng của con lắc đơn ở ly độ góc $\alpha$ (hoặc $a$) được tính bằng công thức tổng quát: ${{W}_{t}}=mgl(1-\cos \alpha)$.
  Khi con lắc dao động với biên độ góc nhỏ, ta có thể sử dụng khai triển Taylor cho $\cos \alpha$: $\cos \alpha \approx 1 - \dfrac{{\alpha }^{2}}{2}$.
  Thay vào công thức thế năng, ta được:
  ${{W}_{t}}=mgl\left( 1-\left( 1-\dfrac{{{\alpha }^{2}}}{2} \right) \right)=mgl\left( \dfrac{{{\alpha }^{2}}}{2} \right)=\dfrac{1}{2}mgl.{{\alpha }^{2}}$.
  Do đó, phương án $A$ là đúng.
}
\end{ex}

% ===== Câu 121 =====
% ID: L11C1B5-30-DS
% Mức: VDC
\begin{ex}
% ID: L11C1B5-30-DS
% Mức: VDC
Một con lắc đơn dao động điều hòa tại nơi có gia tốc trọng trường $g = 9,8$ m/s $^2$. Chiều dài dây treo con lắc là $l = 1$ m, quả nặng có khối lượng $m = 100$ g. Kéo quả nặng lệch khỏi phương thẳng đứng góc $\alpha_0 = 6^\circ$ rồi buông nhẹ. Chọn mốc thế năng tại vị trí cân bằng.
\choiceTF
{\True Thế năng trọng trường cực đại của con lắc đơn được xác định khi quả cầu nặng đạt vị trí biên.}
{\True Biên độ góc của con lắc đơn khi đổi ra đơn vị radian có giá trị xấp xỉ bằng $\dfrac{\pi}{30}$ rad.}
{\True Cơ năng dao động của con lắc đơn có giá trị xấp xỉ bằng $5,37 \cdot 10^{-3}$ J.}
{Tốc độ của quả cầu nặng khi đi qua vị trí cân bằng đạt giá trị xấp xỉ bằng $1$ m/s.}
\loigiai{
\begin{itemchoice}
\item (Đúng) Thế năng trọng trường cực đại của con lắc đơn được xác định khi quả cầu nặng đạt vị trí biên. (Vì): Mệnh đề này đúng vì thế năng con lắc đạt cực đại tại vị trí biên, nơi li độ góc đạt giá trị cực đại $\alpha_0$
\item (Đúng) Biên độ góc của con lắc đơn khi đổi ra đơn vị radian có giá trị xấp xỉ bằng $\dfrac{\pi}{30}$ rad. (Vì): Mệnh đề này đúng vì $\alpha_0 = 6^\circ = 6 \cdot \dfrac{\pi}{180} = \dfrac{\pi}{30}$ rad $\approx 0,1047$ rad
\item (Đúng) Cơ năng dao động của con lắc đơn có giá trị xấp xỉ bằng $5,37 \cdot 10^{-3}$ J. (Vì): Mệnh đề này đúng vì cơ năng của con lắc đơn dao động với góc lệch nhỏ là $W \approx \dfrac{1}{2}mgl\alpha_0^2 = 0,5 \cdot 0,1 \cdot 9,8 \cdot 1 \cdot 0,1047^2 \approx 5,37 \cdot 10^{-3}$ J
\item (Sai) Tốc độ của quả cầu nặng khi đi qua vị trí cân bằng đạt giá trị xấp xỉ bằng $1$ m/s. (Vì): Mệnh đề này sai vì tốc độ cực đại khi qua vị trí cân bằng là $v_{\text{max}} = \sqrt{\dfrac{2W}{m}} = \sqrt{\dfrac{2 \cdot 5,37 \cdot 10^{-3}}{0,1}} \approx 0,328$ m/s
\end{itemchoice}
}
\end{ex}

% ===== Câu 122 =====
% ID: L11C1B5-30-TL
% Mức: VD
\dangbt{Tính động năng, thế năng và cơ năng tại một li độ hoặc thời điểm xác định}
\begin{bt}
% ID: L11C1B5-30-TL
% Mức: VD
Con lắc lò xo có $k=100$ N/m, biên độ $8\,\text{cm}$. Tính cơ năng; động năng và thế năng tại $x=4$ cm.
\end{bt}

% ===== Câu 123 =====
% ID: L11C1B5-30-TLN
% Mức: VD
\dangbt{Tính cơ năng của con lắc đơn}
\begin{ex}
% ID: L11C1B5-30-TLN
% Mức: VD
Một con lắc đơn có vật nặng $0,5\,\text{kg}$, khi qua vị trí cân bằng đạt tốc độ $0,4\,\text{m}$ /s. Cơ năng của con lắc theo đơn vị $J$ bằng bao nhiêu?
\shortans{0,04}
\loigiai{
$W = mvmax²/2 = 0,5·0,4²/2 = 0,04 J.$
}
\end{ex}

% ===== Câu 124 =====
% ID: L11C1B5-30-TN
% Mức: VD
\dangbt{Năng lượng của con lắc đơn dao động điều hòa}
\begin{ex}
% ID: L11C1B5-30-TN
% Mức: VD
Một con lắc đơn có khối lượng $m=5kg$ và độ dài $l=1m$. Góc lệch cực đại của con lắc so với đường thẳng đứng là $\alpha_0=6^\circ \approx 0,1$ rad. Cho $g=10m/s^2$. Tính cơ năng của con lắc:
\choice
{$2,5J$}
{$0,75J$}
{\True $0,25J$}
{$0,5J$}
\loigiai{
Cơ năng của con lắc đơn được tính bằng công thức: $E = mgh_{max}$, trong đó $h_{max}$ là độ cao cực đại mà con lắc đạt được so với vị trí cân bằng.
  Ta có $h_{max} = l(1 - \cos \alpha_0)$.
  Vì $\alpha_0$ nhỏ ($6^\circ \approx 0,1$ rad), ta có thể sử dụng công thức gần đúng cho $\cos \alpha_0 \approx 1 - \dfrac{\alpha_0^2}{2}$.
  Khi đó, $h_{max} \approx l \left(1 - \left(1 - \dfrac{\alpha_0^2}{2}\right)\right) = l \dfrac{\alpha_0^2}{2}$.
  Vậy cơ năng của con lắc là $E = mg \left(l \dfrac{\alpha_0^2}{2}\right) = \dfrac{1}{2} mgl\alpha_0^2$.
  Thay số: $m=5kg$, $g=10m/s^2$, $l=1m$, $\alpha_0=0,1$ rad.
  $E = \dfrac{1}{2} \cdot 5 \cdot 10 \cdot 1 \cdot (0,1)^2 = \dfrac{1}{2} \cdot 50 \cdot 0,01 = 25 \cdot 0,01 = 0,25 J$.
}
\end{ex}

% ===== Câu 125 =====
% ID: L11C1B5-31-DS
% Mức: VDC
\begin{ex}
% ID: L11C1B5-31-DS
% Mức: VDC
Một con lắc đơn dao động điều hòa tại nơi có gia tốc trọng trường $g = 9,8$ m/s $^2$. Chiều dài dây treo con lắc là $l = 1$ m, quả nặng có khối lượng $m = 100$ g. Kéo quả nặng lệch khỏi phương thẳng đứng góc $\alpha_0 = 6^\circ$ rồi buông nhẹ. Chọn mốc thế năng tại vị trí cân bằng.
\choiceTF
{\True Thế năng trọng trường cực đại của con lắc đơn được xác định khi quả cầu nặng đạt vị trí biên.}
{\True Biên độ góc của con lắc đơn khi đổi ra đơn vị radian có giá trị xấp xỉ bằng $\dfrac{\pi}{30}$ rad.}
{\True Cơ năng dao động của con lắc đơn có giá trị xấp xỉ bằng $5,37 \cdot 10^{-3}$ J.}
{Tốc độ của quả cầu nặng khi đi qua vị trí cân bằng đạt giá trị xấp xỉ bằng $1$ m/s.}
\loigiai{
\begin{itemchoice}
\item (Đúng) Thế năng trọng trường cực đại của con lắc đơn được xác định khi quả cầu nặng đạt vị trí biên. (Vì): Mệnh đề này đúng vì thế năng con lắc đạt cực đại tại vị trí biên, nơi li độ góc đạt giá trị cực đại $\alpha_0$
\item (Đúng) Biên độ góc của con lắc đơn khi đổi ra đơn vị radian có giá trị xấp xỉ bằng $\dfrac{\pi}{30}$ rad. (Vì): Mệnh đề này đúng vì $\alpha_0 = 6^\circ = 6 \cdot \dfrac{\pi}{180} = \dfrac{\pi}{30}$ rad $\approx 0,1047$ rad
\item (Đúng) Cơ năng dao động của con lắc đơn có giá trị xấp xỉ bằng $5,37 \cdot 10^{-3}$ J. (Vì): Mệnh đề này đúng vì cơ năng của con lắc đơn dao động với góc lệch nhỏ là $W \approx \dfrac{1}{2}mgl\alpha_0^2 = 0,5 \cdot 0,1 \cdot 9,8 \cdot 1 \cdot 0,1047^2 \approx 5,37 \cdot 10^{-3}$ J
\item (Sai) Tốc độ của quả cầu nặng khi đi qua vị trí cân bằng đạt giá trị xấp xỉ bằng $1$ m/s. (Vì): Mệnh đề này sai vì tốc độ cực đại khi qua vị trí cân bằng là $v_{\text{max}} = \sqrt{\dfrac{2W}{m}} = \sqrt{\dfrac{2 \cdot 5,37 \cdot 10^{-3}}{0,1}} \approx 0,328$ m/s
\end{itemchoice}
}
\end{ex}

% ===== Câu 126 =====
% ID: L11C1B5-31-TL
% Mức: TH
\dangbt{Xác định biểu thức cơ năng trong dao động điều hòa}
\begin{bt}
% ID: L11C1B5-31-TL
% Mức: TH
Con lắc lò xo có $k=50$ N/m, biên độ $4\,\text{cm}$. Tính cơ năng; động năng và thế năng tại $x=2$ cm.
\end{bt}

% ===== Câu 127 =====
% ID: L11C1B5-31-TLN
% Mức: NB
\dangbt{Tính động năng, thế năng và cơ năng tại một li độ hoặc thời điểm xác định}
\begin{ex}
% ID: L11C1B5-31-TLN
% Mức: NB
Con lắc đơn có chiều dài $\ell=1\,\text{m}$, lấy $g=10\,\text{m/s}^2$. Chọn gốc thế năng tại vị trí cân bằng. Con lắc dao động với biên độ góc $\alpha_0=9^\circ$. Giá trị vận tốc của vật tại vị trí mà động năng bằng thế năng là bao nhiêu m/s?
\shortans{0.35}
\loigiai{
Tốc độ cực đại của con lắc:
$v_{\max}=\alpha_0\sqrt{g\ell} =\dfrac{9\pi}{180}\sqrt{10\cdot1} =0,5\,\text{m/s}.$
Tại vị trí động năng bằng thế năng:
$v=\dfrac{v_{\max}}{\sqrt{2}} =\dfrac{0,5}{\sqrt{2}} \approx0,35\,\text{m/s}.$
}
\end{ex}

% ===== Câu 128 =====
% ID: L11C1B5-31-TN
% Mức: NB
\dangbt{Phân tích đồ thị năng lượng}
\begin{ex}
% ID: L11C1B5-31-TN
% Mức: NB
Đồ thị biểu diễn sự biến thiên của thế năng $W_{\text{t}}$ theo li độ $x$ của một vật dao động điều hòa có dạng là một:
\choice
{Đường thẳng đi qua gốc tọa độ}
{Đường sóng hình sin dao động tuần hoàn}
{\True Đường parabol có bề lõm hướng lên phía trên}
{Đường parabol có bề lõm hướng xuống dưới}
\loigiai{
• Công thức liên hệ giữa thế năng và li độ:
 $W_{\text{t}} = \dfrac{1}{2}m\omega^2 x^2.$
 
• Đây là hàm số bậc hai đối với biến $x$, dạng $y = ax^2$ với hệ số $a = \dfrac{1}{2}m\omega^2 \gt  0$.
 
• Do đó đồ thị là một đường parabol có đỉnh trùng gốc tọa độ $O$ và bề lõm hướng lên phía trên.
}
\end{ex}

% ===== Câu 129 =====
% ID: L11C1B5-32-DS
% Mức: VDC
\dangbt{Năng lượng của con lắc đơn dao động điều hòa}
\begin{ex}
% ID: L11C1B5-32-DS
% Mức: VDC
Một con lắc đơn dao động điều hòa tại nơi có gia tốc trọng trường $g = 9,8$ m/s $^2$. Chiều dài dây treo con lắc là $l = 1$ m, quả nặng có khối lượng $m = 100$ g. Kéo quả nặng lệch khỏi phương thẳng đứng góc $\alpha_0 = 6^\circ$ rồi buông nhẹ. Chọn mốc thế năng tại vị trí cân bằng.
\choiceTF
{\True Thế năng trọng trường cực đại của con lắc đơn được xác định khi quả cầu nặng đạt vị trí biên.}
{\True Biên độ góc của con lắc đơn khi đổi ra đơn vị radian có giá trị xấp xỉ bằng $\dfrac{\pi}{30}$ rad.}
{\True Cơ năng dao động của con lắc đơn có giá trị xấp xỉ bằng $5,37 \cdot 10^{-3}$ J.}
{Tốc độ của quả cầu nặng khi đi qua vị trí cân bằng đạt giá trị xấp xỉ bằng $1$ m/s.}
\loigiai{
\begin{itemchoice}
\item (Đúng) Thế năng trọng trường cực đại của con lắc đơn được xác định khi quả cầu nặng đạt vị trí biên. (Vì): Mệnh đề này đúng vì thế năng con lắc đạt cực đại tại vị trí biên, nơi li độ góc đạt giá trị cực đại $\alpha_0$
\item (Đúng) Biên độ góc của con lắc đơn khi đổi ra đơn vị radian có giá trị xấp xỉ bằng $\dfrac{\pi}{30}$ rad. (Vì): Mệnh đề này đúng vì $\alpha_0 = 6^\circ = 6 \cdot \dfrac{\pi}{180} = \dfrac{\pi}{30}$ rad $\approx 0,1047$ rad
\item (Đúng) Cơ năng dao động của con lắc đơn có giá trị xấp xỉ bằng $5,37 \cdot 10^{-3}$ J. (Vì): Mệnh đề này đúng vì cơ năng của con lắc đơn dao động với góc lệch nhỏ là $W \approx \dfrac{1}{2}mgl\alpha_0^2 = 0,5 \cdot 0,1 \cdot 9,8 \cdot 1 \cdot 0,1047^2 \approx 5,37 \cdot 10^{-3}$ J
\item (Sai) Tốc độ của quả cầu nặng khi đi qua vị trí cân bằng đạt giá trị xấp xỉ bằng $1$ m/s. (Vì): Mệnh đề này sai vì tốc độ cực đại khi qua vị trí cân bằng là $v_{\text{max}} = \sqrt{\dfrac{2W}{m}} = \sqrt{\dfrac{2 \cdot 5,37 \cdot 10^{-3}}{0,1}} \approx 0,328$ m/s
\end{itemchoice}
}
\end{ex}

% ===== Câu 130 =====
% ID: L11C1B5-32-TL
% Mức: VD
\dangbt{Xác định biểu thức cơ năng trong dao động điều hòa}
\begin{bt}
% ID: L11C1B5-32-TL
% Mức: VD
Con lắc lò xo có $k=80$ N/m, biên độ $6\,\text{cm}$. Tính cơ năng; động năng và thế năng tại $x=3$ cm.
\end{bt}

% ===== Câu 131 =====
% ID: L11C1B5-32-TLN
% Mức: NB
\dangbt{Tính động năng, thế năng và cơ năng tại một li độ hoặc thời điểm xác định}
\begin{ex}
% ID: L11C1B5-32-TLN
% Mức: NB
Con lắc đơn có chiều dài $\ell=1\,\text{m}$, lấy $g=10\,\text{m/s}^2$. Chọn gốc thế năng tại vị trí cân bằng. Con lắc dao động với biên độ góc $\alpha_0=9^\circ$. Giá trị vận tốc của vật tại vị trí mà động năng bằng thế năng là bao nhiêu m/s?
\shortans{0.35}
\loigiai{
Tốc độ cực đại của con lắc:
$v_{\max}=\alpha_0\sqrt{g\ell} =\dfrac{9\pi}{180}\sqrt{10\cdot1} =0,5\,\text{m/s}.$
Tại vị trí động năng bằng thế năng:
$v=\dfrac{v_{\max}}{\sqrt{2}} =\dfrac{0,5}{\sqrt{2}} \approx0,35\,\text{m/s}.$
}
\end{ex}

% ===== Câu 132 =====
% ID: L11C1B5-32-TN
% Mức: NB
\dangbt{Phân tích đồ thị năng lượng}
\begin{ex}
% ID: L11C1B5-32-TN
% Mức: NB
Đồ thị biểu diễn sự biến thiên của thế năng $W_{\text{t}}$ theo li độ $x$ của một vật dao động điều hòa có dạng là một:
\choice
{Đường thẳng đi qua gốc tọa độ}
{Đường sóng hình sin dao động tuần hoàn}
{\True Đường parabol có bề lõm hướng lên phía trên}
{Đường parabol có bề lõm hướng xuống dưới}
\loigiai{
• Công thức liên hệ giữa thế năng và li độ:
 $W_{\text{t}} = \dfrac{1}{2}m\omega^2 x^2.$
 
• Đây là hàm số bậc hai đối với biến $x$, dạng $y = ax^2$ với hệ số $a = \dfrac{1}{2}m\omega^2 \gt  0$.
 
• Do đó đồ thị là một đường parabol có đỉnh trùng gốc tọa độ $O$ và bề lõm hướng lên phía trên.
}
\end{ex}

% ===== Câu 133 =====
% ID: L11C1B5-33-DS
% Mức: TH
\begin{ex}
% ID: L11C1B5-33-DS
% Mức: TH
Một vật dao động điều hòa dọc theo trục Ox quanh vị trí cân bằng O. Khi biểu diễn đồ thị các loại năng lượng của vật theo li độ $x$ trên hệ trục tọa độ, ta thu được các đường đồ thị tương ứng:
\choiceTF
{\True Đồ thị biểu diễn sự phụ thuộc của thế năng vào li độ $x$ là một đường parabol có bề lõm hướng lên.}
{\True Đồ thị biểu diễn sự phụ thuộc của động năng vào li độ $x$ là một đường parabol có bề lõm hướng xuống.}
{\True Đồ thị biểu diễn sự phụ thuộc của cơ năng toàn phần vào li độ $x$ là một đường thẳng song song với trục hoành Ox.}
{Đỉnh parabol của đồ thị động năng theo li độ $W_d(x)$ nằm tại các vị trí biên $x = \pm A$.}
\loigiai{
\begin{itemchoice}
\item (Đúng) Đồ thị biểu diễn sự phụ thuộc của thế năng vào li độ $x$ là một đường parabol có bề lõm hướng lên. (Vì): Mệnh đề này đúng vì thế năng $W_t = \dfrac{1}{2}m\omega^2 x^2$ là hàm bậc hai có hệ số $a \gt  0$, đồ thị là parabol hướng lên
\item (Đúng) Đồ thị biểu diễn sự phụ thuộc của động năng vào li độ $x$ là một đường parabol có bề lõm hướng xuống. (Vì): Mệnh đề này đúng vì động năng $W_d = W - W_t = \dfrac{1}{2}m\omega^2 A^2 - \dfrac{1}{2}m\omega^2 x^2$ là hàm bậc hai có hệ số của $x^2$ là âm, đồ thị là parabol hướng xuống
\item (Đúng) Đồ thị biểu diễn sự phụ thuộc của cơ năng toàn phần vào li độ $x$ là một đường thẳng song song với trục hoành Ox. (Vì): Mệnh đề này đúng vì cơ năng $W = \dfrac{1}{2}m\omega^2 A^2$ là một hằng số không phụ thuộc vào tọa độ $x$, đồ thị là đường thẳng nằm ngang
\item (Sai) Đỉnh parabol của đồ thị động năng theo li độ $W_d(x)$ nằm tại các vị trí biên $x = \pm A$.
\end{itemchoice}
}
\end{ex}

% ===== Câu 134 =====
% ID: L11C1B5-33-TL
% Mức: VD
\dangbt{Xác định biểu thức cơ năng trong dao động điều hòa}
\begin{bt}
% ID: L11C1B5-33-TL
% Mức: VD
Con lắc lò xo có $k=100$ N/m, biên độ $8\,\text{cm}$. Tính cơ năng; động năng và thế năng tại $x=4$ cm.
\end{bt}

% ===== Câu 135 =====
% ID: L11C1B5-33-TLN
% Mức: NB
\dangbt{Tính động năng, thế năng và cơ năng tại một li độ hoặc thời điểm xác định}
\begin{ex}
% ID: L11C1B5-33-TLN
% Mức: NB
Con lắc đơn có chiều dài $\ell=1\,\text{m}$, lấy $g=10\,\text{m/s}^2$. Chọn gốc thế năng tại vị trí cân bằng. Con lắc dao động với biên độ góc $\alpha_0=9^\circ$. Giá trị vận tốc của vật tại vị trí mà động năng bằng thế năng là bao nhiêu m/s?
\shortans{0.35}
\loigiai{
Tốc độ cực đại của con lắc:
$v_{\max}=\alpha_0\sqrt{g\ell} =\dfrac{9\pi}{180}\sqrt{10\cdot1} =0,5\,\text{m/s}.$
Tại vị trí động năng bằng thế năng:
$v=\dfrac{v_{\max}}{\sqrt{2}} =\dfrac{0,5}{\sqrt{2}} \approx0,35\,\text{m/s}.$
}
\end{ex}

% ===== Câu 136 =====
% ID: L11C1B5-33-TN
% Mức: NB
\dangbt{Phân tích đồ thị năng lượng}
\begin{ex}
% ID: L11C1B5-33-TN
% Mức: NB
Đồ thị biểu diễn sự biến thiên của thế năng $W_{\text{t}}$ theo li độ $x$ của một vật dao động điều hòa có dạng là một:
\choice
{Đường thẳng đi qua gốc tọa độ}
{Đường sóng hình sin dao động tuần hoàn}
{\True Đường parabol có bề lõm hướng lên phía trên}
{Đường parabol có bề lõm hướng xuống dưới}
\loigiai{
• Công thức liên hệ giữa thế năng và li độ:
 $W_{\text{t}} = \dfrac{1}{2}m\omega^2 x^2.$
 
• Đây là hàm số bậc hai đối với biến $x$, dạng $y = ax^2$ với hệ số $a = \dfrac{1}{2}m\omega^2 \gt  0$.
 
• Do đó đồ thị là một đường parabol có đỉnh trùng gốc tọa độ $O$ và bề lõm hướng lên phía trên.
}
\end{ex}

% ===== Câu 137 =====
% ID: L11C1B5-34-DS
% Mức: TH
\begin{ex}
% ID: L11C1B5-34-DS
% Mức: TH
Một vật dao động điều hòa dọc theo trục Ox quanh vị trí cân bằng O. Khi biểu diễn đồ thị các loại năng lượng của vật theo li độ $x$ trên hệ trục tọa độ, ta thu được các đường đồ thị tương ứng:
\choiceTF
{\True Đồ thị biểu diễn sự phụ thuộc của thế năng vào li độ $x$ là một đường parabol có bề lõm hướng lên.}
{\True Đồ thị biểu diễn sự phụ thuộc của động năng vào li độ $x$ là một đường parabol có bề lõm hướng xuống.}
{\True Đồ thị biểu diễn sự phụ thuộc của cơ năng toàn phần vào li độ $x$ là một đường thẳng song song với trục hoành Ox.}
{Đỉnh parabol của đồ thị động năng theo li độ $W_d(x)$ nằm tại các vị trí biên $x = \pm A$.}
\loigiai{
\begin{itemchoice}
\item (Đúng) Đồ thị biểu diễn sự phụ thuộc của thế năng vào li độ $x$ là một đường parabol có bề lõm hướng lên. (Vì): Mệnh đề này đúng vì thế năng $W_t = \dfrac{1}{2}m\omega^2 x^2$ là hàm bậc hai có hệ số $a \gt  0$, đồ thị là parabol hướng lên
\item (Đúng) Đồ thị biểu diễn sự phụ thuộc của động năng vào li độ $x$ là một đường parabol có bề lõm hướng xuống. (Vì): Mệnh đề này đúng vì động năng $W_d = W - W_t = \dfrac{1}{2}m\omega^2 A^2 - \dfrac{1}{2}m\omega^2 x^2$ là hàm bậc hai có hệ số của $x^2$ là âm, đồ thị là parabol hướng xuống
\item (Đúng) Đồ thị biểu diễn sự phụ thuộc của cơ năng toàn phần vào li độ $x$ là một đường thẳng song song với trục hoành Ox. (Vì): Mệnh đề này đúng vì cơ năng $W = \dfrac{1}{2}m\omega^2 A^2$ là một hằng số không phụ thuộc vào tọa độ $x$, đồ thị là đường thẳng nằm ngang
\item (Sai) Đỉnh parabol của đồ thị động năng theo li độ $W_d(x)$ nằm tại các vị trí biên $x = \pm A$.
\end{itemchoice}
}
\end{ex}

% ===== Câu 138 =====
% ID: L11C1B5-34-TL
% Mức: TH
\dangbt{Xác định và tính toán các đại lượng vận tốc, gia tốc tại thời điểm xác định}
\begin{ex}
% ID: L11C1B5-34-TL
% Mức: TH
Từ $x=6\cos(5t+\pi/3)$ cm, hãy lập biểu thức vận tốc và gia tốc.
\choiceTF

\end{ex}

% ===== Câu 139 =====
% ID: L11C1B5-34-TLN
% Mức: NB
\dangbt{Tính động năng, thế năng và cơ năng tại một li độ hoặc thời điểm xác định}
\begin{ex}
% ID: L11C1B5-34-TLN
% Mức: NB
Con lắc lò xo mà vật dao động có khối lượng $1\,\text{kg}$, dao động điều hòa với cơ năng $125\,\text{mJ}$. Tại thời điểm ban đầu, vật có vận tốc $25\,\text{cm/s}$ và gia tốc $-6,25\sqrt{3}\,\text{m/s}^2$. Biên độ dao động của vật là bao nhiêu cm?
\shortans{2}
\loigiai{
Ta có:
$W=\dfrac{(ma)^2}{2k}+\dfrac{mv^2}{2}.$
Thay số:
$125\cdot10^{-3} =\dfrac{\left(-6,25\sqrt{3}\right)^2}{2k} +\dfrac{1\cdot(0,25)^2}{2} \Rightarrow k=625\,\text{N/m}.$
Biên độ dao động:
$A=\sqrt{\dfrac{2W}{k}} =\sqrt{\dfrac{2\cdot125\cdot10^{-3}}{625}} =0,02\,\text{m}=2\,\text{cm}.$
}
\end{ex}

% ===== Câu 140 =====
% ID: L11C1B5-34-TN
% Mức: NB
\dangbt{Phân tích đồ thị năng lượng}
\begin{ex}
% ID: L11C1B5-34-TN
% Mức: NB
Khi phân tích đồ thị biểu diễn động năng $W_{\text{đ}}$ và thế năng $W_{\text{t}}$ theo li độ $x$ của một vật dao động điều hòa, tọa độ giao điểm của hai đồ thị này cho biết:
\choice
{Vị trí vật đạt tốc độ cực đại}
{Vị trí vật đạt gia tốc cực đại}
{\True Vị trí tại đó động năng bằng thế năng}
{Vị trí cân bằng của vật}
\loigiai{
• Giao điểm của hai đường đồ thị ứng với điểm tại đó hai giá trị hàm số bằng nhau: $W_{\text{đ}} = W_{\text{t}}$.
 
• Tại vị trí này, năng lượng của hệ được chia đều cho hai thành phần động năng và thế năng.
 
• Tọa độ giao điểm trên trục hoành chính là các li độ $x = \pm \dfrac{A}{\sqrt{2}}$.
}
\end{ex}

% ===== Câu 141 =====
% ID: L11C1B5-35-DS
% Mức: TH
\begin{ex}
% ID: L11C1B5-35-DS
% Mức: TH
Một vật dao động điều hòa dọc theo trục Ox quanh vị trí cân bằng O. Khi biểu diễn đồ thị các loại năng lượng của vật theo li độ $x$ trên hệ trục tọa độ, ta thu được các đường đồ thị tương ứng:
\choiceTF
{\True Đồ thị biểu diễn sự phụ thuộc của thế năng vào li độ $x$ là một đường parabol có bề lõm hướng lên.}
{\True Đồ thị biểu diễn sự phụ thuộc của động năng vào li độ $x$ là một đường parabol có bề lõm hướng xuống.}
{\True Đồ thị biểu diễn sự phụ thuộc của cơ năng toàn phần vào li độ $x$ là một đường thẳng song song với trục hoành Ox.}
{Đỉnh parabol của đồ thị động năng theo li độ $W_d(x)$ nằm tại các vị trí biên $x = \pm A$.}
\loigiai{
\begin{itemchoice}
\item (Đúng) Đồ thị biểu diễn sự phụ thuộc của thế năng vào li độ $x$ là một đường parabol có bề lõm hướng lên. (Vì): Mệnh đề này đúng vì thế năng $W_t = \dfrac{1}{2}m\omega^2 x^2$ là hàm bậc hai có hệ số $a \gt  0$, đồ thị là parabol hướng lên
\item (Đúng) Đồ thị biểu diễn sự phụ thuộc của động năng vào li độ $x$ là một đường parabol có bề lõm hướng xuống. (Vì): Mệnh đề này đúng vì động năng $W_d = W - W_t = \dfrac{1}{2}m\omega^2 A^2 - \dfrac{1}{2}m\omega^2 x^2$ là hàm bậc hai có hệ số của $x^2$ là âm, đồ thị là parabol hướng xuống
\item (Đúng) Đồ thị biểu diễn sự phụ thuộc của cơ năng toàn phần vào li độ $x$ là một đường thẳng song song với trục hoành Ox. (Vì): Mệnh đề này đúng vì cơ năng $W = \dfrac{1}{2}m\omega^2 A^2$ là một hằng số không phụ thuộc vào tọa độ $x$, đồ thị là đường thẳng nằm ngang
\item (Sai) Đỉnh parabol của đồ thị động năng theo li độ $W_d(x)$ nằm tại các vị trí biên $x = \pm A$.
\end{itemchoice}
}
\end{ex}

% ===== Câu 142 =====
% ID: L11C1B5-35-TL
% Mức: VD
\dangbt{Xác định và tính toán các đại lượng vận tốc, gia tốc tại thời điểm xác định}
\begin{ex}
% ID: L11C1B5-35-TL
% Mức: VD
Từ $x=7\cos(2t+\pi/3)$ cm, hãy lập biểu thức vận tốc và gia tốc.
\choiceTF

\end{ex}

% ===== Câu 143 =====
% ID: L11C1B5-35-TLN
% Mức: NB
\dangbt{Tính động năng, thế năng và cơ năng tại một li độ hoặc thời điểm xác định}
\begin{ex}
% ID: L11C1B5-35-TLN
% Mức: NB
Con lắc lò xo mà vật dao động có khối lượng $1\,\text{kg}$, dao động điều hòa với cơ năng $125\,\text{mJ}$. Tại thời điểm ban đầu, vật có vận tốc $25\,\text{cm/s}$ và gia tốc $-6,25\sqrt{3}\,\text{m/s}^2$. Biên độ dao động của vật là bao nhiêu cm?
\shortans{2}
\loigiai{
Ta có:
$W=\dfrac{(ma)^2}{2k}+\dfrac{mv^2}{2}.$
Thay số:
$125\cdot10^{-3} =\dfrac{\left(-6,25\sqrt{3}\right)^2}{2k} +\dfrac{1\cdot(0,25)^2}{2} \Rightarrow k=625\,\text{N/m}.$
Biên độ dao động:
$A=\sqrt{\dfrac{2W}{k}} =\sqrt{\dfrac{2\cdot125\cdot10^{-3}}{625}} =0,02\,\text{m}=2\,\text{cm}.$
}
\end{ex}

% ===== Câu 144 =====
% ID: L11C1B5-35-TN
% Mức: NB
\dangbt{Phân tích đồ thị năng lượng}
\begin{ex}
% ID: L11C1B5-35-TN
% Mức: NB
Khi phân tích đồ thị biểu diễn động năng $W_{\text{đ}}$ và thế năng $W_{\text{t}}$ theo li độ $x$ của một vật dao động điều hòa, tọa độ giao điểm của hai đồ thị này cho biết:
\choice
{Vị trí vật đạt tốc độ cực đại}
{Vị trí vật đạt gia tốc cực đại}
{\True Vị trí tại đó động năng bằng thế năng}
{Vị trí cân bằng của vật}
\loigiai{
• Giao điểm của hai đường đồ thị ứng với điểm tại đó hai giá trị hàm số bằng nhau: $W_{\text{đ}} = W_{\text{t}}$.
 
• Tại vị trí này, năng lượng của hệ được chia đều cho hai thành phần động năng và thế năng.
 
• Tọa độ giao điểm trên trục hoành chính là các li độ $x = \pm \dfrac{A}{\sqrt{2}}$.
}
\end{ex}

% ===== Câu 145 =====
% ID: L11C1B5-36-DS
% Mức: TH
\begin{ex}
% ID: L11C1B5-36-DS
% Mức: TH
Đồ thị động năng theo li độ $W_d(x)$ của một vật dao động điều hòa là một đường parabol có bề lõm hướng xuống dưới. Phân tích các thông số động học từ đồ thị:
\choiceTF
{\True Điểm cao nhất của parabol (đỉnh parabol) nằm trên trục tung tại vị trí có tung độ bằng cơ năng $W$ của hệ.}
{\True Đồ thị động năng cắt trục hoành nằm ngang tại hai điểm có tọa độ li độ tương ứng $x = -A$ và $x = A$.}
{Khi li độ của vật dao động tăng dần từ vị trí biên âm $-A$ đến vị trí cân bằng $0$, động năng của vật giảm dần.}
{\True Động năng của vật luôn không âm nên toàn bộ đường đồ thị động năng luôn nằm phía trên trục hoành Ox.}
\loigiai{
\begin{itemchoice}
\item (Đúng) Điểm cao nhất của parabol (đỉnh parabol) nằm trên trục tung tại vị trí có tung độ bằng cơ năng $W$ của hệ. (Vì): Mệnh đề này đúng vì động năng cực đại tại vị trí cân bằng $x = 0$ có giá trị bằng cơ năng, tương ứng đỉnh của parabol nằm tại trục tung Oy
\item (Đúng) Đồ thị động năng cắt trục hoành nằm ngang tại hai điểm có tọa độ li độ tương ứng $x = -A$ và $x = A$. (Vì): Mệnh đề này đúng vì tại các vị trí biên $x = \pm A$, động năng bằng 0 nên đồ thị cắt trục hoành tại các giao điểm này
\item (Sai) Khi li độ của vật dao động tăng dần từ vị trí biên âm $-A$ đến vị trí cân bằng $0$, động năng của vật giảm dần. (Vì): Mệnh đề này sai vì khi đi từ biên âm về vị trí cân bằng, tốc độ của vật tăng dần từ không đến cực đại, nên động năng của vật tăng dần
\item (Đúng) Động năng của vật luôn không âm nên toàn bộ đường đồ thị động năng luôn nằm phía trên trục hoành Ox. (Vì): Mệnh đề này đúng vì động năng $W_d = \dfrac{1}{2}mv^2 \ge 0$, do đó đồ thị động năng không bao giờ đi xuống vùng tung độ âm
\end{itemchoice}
}
\end{ex}

% ===== Câu 146 =====
% ID: L11C1B5-36-TL
% Mức: VD
\dangbt{Xác định và tính toán các đại lượng vận tốc, gia tốc tại thời điểm xác định}
\begin{ex}
% ID: L11C1B5-36-TL
% Mức: VD
Từ $x=8\cos(3t+\pi/3)$ cm, hãy lập biểu thức vận tốc và gia tốc.
\choiceTF

\end{ex}

% ===== Câu 147 =====
% ID: L11C1B5-36-TLN
% Mức: NB
\dangbt{Tính động năng, thế năng và cơ năng tại một li độ hoặc thời điểm xác định}
\begin{ex}
% ID: L11C1B5-36-TLN
% Mức: NB
Con lắc lò xo mà vật dao động có khối lượng $1\,\text{kg}$, dao động điều hòa với cơ năng $125\,\text{mJ}$. Tại thời điểm ban đầu, vật có vận tốc $25\,\text{cm/s}$ và gia tốc $-6,25\sqrt{3}\,\text{m/s}^2$. Biên độ dao động của vật là bao nhiêu cm?
\shortans{2}
\loigiai{
Ta có:
$W=\dfrac{(ma)^2}{2k}+\dfrac{mv^2}{2}.$
Thay số:
$125\cdot10^{-3} =\dfrac{\left(-6,25\sqrt{3}\right)^2}{2k} +\dfrac{1\cdot(0,25)^2}{2} \Rightarrow k=625\,\text{N/m}.$
Biên độ dao động:
$A=\sqrt{\dfrac{2W}{k}} =\sqrt{\dfrac{2\cdot125\cdot10^{-3}}{625}} =0,02\,\text{m}=2\,\text{cm}.$
}
\end{ex}

% ===== Câu 148 =====
% ID: L11C1B5-36-TN
% Mức: NB
\dangbt{Phân tích đồ thị năng lượng}
\begin{ex}
% ID: L11C1B5-36-TN
% Mức: NB
Khi phân tích đồ thị biểu diễn động năng $W_{\text{đ}}$ và thế năng $W_{\text{t}}$ theo li độ $x$ của một vật dao động điều hòa, tọa độ giao điểm của hai đồ thị này cho biết:
\choice
{Vị trí vật đạt tốc độ cực đại}
{Vị trí vật đạt gia tốc cực đại}
{\True Vị trí tại đó động năng bằng thế năng}
{Vị trí cân bằng của vật}
\loigiai{
• Giao điểm của hai đường đồ thị ứng với điểm tại đó hai giá trị hàm số bằng nhau: $W_{\text{đ}} = W_{\text{t}}$.
 
• Tại vị trí này, năng lượng của hệ được chia đều cho hai thành phần động năng và thế năng.
 
• Tọa độ giao điểm trên trục hoành chính là các li độ $x = \pm \dfrac{A}{\sqrt{2}}$.
}
\end{ex}

% ===== Câu 149 =====
% ID: L11C1B5-37-DS
% Mức: TH
\begin{ex}
% ID: L11C1B5-37-DS
% Mức: TH
Đồ thị động năng theo li độ $W_d(x)$ của một vật dao động điều hòa là một đường parabol có bề lõm hướng xuống dưới. Phân tích các thông số động học từ đồ thị:
\choiceTF
{\True Điểm cao nhất của parabol (đỉnh parabol) nằm trên trục tung tại vị trí có tung độ bằng cơ năng $W$ của hệ.}
{\True Đồ thị động năng cắt trục hoành nằm ngang tại hai điểm có tọa độ li độ tương ứng $x = -A$ và $x = A$.}
{Khi li độ của vật dao động tăng dần từ vị trí biên âm $-A$ đến vị trí cân bằng $0$, động năng của vật giảm dần.}
{\True Động năng của vật luôn không âm nên toàn bộ đường đồ thị động năng luôn nằm phía trên trục hoành Ox.}
\loigiai{
\begin{itemchoice}
\item (Đúng) Điểm cao nhất của parabol (đỉnh parabol) nằm trên trục tung tại vị trí có tung độ bằng cơ năng $W$ của hệ. (Vì): Mệnh đề này đúng vì động năng cực đại tại vị trí cân bằng $x = 0$ có giá trị bằng cơ năng, tương ứng đỉnh của parabol nằm tại trục tung Oy
\item (Đúng) Đồ thị động năng cắt trục hoành nằm ngang tại hai điểm có tọa độ li độ tương ứng $x = -A$ và $x = A$. (Vì): Mệnh đề này đúng vì tại các vị trí biên $x = \pm A$, động năng bằng 0 nên đồ thị cắt trục hoành tại các giao điểm này
\item (Sai) Khi li độ của vật dao động tăng dần từ vị trí biên âm $-A$ đến vị trí cân bằng $0$, động năng của vật giảm dần. (Vì): Mệnh đề này sai vì khi đi từ biên âm về vị trí cân bằng, tốc độ của vật tăng dần từ không đến cực đại, nên động năng của vật tăng dần
\item (Đúng) Động năng của vật luôn không âm nên toàn bộ đường đồ thị động năng luôn nằm phía trên trục hoành Ox. (Vì): Mệnh đề này đúng vì động năng $W_d = \dfrac{1}{2}mv^2 \ge 0$, do đó đồ thị động năng không bao giờ đi xuống vùng tung độ âm
\end{itemchoice}
}
\end{ex}

% ===== Câu 150 =====
% ID: L11C1B5-37-TL
% Mức: TH
\dangbt{Xác định vị trí dựa vào tỉ số động năng và thế năng}
\begin{bt}
% ID: L11C1B5-37-TL
% Mức: TH
Con lắc lò xo có $k=50$ N/m, biên độ $4\,\text{cm}$. Tính cơ năng; động năng và thế năng tại $x=2$ cm.
\end{bt}

% ===== Câu 151 =====
% ID: L11C1B5-37-TLN
% Mức: NB
\dangbt{Tính động năng, thế năng và cơ năng tại một li độ hoặc thời điểm xác định}
\begin{ex}
% ID: L11C1B5-37-TLN
% Mức: NB
Một con lắc đơn có khối lượng của vật nặng là $200\,\text{g}$ dao động với phương trình $s=10\sin(2\pi t)\,\text{cm}$. Ở thời điểm $t=\dfrac{\pi}{6}\,\text{s}$, con lắc có động năng là bao nhiêu mJ?
\shortans{1}
\loigiai{
Vận tốc của vật:
$v=s'=0,1\cdot2\pi\cos(2\pi t).$
Theo kết quả tính toán:
$v=-0,1\,\text{m/s}.$
Động năng của con lắc:
$W_{\text{đ}}=\dfrac{1}{2}mv^2 =\dfrac{1}{2}\cdot0,2\cdot(-0,1)^2 =0,001\,\text{J}=1\,\text{mJ}.$
}
\end{ex}

% ===== Câu 152 =====
% ID: L11C1B5-37-TN
% Mức: NB
\dangbt{Phân tích đồ thị năng lượng}
\begin{ex}
% ID: L11C1B5-37-TN
% Mức: NB
Đồ thị động năng theo li độ $W_{\text{đ}}(x)$ của một vật dao động điều hòa là một đường parabol có bề lõm hướng xuống. Điểm cao nhất của đồ thị này (đỉnh parabol) ứng với vị trí:
\choice
{Vị trí biên dương $x = A$}
{Vị trí biên âm $x = -A$}
{\True Vị trí cân bằng $x = 0$}
{Vị trí có li độ $x = \pm \dfrac{A}{2}$}
\loigiai{
• Công thức xác định động năng theo li độ:
 $W_{\text{đ}} = W - W_{\text{t}} = \dfrac{1}{2}m\omega^2(A^2 - x^2).$
 
• Đỉnh của đường parabol ứng với giá trị cực đại của động năng $W_{\text{đ}}$.
 
• Động năng cực đại khi $x^2 = 0 \Rightarrow x = 0$, tức là khi vật đi qua vị trí cân bằng.
}
\end{ex}

% ===== Câu 153 =====
% ID: L11C1B5-38-DS
% Mức: TH
\begin{ex}
% ID: L11C1B5-38-DS
% Mức: TH
Đồ thị động năng theo li độ $W_d(x)$ của một vật dao động điều hòa là một đường parabol có bề lõm hướng xuống dưới. Phân tích các thông số động học từ đồ thị:
\choiceTF
{\True Điểm cao nhất của parabol (đỉnh parabol) nằm trên trục tung tại vị trí có tung độ bằng cơ năng $W$ của hệ.}
{\True Đồ thị động năng cắt trục hoành nằm ngang tại hai điểm có tọa độ li độ tương ứng $x = -A$ và $x = A$.}
{Khi li độ của vật dao động tăng dần từ vị trí biên âm $-A$ đến vị trí cân bằng $0$, động năng của vật giảm dần.}
{\True Động năng của vật luôn không âm nên toàn bộ đường đồ thị động năng luôn nằm phía trên trục hoành Ox.}
\loigiai{
\begin{itemchoice}
\item (Đúng) Điểm cao nhất của parabol (đỉnh parabol) nằm trên trục tung tại vị trí có tung độ bằng cơ năng $W$ của hệ. (Vì): Mệnh đề này đúng vì động năng cực đại tại vị trí cân bằng $x = 0$ có giá trị bằng cơ năng, tương ứng đỉnh của parabol nằm tại trục tung Oy
\item (Đúng) Đồ thị động năng cắt trục hoành nằm ngang tại hai điểm có tọa độ li độ tương ứng $x = -A$ và $x = A$. (Vì): Mệnh đề này đúng vì tại các vị trí biên $x = \pm A$, động năng bằng 0 nên đồ thị cắt trục hoành tại các giao điểm này
\item (Sai) Khi li độ của vật dao động tăng dần từ vị trí biên âm $-A$ đến vị trí cân bằng $0$, động năng của vật giảm dần. (Vì): Mệnh đề này sai vì khi đi từ biên âm về vị trí cân bằng, tốc độ của vật tăng dần từ không đến cực đại, nên động năng của vật tăng dần
\item (Đúng) Động năng của vật luôn không âm nên toàn bộ đường đồ thị động năng luôn nằm phía trên trục hoành Ox. (Vì): Mệnh đề này đúng vì động năng $W_d = \dfrac{1}{2}mv^2 \ge 0$, do đó đồ thị động năng không bao giờ đi xuống vùng tung độ âm
\end{itemchoice}
}
\end{ex}

% ===== Câu 154 =====
% ID: L11C1B5-38-TL
% Mức: VD
\dangbt{Xác định vị trí dựa vào tỉ số động năng và thế năng}
\begin{bt}
% ID: L11C1B5-38-TL
% Mức: VD
Con lắc lò xo có $k=80$ N/m, biên độ $6\,\text{cm}$. Tính cơ năng; động năng và thế năng tại $x=3$ cm.
\end{bt}

% ===== Câu 155 =====
% ID: L11C1B5-38-TLN
% Mức: NB
\dangbt{Tính động năng, thế năng và cơ năng tại một li độ hoặc thời điểm xác định}
\begin{ex}
% ID: L11C1B5-38-TLN
% Mức: NB
Một con lắc đơn có khối lượng của vật nặng là $200\,\text{g}$ dao động với phương trình $s=10\sin(2\pi t)\,\text{cm}$. Ở thời điểm $t=\dfrac{\pi}{6}\,\text{s}$, con lắc có động năng là bao nhiêu mJ?
\shortans{1}
\loigiai{
Vận tốc của vật:
$v=s'=0,1\cdot2\pi\cos(2\pi t).$
Theo kết quả tính toán:
$v=-0,1\,\text{m/s}.$
Động năng của con lắc:
$W_{\text{đ}}=\dfrac{1}{2}mv^2 =\dfrac{1}{2}\cdot0,2\cdot(-0,1)^2 =0,001\,\text{J}=1\,\text{mJ}.$
}
\end{ex}

% ===== Câu 156 =====
% ID: L11C1B5-38-TN
% Mức: NB
\dangbt{Phân tích đồ thị năng lượng}
\begin{ex}
% ID: L11C1B5-38-TN
% Mức: NB
Đồ thị động năng theo li độ $W_{\text{đ}}(x)$ của một vật dao động điều hòa là một đường parabol có bề lõm hướng xuống. Điểm cao nhất của đồ thị này (đỉnh parabol) ứng với vị trí:
\choice
{Vị trí biên dương $x = A$}
{Vị trí biên âm $x = -A$}
{\True Vị trí cân bằng $x = 0$}
{Vị trí có li độ $x = \pm \dfrac{A}{2}$}
\loigiai{
• Công thức xác định động năng theo li độ:
 $W_{\text{đ}} = W - W_{\text{t}} = \dfrac{1}{2}m\omega^2(A^2 - x^2).$
 
• Đỉnh của đường parabol ứng với giá trị cực đại của động năng $W_{\text{đ}}$.
 
• Động năng cực đại khi $x^2 = 0 \Rightarrow x = 0$, tức là khi vật đi qua vị trí cân bằng.
}
\end{ex}

% ===== Câu 157 =====
% ID: L11C1B5-39-DS
% Mức: VD
\begin{ex}
% ID: L11C1B5-39-DS
% Mức: VD
Cho đồ thị biểu diễn sự phụ thuộc của thế năng $W_t$ và động năng $W_d$ vào li độ $x$ của một con lắc lò xo. Trên đồ thị, điểm giao nhau của hai đường có tọa độ li độ là $x = \pm 4$ cm và giá trị năng lượng là $0,16$ J.
\choiceTF
{\True Tại điểm giao nhau của hai đồ thị, thế năng bằng động năng và cùng bằng $0,16$ J.}
{\True Cơ năng dao động của con lắc lò xo được xác định bằng $0,32$ J.}
{Biên độ dao động của con lắc là $8$ cm.}
{\True Độ cứng của lò xo con lắc dùng trong khảo sát bằng $200$ N/m.}
\loigiai{
\begin{itemchoice}
\item (Đúng) Tại điểm giao nhau của hai đồ thị, thế năng bằng động năng và cùng bằng $0,16$ J. (Vì): Mệnh đề này đúng vì giao điểm của hai đồ thị là vị trí thế năng bằng động năng, tung độ giao điểm cho biết giá trị năng lượng thành phần bằng $0,16$ J
\item (Đúng) Cơ năng dao động của con lắc lò xo được xác định bằng $0,32$ J. (Vì): Mệnh đề này đúng vì tại vị trí thế năng bằng động năng, cơ năng của vật được tính là: $W = W_d + W_t = 2W_t = 2 \cdot 0,16 = 0,32$ J
\item (Sai) Biên độ dao động của con lắc là $8$ cm. (Vì): Mệnh đề này sai vì tại vị trí $W_d = W_t \Rightarrow x = \pm \dfrac{A}{\sqrt{2}} \Rightarrow 4 = \dfrac{A}{\sqrt{2}} \Rightarrow A = 4\sqrt{2}$ cm $\approx 5,66$ cm
\item (Đúng) Độ cứng của lò xo con lắc dùng trong khảo sát bằng $200$ N/m. (Vì): Mệnh đề này đúng vì đổi $x = 0,04$ m; áp dụng thế năng $W_t = \dfrac{1}{2}kx^2 \Rightarrow 0,16 = 0,5 \cdot k \cdot 0,04^2 \Rightarrow k = 200$ N/m
\end{itemchoice}
}
\end{ex}

% ===== Câu 158 =====
% ID: L11C1B5-39-TL
% Mức: VD
\dangbt{Xác định vị trí dựa vào tỉ số động năng và thế năng}
\begin{bt}
% ID: L11C1B5-39-TL
% Mức: VD
Con lắc lò xo có $k=100$ N/m, biên độ $8\,\text{cm}$. Tính cơ năng; động năng và thế năng tại $x=4$ cm.
\end{bt}

% ===== Câu 159 =====
% ID: L11C1B5-39-TLN
% Mức: NB
\dangbt{Tính động năng, thế năng và cơ năng tại một li độ hoặc thời điểm xác định}
\begin{ex}
% ID: L11C1B5-39-TLN
% Mức: NB
Một con lắc đơn có khối lượng của vật nặng là $200\,\text{g}$ dao động với phương trình $s=10\sin(2\pi t)\,\text{cm}$. Ở thời điểm $t=\dfrac{\pi}{6}\,\text{s}$, con lắc có động năng là bao nhiêu mJ?
\shortans{1}
\loigiai{
Vận tốc của vật:
$v=s'=0,1\cdot2\pi\cos(2\pi t).$
Theo kết quả tính toán:
$v=-0,1\,\text{m/s}.$
Động năng của con lắc:
$W_{\text{đ}}=\dfrac{1}{2}mv^2 =\dfrac{1}{2}\cdot0,2\cdot(-0,1)^2 =0,001\,\text{J}=1\,\text{mJ}.$
}
\end{ex}

% ===== Câu 160 =====
% ID: L11C1B5-39-TN
% Mức: NB
\dangbt{Phân tích đồ thị năng lượng}
\begin{ex}
% ID: L11C1B5-39-TN
% Mức: NB
Đồ thị động năng theo li độ $W_{\text{đ}}(x)$ của một vật dao động điều hòa là một đường parabol có bề lõm hướng xuống. Điểm cao nhất của đồ thị này (đỉnh parabol) ứng với vị trí:
\choice
{Vị trí biên dương $x = A$}
{Vị trí biên âm $x = -A$}
{\True Vị trí cân bằng $x = 0$}
{Vị trí có li độ $x = \pm \dfrac{A}{2}$}
\loigiai{
• Công thức xác định động năng theo li độ:
 $W_{\text{đ}} = W - W_{\text{t}} = \dfrac{1}{2}m\omega^2(A^2 - x^2).$
 
• Đỉnh của đường parabol ứng với giá trị cực đại của động năng $W_{\text{đ}}$.
 
• Động năng cực đại khi $x^2 = 0 \Rightarrow x = 0$, tức là khi vật đi qua vị trí cân bằng.
}
\end{ex}

% ===== Câu 161 =====
% ID: L11C1B5-40-DS
% Mức: VD
\begin{ex}
% ID: L11C1B5-40-DS
% Mức: VD
Cho đồ thị biểu diễn sự phụ thuộc của thế năng $W_t$ và động năng $W_d$ vào li độ $x$ của một con lắc lò xo. Trên đồ thị, điểm giao nhau của hai đường có tọa độ li độ là $x = \pm 4$ cm và giá trị năng lượng là $0,16$ J.
\choiceTF
{\True Tại điểm giao nhau của hai đồ thị, thế năng bằng động năng và cùng bằng $0,16$ J.}
{\True Cơ năng dao động của con lắc lò xo được xác định bằng $0,32$ J.}
{Biên độ dao động của con lắc là $8$ cm.}
{\True Độ cứng của lò xo con lắc dùng trong khảo sát bằng $200$ N/m.}
\loigiai{
\begin{itemchoice}
\item (Đúng) Tại điểm giao nhau của hai đồ thị, thế năng bằng động năng và cùng bằng $0,16$ J. (Vì): Mệnh đề này đúng vì giao điểm của hai đồ thị là vị trí thế năng bằng động năng, tung độ giao điểm cho biết giá trị năng lượng thành phần bằng $0,16$ J
\item (Đúng) Cơ năng dao động của con lắc lò xo được xác định bằng $0,32$ J. (Vì): Mệnh đề này đúng vì tại vị trí thế năng bằng động năng, cơ năng của vật được tính là: $W = W_d + W_t = 2W_t = 2 \cdot 0,16 = 0,32$ J
\item (Sai) Biên độ dao động của con lắc là $8$ cm. (Vì): Mệnh đề này sai vì tại vị trí $W_d = W_t \Rightarrow x = \pm \dfrac{A}{\sqrt{2}} \Rightarrow 4 = \dfrac{A}{\sqrt{2}} \Rightarrow A = 4\sqrt{2}$ cm $\approx 5,66$ cm
\item (Đúng) Độ cứng của lò xo con lắc dùng trong khảo sát bằng $200$ N/m. (Vì): Mệnh đề này đúng vì đổi $x = 0,04$ m; áp dụng thế năng $W_t = \dfrac{1}{2}kx^2 \Rightarrow 0,16 = 0,5 \cdot k \cdot 0,04^2 \Rightarrow k = 200$ N/m
\end{itemchoice}
}
\end{ex}

% ===== Câu 162 =====
% ID: L11C1B5-40-TL
% Mức: TH
\dangbt{Xác định độ cứng lò xo từ chu kì biến thiên năng lượng}
\begin{ex}
% ID: L11C1B5-40-TL
% Mức: TH
Con lắc lò xo có cơ năng 0.04 $J$ và biên độ $4\,\text{cm}$. Tính độ cứng lò xo và nêu công thức sử dụng.
\choiceTF

\end{ex}

% ===== Câu 163 =====
% ID: L11C1B5-40-TLN
% Mức: NB
\dangbt{Tính động năng, thế năng và cơ năng tại một li độ hoặc thời điểm xác định}
\begin{ex}
% ID: L11C1B5-40-TLN
% Mức: NB
Con lắc lò xo nằm ngang có $k=100\,\text{N/m}$, $m=1\,\text{kg}$ dao động điều hòa. Khi vật có động năng $10\,\text{mJ}$ thì cách vị trí cân bằng $1\,\text{cm}$; khi có động năng $5\,\text{mJ}$ thì cách vị trí cân bằng một đoạn là bao nhiêu cm? Kết quả làm tròn đến hai chữ số thập phân.
\shortans{1.41}
\loigiai{
Áp dụng bảo toàn cơ năng:
$W_{\text{đ}1}+W_{\text{t}1}=W_{\text{đ}2}+W_{\text{t}2}.$
Suy ra
$10\cdot10^{-3}+\dfrac{1}{2}\cdot100\cdot(0,01)^2 =5\cdot10^{-3}+\dfrac{1}{2}\cdot100\cdot x_2^2.$
Do đó
$x_2=\sqrt{2}\,\text{cm}\approx1,41\,\text{cm}.$
}
\end{ex}

% ===== Câu 164 =====
% ID: L11C1B5-40-TN
% Mức: TH
\dangbt{Sự chuyển hóa năng lượng trong dao động điều hòa}
\begin{ex}
% ID: L11C1B5-40-TN
% Mức: TH
Một vật dao động điều hòa đi từ vị trí biên về vị trí cân bằng. Trong quá trình đó, động năng và thế năng của vật biến đổi như thế nào?
\choice
{Động năng giảm, thế năng tăng.}
{\True Động năng tăng, thế năng giảm.}
{Cả động năng và thế năng đều tăng.}
{Cả động năng và thế năng đều giảm.}
\loigiai{
Khi vật đi từ biên về vị trí cân bằng, thế năng giảm dần và chuyển hóa thành động năng nên động năng tăng dần.
}
\end{ex}

% ===== Câu 165 =====
% ID: L11C1B5-41-DS
% Mức: VD
\dangbt{Phân tích đồ thị năng lượng}
\begin{ex}
% ID: L11C1B5-41-DS
% Mức: VD
Cho đồ thị biểu diễn sự phụ thuộc của thế năng $W_t$ và động năng $W_d$ vào li độ $x$ của một con lắc lò xo. Trên đồ thị, điểm giao nhau của hai đường có tọa độ li độ là $x = \pm 4$ cm và giá trị năng lượng là $0,16$ J.
\choiceTF
{\True Tại điểm giao nhau của hai đồ thị, thế năng bằng động năng và cùng bằng $0,16$ J.}
{\True Cơ năng dao động của con lắc lò xo được xác định bằng $0,32$ J.}
{Biên độ dao động của con lắc là $8$ cm.}
{\True Độ cứng của lò xo con lắc dùng trong khảo sát bằng $200$ N/m.}
\loigiai{
\begin{itemchoice}
\item (Đúng) Tại điểm giao nhau của hai đồ thị, thế năng bằng động năng và cùng bằng $0,16$ J. (Vì): Mệnh đề này đúng vì giao điểm của hai đồ thị là vị trí thế năng bằng động năng, tung độ giao điểm cho biết giá trị năng lượng thành phần bằng $0,16$ J
\item (Đúng) Cơ năng dao động của con lắc lò xo được xác định bằng $0,32$ J. (Vì): Mệnh đề này đúng vì tại vị trí thế năng bằng động năng, cơ năng của vật được tính là: $W = W_d + W_t = 2W_t = 2 \cdot 0,16 = 0,32$ J
\item (Sai) Biên độ dao động của con lắc là $8$ cm. (Vì): Mệnh đề này sai vì tại vị trí $W_d = W_t \Rightarrow x = \pm \dfrac{A}{\sqrt{2}} \Rightarrow 4 = \dfrac{A}{\sqrt{2}} \Rightarrow A = 4\sqrt{2}$ cm $\approx 5,66$ cm
\item (Đúng) Độ cứng của lò xo con lắc dùng trong khảo sát bằng $200$ N/m. (Vì): Mệnh đề này đúng vì đổi $x = 0,04$ m; áp dụng thế năng $W_t = \dfrac{1}{2}kx^2 \Rightarrow 0,16 = 0,5 \cdot k \cdot 0,04^2 \Rightarrow k = 200$ N/m
\end{itemchoice}
}
\end{ex}

% ===== Câu 166 =====
% ID: L11C1B5-41-TL
% Mức: VD
\dangbt{Xác định độ cứng lò xo từ chu kì biến thiên năng lượng}
\begin{ex}
% ID: L11C1B5-41-TL
% Mức: VD
Con lắc lò xo có cơ năng 0.09 $J$ và biên độ $6\,\text{cm}$. Tính độ cứng lò xo và nêu công thức sử dụng.
\choiceTF

\end{ex}

% ===== Câu 167 =====
% ID: L11C1B5-41-TLN
% Mức: NB
\dangbt{Tính động năng, thế năng và cơ năng tại một li độ hoặc thời điểm xác định}
\begin{ex}
% ID: L11C1B5-41-TLN
% Mức: NB
Con lắc lò xo nằm ngang có $k=100\,\text{N/m}$, $m=1\,\text{kg}$ dao động điều hòa. Khi vật có động năng $10\,\text{mJ}$ thì cách vị trí cân bằng $1\,\text{cm}$; khi có động năng $5\,\text{mJ}$ thì cách vị trí cân bằng một đoạn là bao nhiêu cm? Kết quả làm tròn đến hai chữ số thập phân.
\shortans{1.41}
\loigiai{
Áp dụng bảo toàn cơ năng:
$W_{\text{đ}1}+W_{\text{t}1}=W_{\text{đ}2}+W_{\text{t}2}.$
Suy ra
$10\cdot10^{-3}+\dfrac{1}{2}\cdot100\cdot(0,01)^2 =5\cdot10^{-3}+\dfrac{1}{2}\cdot100\cdot x_2^2.$
Do đó
$x_2=\sqrt{2}\,\text{cm}\approx1,41\,\text{cm}.$
}
\end{ex}

% ===== Câu 168 =====
% ID: L11C1B5-41-TN
% Mức: TH
\dangbt{Sự chuyển hóa năng lượng trong dao động điều hòa}
\begin{ex}
% ID: L11C1B5-41-TN
% Mức: TH
Một vật dao động điều hòa đi từ vị trí biên về vị trí cân bằng. Trong quá trình đó, động năng và thế năng của vật biến đổi như thế nào?
\choice
{Động năng giảm, thế năng tăng.}
{\True Động năng tăng, thế năng giảm.}
{Cả động năng và thế năng đều tăng.}
{Cả động năng và thế năng đều giảm.}
\loigiai{
Khi vật đi từ biên về vị trí cân bằng, thế năng giảm dần và chuyển hóa thành động năng nên động năng tăng dần.
}
\end{ex}

% ===== Câu 169 =====
% ID: L11C1B5-42-DS
% Mức: TH
\begin{ex}
% ID: L11C1B5-42-DS
% Mức: TH
Một vật dao động điều hòa có cơ năng $W$. Tại li độ $x=A/2$, xác định tính đúng sai.
\choiceTF
{\True Thế năng bằng $W/4$.}
{\True Động năng bằng $3W/4$.}
{\True Động năng gấp ba lần thế năng.}
{Cơ năng giảm còn $W/2$.}
\end{ex}

% ===== Câu 170 =====
% ID: L11C1B5-42-TL
% Mức: VD
\dangbt{Xác định độ cứng lò xo từ chu kì biến thiên năng lượng}
\begin{ex}
% ID: L11C1B5-42-TL
% Mức: VD
Con lắc lò xo có cơ năng 0.2 $J$ và biên độ $10\,\text{cm}$. Tính độ cứng lò xo và nêu công thức sử dụng.
\choiceTF

\end{ex}

% ===== Câu 171 =====
% ID: L11C1B5-42-TLN
% Mức: NB
\dangbt{Tính động năng, thế năng và cơ năng tại một li độ hoặc thời điểm xác định}
\begin{ex}
% ID: L11C1B5-42-TLN
% Mức: NB
Con lắc lò xo nằm ngang có $k=100\,\text{N/m}$, $m=1\,\text{kg}$ dao động điều hòa. Khi vật có động năng $10\,\text{mJ}$ thì cách vị trí cân bằng $1\,\text{cm}$; khi có động năng $5\,\text{mJ}$ thì cách vị trí cân bằng một đoạn là bao nhiêu cm? Kết quả làm tròn đến hai chữ số thập phân.
\shortans{1.41}
\loigiai{
Áp dụng bảo toàn cơ năng:
$W_{\text{đ}1}+W_{\text{t}1}=W_{\text{đ}2}+W_{\text{t}2}.$
Suy ra
$10\cdot10^{-3}+\dfrac{1}{2}\cdot100\cdot(0,01)^2 =5\cdot10^{-3}+\dfrac{1}{2}\cdot100\cdot x_2^2.$
Do đó
$x_2=\sqrt{2}\,\text{cm}\approx1,41\,\text{cm}.$
}
\end{ex}

% ===== Câu 172 =====
% ID: L11C1B5-42-TN
% Mức: TH
\dangbt{Sự chuyển hóa năng lượng trong dao động điều hòa}
\begin{ex}
% ID: L11C1B5-42-TN
% Mức: TH
Một vật dao động điều hòa đi từ vị trí biên về vị trí cân bằng. Trong quá trình đó, động năng và thế năng của vật biến đổi như thế nào?
\choice
{Động năng giảm, thế năng tăng.}
{\True Động năng tăng, thế năng giảm.}
{Cả động năng và thế năng đều tăng.}
{Cả động năng và thế năng đều giảm.}
\loigiai{
Khi vật đi từ biên về vị trí cân bằng, thế năng giảm dần và chuyển hóa thành động năng nên động năng tăng dần.
}
\end{ex}

% ===== Câu 173 =====
% ID: L11C1B5-43-DS
% Mức: VD
\begin{ex}
% ID: L11C1B5-43-DS
% Mức: VD
Một vật dao động điều hòa có cơ năng $W$. Tại li độ $x=A/2$, xác định tính đúng sai.
\choiceTF
{\True Thế năng bằng $W/4$.}
{\True Động năng bằng $3W/4$.}
{\True Động năng gấp ba lần thế năng.}
{Cơ năng giảm còn $W/2$.}
\end{ex}

% ===== Câu 174 =====
% ID: L11C1B5-43-TL
% Mức: TH
\dangbt{Xác định độ cứng lò xo từ cơ năng và thông số dao động}
\begin{ex}
% ID: L11C1B5-43-TL
% Mức: TH
Con lắc lò xo có cơ năng 0.04 $J$ và biên độ $4\,\text{cm}$. Tính độ cứng lò xo và nêu công thức sử dụng.
\choiceTF

\end{ex}

% ===== Câu 175 =====
% ID: L11C1B5-43-TLN
% Mức: NB
\dangbt{Tính động năng, thế năng và cơ năng tại một li độ hoặc thời điểm xác định}
\begin{ex}
% ID: L11C1B5-43-TLN
% Mức: NB
Vật nhỏ của một con lắc lò xo dao động điều hòa theo phương ngang, mốc thế năng tại vị trí cân bằng. Khi gia tốc của vật có độ lớn bằng một nửa gia tốc cực đại thì tỉ số giữa động năng và thế năng của vật là bao nhiêu?
\shortans{3}
\loigiai{
Ta có:
$|a|=\dfrac{a_{\max}}{2} \Rightarrow \omega^2|x|=\dfrac{\omega^2A}{2} \Rightarrow |x|=\dfrac{A}{2}.$
Tỉ số giữa động năng và thế năng:
$\dfrac{W_{\text{đ}}}{W_t} =\dfrac{\dfrac{1}{2}k(A^2-x^2)}{\dfrac{1}{2}kx^2} =\dfrac{A^2-x^2}{x^2} =\dfrac{A^2-\left(\dfrac{A}{2}\right)^2}{\left(\dfrac{A}{2}\right)^2} =3.$
}
\end{ex}

% ===== Câu 176 =====
% ID: L11C1B5-43-TN
% Mức: TH
\dangbt{Sự chuyển hóa năng lượng trong dao động điều hòa}
\begin{ex}
% ID: L11C1B5-43-TN
% Mức: TH
Tại vị trí biên của một vật dao động điều hòa, nhận xét nào sau đây đúng?
\choice
{Động năng cực đại, thế năng bằng 0.}
{Động năng bằng thế năng.}
{Động năng và thế năng đều bằng 0.}
{\True Động năng bằng 0, thế năng cực đại.}
\loigiai{
Ở vị trí biên, vận tốc bằng 0 nên động năng bằng 0; thế năng bằng cơ năng và đạt cực đại.
}
\end{ex}

% ===== Câu 177 =====
% ID: L11C1B5-44-DS
% Mức: VD
\begin{ex}
% ID: L11C1B5-44-DS
% Mức: VD
Một vật dao động điều hòa có cơ năng $W$. Tại li độ $x=A/2$, xác định tính đúng sai.
\choiceTF
{\True Thế năng bằng $W/4$.}
{\True Động năng bằng $3W/4$.}
{\True Động năng gấp ba lần thế năng.}
{Cơ năng giảm còn $W/2$.}
\end{ex}

% ===== Câu 178 =====
% ID: L11C1B5-44-TL
% Mức: VD
\dangbt{Xác định độ cứng lò xo từ cơ năng và thông số dao động}
\begin{ex}
% ID: L11C1B5-44-TL
% Mức: VD
Con lắc lò xo có cơ năng 0.09 $J$ và biên độ $6\,\text{cm}$. Tính độ cứng lò xo và nêu công thức sử dụng.
\choiceTF

\end{ex}

% ===== Câu 179 =====
% ID: L11C1B5-44-TLN
% Mức: NB
\dangbt{Tính động năng, thế năng và cơ năng tại một li độ hoặc thời điểm xác định}
\begin{ex}
% ID: L11C1B5-44-TLN
% Mức: NB
Vật nhỏ của một con lắc lò xo dao động điều hòa theo phương ngang, mốc thế năng tại vị trí cân bằng. Khi gia tốc của vật có độ lớn bằng một nửa gia tốc cực đại thì tỉ số giữa động năng và thế năng của vật là bao nhiêu?
\shortans{3}
\loigiai{
Ta có:
$|a|=\dfrac{a_{\max}}{2} \Rightarrow \omega^2|x|=\dfrac{\omega^2A}{2} \Rightarrow |x|=\dfrac{A}{2}.$
Tỉ số giữa động năng và thế năng:
$\dfrac{W_{\text{đ}}}{W_t} =\dfrac{\dfrac{1}{2}k(A^2-x^2)}{\dfrac{1}{2}kx^2} =\dfrac{A^2-x^2}{x^2} =\dfrac{A^2-\left(\dfrac{A}{2}\right)^2}{\left(\dfrac{A}{2}\right)^2} =3.$
}
\end{ex}

% ===== Câu 180 =====
% ID: L11C1B5-44-TN
% Mức: TH
\dangbt{Sự chuyển hóa năng lượng trong dao động điều hòa}
\begin{ex}
% ID: L11C1B5-44-TN
% Mức: TH
Tại vị trí biên của một vật dao động điều hòa, nhận xét nào sau đây đúng?
\choice
{Động năng cực đại, thế năng bằng 0.}
{Động năng bằng thế năng.}
{Động năng và thế năng đều bằng 0.}
{\True Động năng bằng 0, thế năng cực đại.}
\loigiai{
Ở vị trí biên, vận tốc bằng 0 nên động năng bằng 0; thế năng bằng cơ năng và đạt cực đại.
}
\end{ex}

% ===== Câu 181 =====
% ID: L11C1B5-45-DS
% Mức: TH
\dangbt{Tính cơ năng của con lắc lò xo từ biên độ dao động}
\begin{ex}
% ID: L11C1B5-45-DS
% Mức: TH
Một vật dao động điều hòa có cơ năng $W$. Tại li độ $x=A/2$, xác định tính đúng sai.
\choiceTF
{\True Thế năng bằng $W/4$.}
{\True Động năng bằng $3W/4$.}
{\True Động năng gấp ba lần thế năng.}
{Cơ năng giảm còn $W/2$.}
\end{ex}

% ===== Câu 182 =====
% ID: L11C1B5-45-TL
% Mức: VD
\dangbt{Xác định độ cứng lò xo từ cơ năng và thông số dao động}
\begin{ex}
% ID: L11C1B5-45-TL
% Mức: VD
Con lắc lò xo có cơ năng 0.2 $J$ và biên độ $10\,\text{cm}$. Tính độ cứng lò xo và nêu công thức sử dụng.
\choiceTF

\end{ex}

% ===== Câu 183 =====
% ID: L11C1B5-45-TLN
% Mức: NB
\dangbt{Tính động năng, thế năng và cơ năng tại một li độ hoặc thời điểm xác định}
\begin{ex}
% ID: L11C1B5-45-TLN
% Mức: NB
Vật nhỏ của một con lắc lò xo dao động điều hòa theo phương ngang, mốc thế năng tại vị trí cân bằng. Khi gia tốc của vật có độ lớn bằng một nửa gia tốc cực đại thì tỉ số giữa động năng và thế năng của vật là bao nhiêu?
\shortans{3}
\loigiai{
Ta có:
$|a|=\dfrac{a_{\max}}{2} \Rightarrow \omega^2|x|=\dfrac{\omega^2A}{2} \Rightarrow |x|=\dfrac{A}{2}.$
Tỉ số giữa động năng và thế năng:
$\dfrac{W_{\text{đ}}}{W_t} =\dfrac{\dfrac{1}{2}k(A^2-x^2)}{\dfrac{1}{2}kx^2} =\dfrac{A^2-x^2}{x^2} =\dfrac{A^2-\left(\dfrac{A}{2}\right)^2}{\left(\dfrac{A}{2}\right)^2} =3.$
}
\end{ex}

% ===== Câu 184 =====
% ID: L11C1B5-45-TN
% Mức: TH
\dangbt{Sự chuyển hóa năng lượng trong dao động điều hòa}
\begin{ex}
% ID: L11C1B5-45-TN
% Mức: TH
Tại vị trí biên của một vật dao động điều hòa, nhận xét nào sau đây đúng?
\choice
{Động năng cực đại, thế năng bằng 0.}
{Động năng bằng thế năng.}
{Động năng và thế năng đều bằng 0.}
{\True Động năng bằng 0, thế năng cực đại.}
\loigiai{
Ở vị trí biên, vận tốc bằng 0 nên động năng bằng 0; thế năng bằng cơ năng và đạt cực đại.
}
\end{ex}

% ===== Câu 185 =====
% ID: L11C1B5-46-DS
% Mức: VD
\dangbt{Tính cơ năng của con lắc lò xo từ biên độ dao động}
\begin{ex}
% ID: L11C1B5-46-DS
% Mức: VD
Một vật dao động điều hòa có cơ năng $W$. Tại li độ $x=A/2$, xác định tính đúng sai.
\choiceTF
{\True Thế năng bằng $W/4$.}
{\True Động năng bằng $3W/4$.}
{\True Động năng gấp ba lần thế năng.}
{Cơ năng giảm còn $W/2$.}
\end{ex}

% ===== Câu 186 =====
% ID: L11C1B5-46-TL
% Mức: TH
\dangbt{Xác định độ cứng lò xo từ năng lượng ban đầu}
\begin{ex}
% ID: L11C1B5-46-TL
% Mức: TH
Con lắc lò xo có cơ năng 0.04 $J$ và biên độ $4\,\text{cm}$. Tính độ cứng lò xo và nêu công thức sử dụng.
\choiceTF

\end{ex}

% ===== Câu 187 =====
% ID: L11C1B5-46-TLN
% Mức: NB
\dangbt{Tính động năng, thế năng và cơ năng tại một li độ hoặc thời điểm xác định}
\begin{ex}
% ID: L11C1B5-46-TLN
% Mức: NB
Một con lắc lò xo gồm một vật nhỏ có độ cứng $20\,\text{N/m}$ dao động điều hòa với chu kì $2\,\text{s}$. Khi pha của dao động là $\dfrac{\pi}{2}$ thì vận tốc của vật là $-20\sqrt{3}\,\text{cm/s}$. Lấy $\pi^2=10$, khi vật đi qua vị trí có li độ $3\pi\,\text{cm}$ thì động năng của con lắc là bao nhiêu J?
\shortans{0.03}
\loigiai{
Tần số góc:
$\omega=\dfrac{2\pi}{T}=\pi\,\text{rad/s}.$
Khi $\varphi=\dfrac{\pi}{2}$:
$v=-\omega A\sin\dfrac{\pi}{2}=-20\sqrt{3} \Rightarrow A=\dfrac{20\sqrt{3}}{\pi}\,\text{cm}.$
Khi $x=3\pi\,\text{cm}$, động năng:
$W_{\text{đ}}=\dfrac{1}{2}k(A^2-x^2)=0,03\,\text{J}.$
}
\end{ex}

% ===== Câu 188 =====
% ID: L11C1B5-46-TN
% Mức: TH
\dangbt{Tính cơ năng của con lắc lò xo từ biên độ dao động}
\begin{ex}
% ID: L11C1B5-46-TN
% Mức: TH
Con lắc lò xo có $k=50$ N/m, biên độ $4\,\text{cm}$. Cơ năng bằng
\choice
{\True 0.040 $J$}
{0.080 $J$}
{0.020 $J$}
{0.160 $J$}
\end{ex}

% ===== Câu 189 =====
% ID: L11C1B5-47-DS
% Mức: VD
\begin{ex}
% ID: L11C1B5-47-DS
% Mức: VD
Một vật dao động điều hòa có cơ năng $W$. Tại li độ $x=A/2$, xác định tính đúng sai.
\choiceTF
{\True Thế năng bằng $W/4$.}
{\True Động năng bằng $3W/4$.}
{\True Động năng gấp ba lần thế năng.}
{Cơ năng giảm còn $W/2$.}
\end{ex}

% ===== Câu 190 =====
% ID: L11C1B5-47-TL
% Mức: VD
\dangbt{Xác định độ cứng lò xo từ năng lượng ban đầu}
\begin{ex}
% ID: L11C1B5-47-TL
% Mức: VD
Con lắc lò xo có cơ năng 0.09 $J$ và biên độ $6\,\text{cm}$. Tính độ cứng lò xo và nêu công thức sử dụng.
\choiceTF

\end{ex}

% ===== Câu 191 =====
% ID: L11C1B5-47-TLN
% Mức: NB
\dangbt{Tính động năng, thế năng và cơ năng tại một li độ hoặc thời điểm xác định}
\begin{ex}
% ID: L11C1B5-47-TLN
% Mức: NB
Một con lắc lò xo gồm một vật nhỏ có độ cứng $20\,\text{N/m}$ dao động điều hòa với chu kì $2\,\text{s}$. Khi pha của dao động là $\dfrac{\pi}{2}$ thì vận tốc của vật là $-20\sqrt{3}\,\text{cm/s}$. Lấy $\pi^2=10$, khi vật đi qua vị trí có li độ $3\pi\,\text{cm}$ thì động năng của con lắc là bao nhiêu J?
\shortans{0.03}
\loigiai{
Tần số góc:
$\omega=\dfrac{2\pi}{T}=\pi\,\text{rad/s}.$
Khi $\varphi=\dfrac{\pi}{2}$:
$v=-\omega A\sin\dfrac{\pi}{2}=-20\sqrt{3} \Rightarrow A=\dfrac{20\sqrt{3}}{\pi}\,\text{cm}.$
Khi $x=3\pi\,\text{cm}$, động năng:
$W_{\text{đ}}=\dfrac{1}{2}k(A^2-x^2)=0,03\,\text{J}.$
}
\end{ex}

% ===== Câu 192 =====
% ID: L11C1B5-47-TN
% Mức: VD
\dangbt{Tính cơ năng của con lắc lò xo từ biên độ dao động}
\begin{ex}
% ID: L11C1B5-47-TN
% Mức: VD
Con lắc lò xo có $k=80$ N/m, biên độ $6\,\text{cm}$. Cơ năng bằng
\choice
{\True 0.144 $J$}
{0.288 $J$}
{0.072 $J$}
{0.576 $J$}
\end{ex}

% ===== Câu 193 =====
% ID: L11C1B5-48-DS
% Mức: TH
\dangbt{Tính cơ năng của con lắc đơn}
\begin{ex}
% ID: L11C1B5-48-DS
% Mức: TH
Con lắc đơn dao động điều hòa biên độ góc $\alpha_0$. Xác định tính đúng sai.
\choiceTF
{\True Cơ năng xấp xỉ $mgl\alpha_0^2/2$.}
{\True Ở vị trí cân bằng, động năng cực đại.}
{\True Ở vị trí biên, tốc độ bằng 0.}
{Cơ năng biến thiên tuần hoàn.}
\end{ex}

% ===== Câu 194 =====
% ID: L11C1B5-48-TL
% Mức: VD
\dangbt{Xác định độ cứng lò xo từ năng lượng ban đầu}
\begin{ex}
% ID: L11C1B5-48-TL
% Mức: VD
Con lắc lò xo có cơ năng 0.2 $J$ và biên độ $10\,\text{cm}$. Tính độ cứng lò xo và nêu công thức sử dụng.
\choiceTF

\end{ex}

% ===== Câu 195 =====
% ID: L11C1B5-48-TLN
% Mức: NB
\dangbt{Tính động năng, thế năng và cơ năng tại một li độ hoặc thời điểm xác định}
\begin{ex}
% ID: L11C1B5-48-TLN
% Mức: NB
Một con lắc lò xo gồm một vật nhỏ có độ cứng $20\,\text{N/m}$ dao động điều hòa với chu kì $2\,\text{s}$. Khi pha của dao động là $\dfrac{\pi}{2}$ thì vận tốc của vật là $-20\sqrt{3}\,\text{cm/s}$. Lấy $\pi^2=10$, khi vật đi qua vị trí có li độ $3\pi\,\text{cm}$ thì động năng của con lắc là bao nhiêu J?
\shortans{0.03}
\loigiai{
Tần số góc:
$\omega=\dfrac{2\pi}{T}=\pi\,\text{rad/s}.$
Khi $\varphi=\dfrac{\pi}{2}$:
$v=-\omega A\sin\dfrac{\pi}{2}=-20\sqrt{3} \Rightarrow A=\dfrac{20\sqrt{3}}{\pi}\,\text{cm}.$
Khi $x=3\pi\,\text{cm}$, động năng:
$W_{\text{đ}}=\dfrac{1}{2}k(A^2-x^2)=0,03\,\text{J}.$
}
\end{ex}

% ===== Câu 196 =====
% ID: L11C1B5-48-TN
% Mức: VD
\dangbt{Tính cơ năng của con lắc lò xo từ biên độ dao động}
\begin{ex}
% ID: L11C1B5-48-TN
% Mức: VD
Con lắc lò xo có $k=100$ N/m, biên độ $8\,\text{cm}$. Cơ năng bằng
\choice
{\True 0.320 $J$}
{0.640 $J$}
{0.160 $J$}
{1.280 $J$}
\end{ex}

% ===== Câu 197 =====
% ID: L11C1B5-49-DS
% Mức: VD
\dangbt{Tính cơ năng của con lắc đơn}
\begin{ex}
% ID: L11C1B5-49-DS
% Mức: VD
Con lắc đơn dao động điều hòa biên độ góc $\alpha_0$. Xác định tính đúng sai.
\choiceTF
{\True Cơ năng xấp xỉ $mgl\alpha_0^2/2$.}
{\True Ở vị trí cân bằng, động năng cực đại.}
{\True Ở vị trí biên, tốc độ bằng 0.}
{Cơ năng biến thiên tuần hoàn.}
\end{ex}

% ===== Câu 198 =====
% ID: L11C1B5-49-TLN
% Mức: NB
\dangbt{Tính động năng, thế năng và cơ năng tại một li độ hoặc thời điểm xác định}
\begin{ex}
% ID: L11C1B5-49-TLN
% Mức: NB
Con lắc lò xo có khối lượng $m=400\,\text{g}$, độ cứng $k=160\,\text{N/m}$ dao động điều hòa theo phương thẳng đứng. Biết khi vật có li độ $2\,\text{cm}$ thì vận tốc của vật bằng $40\,\text{cm/s}$. Năng lượng dao động của vật là bao nhiêu mJ?
\shortans{64}
\loigiai{
Đổi $m=0,4\,\text{kg}$, $x=0,02\,\text{m}$, $v=0,4\,\text{m/s}$.
Năng lượng dao động của vật:
$W=\dfrac{1}{2}kx^2+\dfrac{1}{2}mv^2 =\dfrac{1}{2}\cdot160\cdot(0,02)^2+\dfrac{1}{2}\cdot0,4\cdot(0,4)^2 =0,064\,\text{J}=64\,\text{mJ}.$
}
\end{ex}

% ===== Câu 199 =====
% ID: L11C1B5-49-TN
% Mức: TH
\dangbt{Tính cơ năng của con lắc đơn}
\begin{ex}
% ID: L11C1B5-49-TN
% Mức: TH
Con lắc đơn có $m=0.2$ kg, $l=1$ m, biên độ góc 0.1 rad, lấy $g=10$ m/s $^2$. Cơ năng bằng
\choice
{\True 0.010 $J$}
{0.020 $J$}
{0.005 $J$}
{0.040 $J$}
\end{ex}

% ===== Câu 200 =====
% ID: L11C1B5-50-DS
% Mức: VD
\begin{ex}
% ID: L11C1B5-50-DS
% Mức: VD
Con lắc đơn dao động điều hòa biên độ góc $\alpha_0$. Xác định tính đúng sai.
\choiceTF
{\True Cơ năng xấp xỉ $mgl\alpha_0^2/2$.}
{\True Ở vị trí cân bằng, động năng cực đại.}
{\True Ở vị trí biên, tốc độ bằng 0.}
{Cơ năng biến thiên tuần hoàn.}
\end{ex}

% ===== Câu 201 =====
% ID: L11C1B5-50-TLN
% Mức: NB
\dangbt{Tính động năng, thế năng và cơ năng tại một li độ hoặc thời điểm xác định}
\begin{ex}
% ID: L11C1B5-50-TLN
% Mức: NB
Con lắc lò xo có khối lượng $m=400\,\text{g}$, độ cứng $k=160\,\text{N/m}$ dao động điều hòa theo phương thẳng đứng. Biết khi vật có li độ $2\,\text{cm}$ thì vận tốc của vật bằng $40\,\text{cm/s}$. Năng lượng dao động của vật là bao nhiêu mJ?
\shortans{64}
\loigiai{
Đổi $m=0,4\,\text{kg}$, $x=0,02\,\text{m}$, $v=0,4\,\text{m/s}$.
Năng lượng dao động của vật:
$W=\dfrac{1}{2}kx^2+\dfrac{1}{2}mv^2 =\dfrac{1}{2}\cdot160\cdot(0,02)^2+\dfrac{1}{2}\cdot0,4\cdot(0,4)^2 =0,064\,\text{J}=64\,\text{mJ}.$
}
\end{ex}

% ===== Câu 202 =====
% ID: L11C1B5-50-TN
% Mức: VD
\dangbt{Tính cơ năng của con lắc đơn}
\begin{ex}
% ID: L11C1B5-50-TN
% Mức: VD
Con lắc đơn có $m=0.25$ kg, $l=0.9$ m, biên độ góc 0.2 rad, lấy $g=10$ m/s $^2$. Cơ năng bằng
\choice
{\True 0.045 $J$}
{0.090 $J$}
{0.023 $J$}
{0.180 $J$}
\end{ex}

% ===== Câu 203 =====
% ID: L11C1B5-51-DS
% Mức: TH
\dangbt{Tính động năng, thế năng và cơ năng tại một li độ hoặc thời điểm xác định}
\begin{ex}
% ID: L11C1B5-51-DS
% Mức: TH
Xét chuyển động và sự chuyển hóa năng lượng qua lại giữa động năng và thế năng của một chất điểm dao động điều hòa trên trục Ox quanh vị trí cân bằng O:
\choiceTF
{Khi chất điểm chuyển động từ vị trí cân bằng ra vị trí biên, thế năng của chất điểm giảm dần.}
{\True Khi chất điểm chuyển động từ vị trí cân bằng ra vị trí biên, tốc độ của vật giảm dần nên động năng giảm dần.}
{Lực kéo về luôn thực hiện công dương trong mọi giai đoạn chuyển động của chất điểm.}
{\True Tại vị trí cân bằng, thế năng của chất điểm bằng không (chọn mốc thế năng tại vị trí cân bằng) và động năng của chất điểm đạt giá trị cực đại.}
\loigiai{
\begin{itemchoice}
\item (Sai) Khi chất điểm chuyển động từ vị trí cân bằng ra vị trí biên, thế năng của chất điểm giảm dần. (Vì): Mệnh đề này sai vì khi vật đi từ vị trí cân bằng ra biên, khoảng cách đến vị trí cân bằng tăng dần nên thế năng của vật tăng dần
\item (Đúng) Khi chất điểm chuyển động từ vị trí cân bằng ra vị trí biên, tốc độ của vật giảm dần nên động năng giảm dần. (Vì): Mệnh đề này đúng vì khi đi từ vị trí cân bằng ra biên, vật chuyển động chậm dần, vận tốc giảm về không kéo theo động năng giảm dần
\item (Sai) Lực kéo về luôn thực hiện công dương trong mọi giai đoạn chuyển động của chất điểm. (Vì): Mệnh đề này sai vì khi vật đi từ vị trí cân bằng ra biên, lực kéo về hướng về vị trí cân bằng nhưng vật chuyển động ra biên, lực ngược chiều chuyển động nên thực hiện công âm
\item (Đúng) Tại vị trí cân bằng, thế năng của chất điểm bằng không (chọn mốc thế năng tại vị trí cân bằng) và động năng của chất điểm đạt giá trị cực đại. (Vì): Mệnh đề này đúng vì tại vị trí cân bằng $x = 0$ nên thế năng bằng $0$, tốc độ vật đạt cực đại nên động năng đạt cực đại
\end{itemchoice}
}
\end{ex}

% ===== Câu 204 =====
% ID: L11C1B5-51-TLN
% Mức: TH
\dangbt{Xác định biểu thức cơ năng trong dao động điều hòa}
\begin{ex}
% ID: L11C1B5-51-TLN
% Mức: TH
Con lắc lò xo có độ cứng $50\,\text{N}$ /m và biên độ $4\,\text{cm}$. Cơ năng theo đơn vị $J$ bằng bao nhiêu?
\shortans{0,04}
\end{ex}

% ===== Câu 205 =====
% ID: L11C1B5-51-TN
% Mức: VD
\dangbt{Tính cơ năng của con lắc đơn}
\begin{ex}
% ID: L11C1B5-51-TN
% Mức: VD
Con lắc đơn có $m=0.5$ kg, $l=1.6$ m, biên độ góc 0.1 rad, lấy $g=10$ m/s $^2$. Cơ năng bằng
\choice
{\True 0.040 $J$}
{0.080 $J$}
{0.020 $J$}
{0.160 $J$}
\end{ex}

% ===== Câu 206 =====
% ID: L11C1B5-52-DS
% Mức: TH
\dangbt{Tính động năng, thế năng và cơ năng tại một li độ hoặc thời điểm xác định}
\begin{ex}
% ID: L11C1B5-52-DS
% Mức: TH
Xét chuyển động và sự chuyển hóa năng lượng qua lại giữa động năng và thế năng của một chất điểm dao động điều hòa trên trục Ox quanh vị trí cân bằng O:
\choiceTF
{Khi chất điểm chuyển động từ vị trí cân bằng ra vị trí biên, thế năng của chất điểm giảm dần.}
{\True Khi chất điểm chuyển động từ vị trí cân bằng ra vị trí biên, tốc độ của vật giảm dần nên động năng giảm dần.}
{Lực kéo về luôn thực hiện công dương trong mọi giai đoạn chuyển động của chất điểm.}
{\True Tại vị trí cân bằng, thế năng của chất điểm bằng không (chọn mốc thế năng tại vị trí cân bằng) và động năng của chất điểm đạt giá trị cực đại.}
\loigiai{
\begin{itemchoice}
\item (Sai) Khi chất điểm chuyển động từ vị trí cân bằng ra vị trí biên, thế năng của chất điểm giảm dần. (Vì): Mệnh đề này sai vì khi vật đi từ vị trí cân bằng ra biên, khoảng cách đến vị trí cân bằng tăng dần nên thế năng của vật tăng dần
\item (Đúng) Khi chất điểm chuyển động từ vị trí cân bằng ra vị trí biên, tốc độ của vật giảm dần nên động năng giảm dần. (Vì): Mệnh đề này đúng vì khi đi từ vị trí cân bằng ra biên, vật chuyển động chậm dần, vận tốc giảm về không kéo theo động năng giảm dần
\item (Sai) Lực kéo về luôn thực hiện công dương trong mọi giai đoạn chuyển động của chất điểm. (Vì): Mệnh đề này sai vì khi vật đi từ vị trí cân bằng ra biên, lực kéo về hướng về vị trí cân bằng nhưng vật chuyển động ra biên, lực ngược chiều chuyển động nên thực hiện công âm
\item (Đúng) Tại vị trí cân bằng, thế năng của chất điểm bằng không (chọn mốc thế năng tại vị trí cân bằng) và động năng của chất điểm đạt giá trị cực đại. (Vì): Mệnh đề này đúng vì tại vị trí cân bằng $x = 0$ nên thế năng bằng $0$, tốc độ vật đạt cực đại nên động năng đạt cực đại
\end{itemchoice}
}
\end{ex}

% ===== Câu 207 =====
% ID: L11C1B5-52-TLN
% Mức: VD
\dangbt{Xác định biểu thức cơ năng trong dao động điều hòa}
\begin{ex}
% ID: L11C1B5-52-TLN
% Mức: VD
Con lắc lò xo có độ cứng $80\,\text{N}$ /m và biên độ $6\,\text{cm}$. Cơ năng theo đơn vị $J$ bằng bao nhiêu?
\shortans{0,144}
\end{ex}

% ===== Câu 208 =====
% ID: L11C1B5-52-TN
% Mức: NB
\dangbt{Tính động năng, thế năng và cơ năng tại một li độ hoặc thời điểm xác định}
\begin{ex}
% ID: L11C1B5-52-TN
% Mức: NB
Khi một vật dao động điều hòa đang chuyển động từ vị trí biên về vị trí cân bằng thì hiện tượng nào sau đây xảy ra?
\choice
{Động năng của vật giảm, thế năng tăng}
{\True Thế năng chuyển hóa dần thành động năng}
{Cơ năng của vật tăng liên tục}
{Động năng chuyển hóa dần thành thế năng}
\loigiai{
• Khi vật chuyển động từ biên về vị trí cân bằng, độ lớn li độ $x$ giảm dần về không, đồng thời tốc độ $v$ tăng dần lên giá trị cực đại.
 
• Do đó thế năng $W_{\text{t}} = \dfrac{1}{2}m\omega^2 x^2$ giảm dần và động năng $W_{\text{đ}} = \dfrac{1}{2}mv^2$ tăng dần.
 
• Theo định luật bảo toàn cơ năng, thế năng giảm đi bao nhiêu thì động năng tăng lên bấy nhiêu, nghĩa là thế năng chuyển hóa dần thành động năng.
}
\end{ex}

% ===== Câu 209 =====
% ID: L11C1B5-53-DS
% Mức: TH
\begin{ex}
% ID: L11C1B5-53-DS
% Mức: TH
Xét chuyển động và sự chuyển hóa năng lượng qua lại giữa động năng và thế năng của một chất điểm dao động điều hòa trên trục Ox quanh vị trí cân bằng O:
\choiceTF
{Khi chất điểm chuyển động từ vị trí cân bằng ra vị trí biên, thế năng của chất điểm giảm dần.}
{\True Khi chất điểm chuyển động từ vị trí cân bằng ra vị trí biên, tốc độ của vật giảm dần nên động năng giảm dần.}
{Lực kéo về luôn thực hiện công dương trong mọi giai đoạn chuyển động của chất điểm.}
{\True Tại vị trí cân bằng, thế năng của chất điểm bằng không (chọn mốc thế năng tại vị trí cân bằng) và động năng của chất điểm đạt giá trị cực đại.}
\loigiai{
\begin{itemchoice}
\item (Sai) Khi chất điểm chuyển động từ vị trí cân bằng ra vị trí biên, thế năng của chất điểm giảm dần. (Vì): Mệnh đề này sai vì khi vật đi từ vị trí cân bằng ra biên, khoảng cách đến vị trí cân bằng tăng dần nên thế năng của vật tăng dần
\item (Đúng) Khi chất điểm chuyển động từ vị trí cân bằng ra vị trí biên, tốc độ của vật giảm dần nên động năng giảm dần. (Vì): Mệnh đề này đúng vì khi đi từ vị trí cân bằng ra biên, vật chuyển động chậm dần, vận tốc giảm về không kéo theo động năng giảm dần
\item (Sai) Lực kéo về luôn thực hiện công dương trong mọi giai đoạn chuyển động của chất điểm. (Vì): Mệnh đề này sai vì khi vật đi từ vị trí cân bằng ra biên, lực kéo về hướng về vị trí cân bằng nhưng vật chuyển động ra biên, lực ngược chiều chuyển động nên thực hiện công âm
\item (Đúng) Tại vị trí cân bằng, thế năng của chất điểm bằng không (chọn mốc thế năng tại vị trí cân bằng) và động năng của chất điểm đạt giá trị cực đại. (Vì): Mệnh đề này đúng vì tại vị trí cân bằng $x = 0$ nên thế năng bằng $0$, tốc độ vật đạt cực đại nên động năng đạt cực đại
\end{itemchoice}
}
\end{ex}

% ===== Câu 210 =====
% ID: L11C1B5-53-TLN
% Mức: VD
\dangbt{Xác định biểu thức cơ năng trong dao động điều hòa}
\begin{ex}
% ID: L11C1B5-53-TLN
% Mức: VD
Con lắc lò xo có độ cứng $100\,\text{N}$ /m và biên độ $8\,\text{cm}$. Cơ năng theo đơn vị $J$ bằng bao nhiêu?
\shortans{0,32}
\end{ex}

% ===== Câu 211 =====
% ID: L11C1B5-53-TN
% Mức: NB
\dangbt{Tính động năng, thế năng và cơ năng tại một li độ hoặc thời điểm xác định}
\begin{ex}
% ID: L11C1B5-53-TN
% Mức: NB
Khi một vật dao động điều hòa đang chuyển động từ vị trí biên về vị trí cân bằng thì hiện tượng nào sau đây xảy ra?
\choice
{Động năng của vật giảm, thế năng tăng}
{\True Thế năng chuyển hóa dần thành động năng}
{Cơ năng của vật tăng liên tục}
{Động năng chuyển hóa dần thành thế năng}
\loigiai{
• Khi vật chuyển động từ biên về vị trí cân bằng, độ lớn li độ $x$ giảm dần về không, đồng thời tốc độ $v$ tăng dần lên giá trị cực đại.
 
• Do đó thế năng $W_{\text{t}} = \dfrac{1}{2}m\omega^2 x^2$ giảm dần và động năng $W_{\text{đ}} = \dfrac{1}{2}mv^2$ tăng dần.
 
• Theo định luật bảo toàn cơ năng, thế năng giảm đi bao nhiêu thì động năng tăng lên bấy nhiêu, nghĩa là thế năng chuyển hóa dần thành động năng.
}
\end{ex}

% ===== Câu 212 =====
% ID: L11C1B5-54-DS
% Mức: TH
\begin{ex}
% ID: L11C1B5-54-DS
% Mức: TH
Đồ thị thế năng theo li độ $W_t(x)$ của một vật dao động điều hòa là một đường parabol trên trục tọa độ Ox. Khảo sát hình học đồ thị này:
\choiceTF
{\True Điểm cực tiểu của đồ thị thế năng (đáy của parabol) nằm trùng khớp với gốc tọa độ $O(0, 0)$.}
{\True Khi li độ $x$ tăng từ $0$ đến vị trí biên dương $+A$, đồ thị biểu diễn thế năng là một đường parabol đi lên.}
{\True Giá trị lớn nhất trên đồ thị thế năng tương ứng với vị trí biên $x = \pm A$ và bằng cơ năng $W$.}
{Tại vị trí biên, động năng của vật đạt giá trị cực đại trên đồ thị.}
\loigiai{
\begin{itemchoice}
\item (Đúng) Điểm cực tiểu của đồ thị thế năng (đáy của parabol) nằm trùng khớp với gốc tọa độ $O(0, 0)$. (Vì): Mệnh đề này đúng vì khi vật ở vị trí cân bằng $x = 0$, thế năng đạt giá trị cực tiểu bằng 0, tương ứng điểm đáy của parabol nằm tại gốc tọa độ
\item (Đúng) Khi li độ $x$ tăng từ $0$ đến vị trí biên dương $+A$, đồ thị biểu diễn thế năng là một đường parabol đi lên. (Vì): Mệnh đề này đúng vì khi li độ tăng từ $0$ đến $A$, độ lớn li độ tăng nên thế năng tăng dần, đồ thị đi lên hướng tới biên dương
\item (Đúng) Giá trị lớn nhất trên đồ thị thế năng tương ứng với vị trí biên $x = \pm A$ và bằng cơ năng $W$. (Vì): Mệnh đề này đúng vì thế năng cực đại tại hai vị trí giới hạn biên có giá trị bằng cơ năng toàn phần
\item (Sai) Tại vị trí biên, động năng của vật đạt giá trị cực đại trên đồ thị. (Vì): Mệnh đề này sai vì tại vị trí biên vận tốc bằng không nên động năng bằng không (đạt giá trị cực tiểu chứ không phải cực đại
\end{itemchoice}
}
\end{ex}

% ===== Câu 213 =====
% ID: L11C1B5-54-TLN
% Mức: TH
\dangbt{Xác định và tính toán các đại lượng vận tốc, gia tốc tại thời điểm xác định}
\begin{ex}
% ID: L11C1B5-54-TLN
% Mức: TH
Vật có $A=9$ cm, $\omega=2$ rad/s. Tốc độ cực đại theo cm/s bằng bao nhiêu?
\shortans{18}
\end{ex}

% ===== Câu 214 =====
% ID: L11C1B5-54-TN
% Mức: NB
\dangbt{Tính động năng, thế năng và cơ năng tại một li độ hoặc thời điểm xác định}
\begin{ex}
% ID: L11C1B5-54-TN
% Mức: NB
Khi một vật dao động điều hòa đang chuyển động từ vị trí biên về vị trí cân bằng thì hiện tượng nào sau đây xảy ra?
\choice
{Động năng của vật giảm, thế năng tăng}
{\True Thế năng chuyển hóa dần thành động năng}
{Cơ năng của vật tăng liên tục}
{Động năng chuyển hóa dần thành thế năng}
\loigiai{
• Khi vật chuyển động từ biên về vị trí cân bằng, độ lớn li độ $x$ giảm dần về không, đồng thời tốc độ $v$ tăng dần lên giá trị cực đại.
 
• Do đó thế năng $W_{\text{t}} = \dfrac{1}{2}m\omega^2 x^2$ giảm dần và động năng $W_{\text{đ}} = \dfrac{1}{2}mv^2$ tăng dần.
 
• Theo định luật bảo toàn cơ năng, thế năng giảm đi bao nhiêu thì động năng tăng lên bấy nhiêu, nghĩa là thế năng chuyển hóa dần thành động năng.
}
\end{ex}

% ===== Câu 215 =====
% ID: L11C1B5-55-DS
% Mức: TH
\begin{ex}
% ID: L11C1B5-55-DS
% Mức: TH
Đồ thị thế năng theo li độ $W_t(x)$ của một vật dao động điều hòa là một đường parabol trên trục tọa độ Ox. Khảo sát hình học đồ thị này:
\choiceTF
{\True Điểm cực tiểu của đồ thị thế năng (đáy của parabol) nằm trùng khớp với gốc tọa độ $O(0, 0)$.}
{\True Khi li độ $x$ tăng từ $0$ đến vị trí biên dương $+A$, đồ thị biểu diễn thế năng là một đường parabol đi lên.}
{\True Giá trị lớn nhất trên đồ thị thế năng tương ứng với vị trí biên $x = \pm A$ và bằng cơ năng $W$.}
{Tại vị trí biên, động năng của vật đạt giá trị cực đại trên đồ thị.}
\loigiai{
\begin{itemchoice}
\item (Đúng) Điểm cực tiểu của đồ thị thế năng (đáy của parabol) nằm trùng khớp với gốc tọa độ $O(0, 0)$. (Vì): Mệnh đề này đúng vì khi vật ở vị trí cân bằng $x = 0$, thế năng đạt giá trị cực tiểu bằng 0, tương ứng điểm đáy của parabol nằm tại gốc tọa độ
\item (Đúng) Khi li độ $x$ tăng từ $0$ đến vị trí biên dương $+A$, đồ thị biểu diễn thế năng là một đường parabol đi lên. (Vì): Mệnh đề này đúng vì khi li độ tăng từ $0$ đến $A$, độ lớn li độ tăng nên thế năng tăng dần, đồ thị đi lên hướng tới biên dương
\item (Đúng) Giá trị lớn nhất trên đồ thị thế năng tương ứng với vị trí biên $x = \pm A$ và bằng cơ năng $W$. (Vì): Mệnh đề này đúng vì thế năng cực đại tại hai vị trí giới hạn biên có giá trị bằng cơ năng toàn phần
\item (Sai) Tại vị trí biên, động năng của vật đạt giá trị cực đại trên đồ thị. (Vì): Mệnh đề này sai vì tại vị trí biên vận tốc bằng không nên động năng bằng không (đạt giá trị cực tiểu chứ không phải cực đại
\end{itemchoice}
}
\end{ex}

% ===== Câu 216 =====
% ID: L11C1B5-55-TLN
% Mức: VD
\dangbt{Xác định và tính toán các đại lượng vận tốc, gia tốc tại thời điểm xác định}
\begin{ex}
% ID: L11C1B5-55-TLN
% Mức: VD
Vật có $A=10$ cm, $\omega=3$ rad/s. Tốc độ cực đại theo cm/s bằng bao nhiêu?
\shortans{30}
\end{ex}

% ===== Câu 217 =====
% ID: L11C1B5-55-TN
% Mức: NB
\dangbt{Tính động năng, thế năng và cơ năng tại một li độ hoặc thời điểm xác định}
\begin{ex}
% ID: L11C1B5-55-TN
% Mức: NB
Một con lắc lò xo gồm lò xo nhẹ có độ cứng $k$ và vật nhỏ dao động điều hòa với biên độ $A$. Cơ năng của con lắc được xác định bằng công thức nào sau đây?
\choice
{$W = \dfrac{1}{2}kA$}
{$W = kA^2$}
{\True $W = \dfrac{1}{2}kA^2$}
{$W = \dfrac{1}{2}mA^2$}
\loigiai{
• Trong con lắc lò xo, cơ năng bằng tổng động năng và thế năng đàn hồi của lò xo.
 
• Tại vị trí biên $x = \pm A$, thế năng đàn hồi đạt cực đại và bằng cơ năng của hệ:
 $W = W_{\text{tmax}} = \dfrac{1}{2}kx_{\max}^2.$
 
• Thay li độ cực đại bằng biên độ $A$:
 $W = \dfrac{1}{2}kA^2.$
}
\end{ex}

% ===== Câu 218 =====
% ID: L11C1B5-56-DS
% Mức: TH
\begin{ex}
% ID: L11C1B5-56-DS
% Mức: TH
Đồ thị thế năng theo li độ $W_t(x)$ của một vật dao động điều hòa là một đường parabol trên trục tọa độ Ox. Khảo sát hình học đồ thị này:
\choiceTF
{\True Điểm cực tiểu của đồ thị thế năng (đáy của parabol) nằm trùng khớp với gốc tọa độ $O(0, 0)$.}
{\True Khi li độ $x$ tăng từ $0$ đến vị trí biên dương $+A$, đồ thị biểu diễn thế năng là một đường parabol đi lên.}
{\True Giá trị lớn nhất trên đồ thị thế năng tương ứng với vị trí biên $x = \pm A$ và bằng cơ năng $W$.}
{Tại vị trí biên, động năng của vật đạt giá trị cực đại trên đồ thị.}
\loigiai{
\begin{itemchoice}
\item (Đúng) Điểm cực tiểu của đồ thị thế năng (đáy của parabol) nằm trùng khớp với gốc tọa độ $O(0, 0)$. (Vì): Mệnh đề này đúng vì khi vật ở vị trí cân bằng $x = 0$, thế năng đạt giá trị cực tiểu bằng 0, tương ứng điểm đáy của parabol nằm tại gốc tọa độ
\item (Đúng) Khi li độ $x$ tăng từ $0$ đến vị trí biên dương $+A$, đồ thị biểu diễn thế năng là một đường parabol đi lên. (Vì): Mệnh đề này đúng vì khi li độ tăng từ $0$ đến $A$, độ lớn li độ tăng nên thế năng tăng dần, đồ thị đi lên hướng tới biên dương
\item (Đúng) Giá trị lớn nhất trên đồ thị thế năng tương ứng với vị trí biên $x = \pm A$ và bằng cơ năng $W$. (Vì): Mệnh đề này đúng vì thế năng cực đại tại hai vị trí giới hạn biên có giá trị bằng cơ năng toàn phần
\item (Sai) Tại vị trí biên, động năng của vật đạt giá trị cực đại trên đồ thị. (Vì): Mệnh đề này sai vì tại vị trí biên vận tốc bằng không nên động năng bằng không (đạt giá trị cực tiểu chứ không phải cực đại
\end{itemchoice}
}
\end{ex}

% ===== Câu 219 =====
% ID: L11C1B5-56-TLN
% Mức: VD
\dangbt{Xác định và tính toán các đại lượng vận tốc, gia tốc tại thời điểm xác định}
\begin{ex}
% ID: L11C1B5-56-TLN
% Mức: VD
Vật có $A=5$ cm, $\omega=4$ rad/s. Tốc độ cực đại theo cm/s bằng bao nhiêu?
\shortans{20}
\end{ex}

% ===== Câu 220 =====
% ID: L11C1B5-56-TN
% Mức: NB
\dangbt{Tính động năng, thế năng và cơ năng tại một li độ hoặc thời điểm xác định}
\begin{ex}
% ID: L11C1B5-56-TN
% Mức: NB
Một con lắc lò xo gồm lò xo nhẹ có độ cứng $k$ và vật nhỏ dao động điều hòa với biên độ $A$. Cơ năng của con lắc được xác định bằng công thức nào sau đây?
\choice
{$W = \dfrac{1}{2}kA$}
{$W = kA^2$}
{\True $W = \dfrac{1}{2}kA^2$}
{$W = \dfrac{1}{2}mA^2$}
\loigiai{
• Trong con lắc lò xo, cơ năng bằng tổng động năng và thế năng đàn hồi của lò xo.
 
• Tại vị trí biên $x = \pm A$, thế năng đàn hồi đạt cực đại và bằng cơ năng của hệ:
 $W = W_{\text{tmax}} = \dfrac{1}{2}kx_{\max}^2.$
 
• Thay li độ cực đại bằng biên độ $A$:
 $W = \dfrac{1}{2}kA^2.$
}
\end{ex}

% ===== Câu 221 =====
% ID: L11C1B5-57-DS
% Mức: TH
\begin{ex}
% ID: L11C1B5-57-DS
% Mức: TH
Xét sự chuyển hóa và bảo toàn năng lượng trong dao động điều hòa của một con lắc lò xo đặt nằm ngang quanh vị trí cân bằng:
\choiceTF
{\True Khi vật di chuyển từ biên dương về vị trí cân bằng, lò xo bớt biến dạng nên thế năng đàn hồi giảm.}
{\True Khi vật di chuyển từ biên dương về vị trí cân bằng, lực đàn hồi sinh công dương làm vận tốc của vật tăng dần, động năng tăng.}
{\True Ở vị trí cân bằng, động năng của vật đạt cực đại, thế năng của vật đạt giá trị nhỏ nhất và bằng không (nếu chọn mốc thế năng tại vị trí cân bằng).}
{Tổng động năng và thế năng của hệ dao động điều hòa biến thiên tuần hoàn theo thời gian.}
\loigiai{
\begin{itemchoice}
\item (Đúng) Khi vật di chuyển từ biên dương về vị trí cân bằng, lò xo bớt biến dạng nên thế năng đàn hồi giảm. (Vì): Mệnh đề này đúng vì khi đi về vị trí cân bằng, độ dãn của lò xo giảm dần về không, kéo theo thế năng đàn hồi giảm dần
\item (Đúng) Khi vật di chuyển từ biên dương về vị trí cân bằng, lực đàn hồi sinh công dương làm vận tốc của vật tăng dần, động năng tăng. (Vì): Mệnh đề này đúng vì lực đàn hồi hướng về vị trí cân bằng, cùng chiều chuyển động của vật nên sinh công dương làm tăng động năng của vật
\item (Đúng) Ở vị trí cân bằng, động năng của vật đạt cực đại, thế năng của vật đạt giá trị nhỏ nhất và bằng không (nếu chọn mốc thế năng tại vị trí cân bằng). (Vì): Mệnh đề này đúng vì tại vị trí cân bằng thế năng đạt giá trị cực tiểu bằng 0 và động năng đạt giá trị cực đại
\item (Sai) Tổng động năng và thế năng của hệ dao động điều hòa biến thiên tuần hoàn theo thời gian. (Vì): Mệnh đề này sai vì tổng động năng và thế năng chính là cơ năng, đây là một đại lượng bảo toàn và không thay đổi theo thời gian
\end{itemchoice}
}
\end{ex}

% ===== Câu 222 =====
% ID: L11C1B5-57-TLN
% Mức: TH
\dangbt{Xác định vị trí dựa vào tỉ số động năng và thế năng}
\begin{ex}
% ID: L11C1B5-57-TLN
% Mức: TH
Con lắc lò xo có độ cứng $50\,\text{N}$ /m và biên độ $4\,\text{cm}$. Cơ năng theo đơn vị $J$ bằng bao nhiêu?
\shortans{0,04}
\end{ex}

% ===== Câu 223 =====
% ID: L11C1B5-57-TN
% Mức: NB
\dangbt{Tính động năng, thế năng và cơ năng tại một li độ hoặc thời điểm xác định}
\begin{ex}
% ID: L11C1B5-57-TN
% Mức: NB
Một con lắc lò xo gồm lò xo nhẹ có độ cứng $k$ và vật nhỏ dao động điều hòa với biên độ $A$. Cơ năng của con lắc được xác định bằng công thức nào sau đây?
\choice
{$W = \dfrac{1}{2}kA$}
{$W = kA^2$}
{\True $W = \dfrac{1}{2}kA^2$}
{$W = \dfrac{1}{2}mA^2$}
\loigiai{
• Trong con lắc lò xo, cơ năng bằng tổng động năng và thế năng đàn hồi của lò xo.
 
• Tại vị trí biên $x = \pm A$, thế năng đàn hồi đạt cực đại và bằng cơ năng của hệ:
 $W = W_{\text{tmax}} = \dfrac{1}{2}kx_{\max}^2.$
 
• Thay li độ cực đại bằng biên độ $A$:
 $W = \dfrac{1}{2}kA^2.$
}
\end{ex}

% ===== Câu 224 =====
% ID: L11C1B5-58-DS
% Mức: TH
\begin{ex}
% ID: L11C1B5-58-DS
% Mức: TH
Xét sự chuyển hóa và bảo toàn năng lượng trong dao động điều hòa của một con lắc lò xo đặt nằm ngang quanh vị trí cân bằng:
\choiceTF
{\True Khi vật di chuyển từ biên dương về vị trí cân bằng, lò xo bớt biến dạng nên thế năng đàn hồi giảm.}
{\True Khi vật di chuyển từ biên dương về vị trí cân bằng, lực đàn hồi sinh công dương làm vận tốc của vật tăng dần, động năng tăng.}
{\True Ở vị trí cân bằng, động năng của vật đạt cực đại, thế năng của vật đạt giá trị nhỏ nhất và bằng không (nếu chọn mốc thế năng tại vị trí cân bằng).}
{Tổng động năng và thế năng của hệ dao động điều hòa biến thiên tuần hoàn theo thời gian.}
\loigiai{
\begin{itemchoice}
\item (Đúng) Khi vật di chuyển từ biên dương về vị trí cân bằng, lò xo bớt biến dạng nên thế năng đàn hồi giảm. (Vì): Mệnh đề này đúng vì khi đi về vị trí cân bằng, độ dãn của lò xo giảm dần về không, kéo theo thế năng đàn hồi giảm dần
\item (Đúng) Khi vật di chuyển từ biên dương về vị trí cân bằng, lực đàn hồi sinh công dương làm vận tốc của vật tăng dần, động năng tăng. (Vì): Mệnh đề này đúng vì lực đàn hồi hướng về vị trí cân bằng, cùng chiều chuyển động của vật nên sinh công dương làm tăng động năng của vật
\item (Đúng) Ở vị trí cân bằng, động năng của vật đạt cực đại, thế năng của vật đạt giá trị nhỏ nhất và bằng không (nếu chọn mốc thế năng tại vị trí cân bằng). (Vì): Mệnh đề này đúng vì tại vị trí cân bằng thế năng đạt giá trị cực tiểu bằng 0 và động năng đạt giá trị cực đại
\item (Sai) Tổng động năng và thế năng của hệ dao động điều hòa biến thiên tuần hoàn theo thời gian. (Vì): Mệnh đề này sai vì tổng động năng và thế năng chính là cơ năng, đây là một đại lượng bảo toàn và không thay đổi theo thời gian
\end{itemchoice}
}
\end{ex}

% ===== Câu 225 =====
% ID: L11C1B5-58-TLN
% Mức: VD
\dangbt{Xác định vị trí dựa vào tỉ số động năng và thế năng}
\begin{ex}
% ID: L11C1B5-58-TLN
% Mức: VD
Con lắc lò xo có độ cứng $80\,\text{N}$ /m và biên độ $6\,\text{cm}$. Cơ năng theo đơn vị $J$ bằng bao nhiêu?
\shortans{0,144}
\end{ex}

% ===== Câu 226 =====
% ID: L11C1B5-58-TN
% Mức: NB
\dangbt{Tính động năng, thế năng và cơ năng tại một li độ hoặc thời điểm xác định}
\begin{ex}
% ID: L11C1B5-58-TN
% Mức: NB
Cơ năng của một chất điểm dao động điều hoà tỉ lệ thuận với
\choice
{\True bình phương biên độ dao động.}
{biên độ dao động.}
{bình phương tần số dao động.}
{tần số dao động.}
\loigiai{
Cơ năng của một chất điểm dao động điều hòa được xác định bởi công thức $E = \frac{1}{2}kA^2$ hoặc $E = \frac{1}{2}m\omega^2A^2$, trong đó $k$ là độ cứng của lò xo (hoặc hằng số tương đương), $m$ là khối lượng của vật, $\omega$ là tần số góc, và $A$ là biên độ dao động. Từ công thức này, ta thấy rằng cơ năng $E$ tỉ lệ thuận với bình phương biên độ dao động $A^2$. Các đại lượng $k$, $m$, $\omega$ là các hằng số đối với một hệ dao động điều hòa cụ thể.
}
\end{ex}

% ===== Câu 227 =====
% ID: L11C1B5-59-DS
% Mức: TH
\begin{ex}
% ID: L11C1B5-59-DS
% Mức: TH
Xét sự chuyển hóa và bảo toàn năng lượng trong dao động điều hòa của một con lắc lò xo đặt nằm ngang quanh vị trí cân bằng:
\choiceTF
{\True Khi vật di chuyển từ biên dương về vị trí cân bằng, lò xo bớt biến dạng nên thế năng đàn hồi giảm.}
{\True Khi vật di chuyển từ biên dương về vị trí cân bằng, lực đàn hồi sinh công dương làm vận tốc của vật tăng dần, động năng tăng.}
{\True Ở vị trí cân bằng, động năng của vật đạt cực đại, thế năng của vật đạt giá trị nhỏ nhất và bằng không (nếu chọn mốc thế năng tại vị trí cân bằng).}
{Tổng động năng và thế năng của hệ dao động điều hòa biến thiên tuần hoàn theo thời gian.}
\loigiai{
\begin{itemchoice}
\item (Đúng) Khi vật di chuyển từ biên dương về vị trí cân bằng, lò xo bớt biến dạng nên thế năng đàn hồi giảm. (Vì): Mệnh đề này đúng vì khi đi về vị trí cân bằng, độ dãn của lò xo giảm dần về không, kéo theo thế năng đàn hồi giảm dần
\item (Đúng) Khi vật di chuyển từ biên dương về vị trí cân bằng, lực đàn hồi sinh công dương làm vận tốc của vật tăng dần, động năng tăng. (Vì): Mệnh đề này đúng vì lực đàn hồi hướng về vị trí cân bằng, cùng chiều chuyển động của vật nên sinh công dương làm tăng động năng của vật
\item (Đúng) Ở vị trí cân bằng, động năng của vật đạt cực đại, thế năng của vật đạt giá trị nhỏ nhất và bằng không (nếu chọn mốc thế năng tại vị trí cân bằng). (Vì): Mệnh đề này đúng vì tại vị trí cân bằng thế năng đạt giá trị cực tiểu bằng 0 và động năng đạt giá trị cực đại
\item (Sai) Tổng động năng và thế năng của hệ dao động điều hòa biến thiên tuần hoàn theo thời gian. (Vì): Mệnh đề này sai vì tổng động năng và thế năng chính là cơ năng, đây là một đại lượng bảo toàn và không thay đổi theo thời gian
\end{itemchoice}
}
\end{ex}

% ===== Câu 228 =====
% ID: L11C1B5-59-TLN
% Mức: VD
\dangbt{Xác định vị trí dựa vào tỉ số động năng và thế năng}
\begin{ex}
% ID: L11C1B5-59-TLN
% Mức: VD
Con lắc lò xo có độ cứng $100\,\text{N}$ /m và biên độ $8\,\text{cm}$. Cơ năng theo đơn vị $J$ bằng bao nhiêu?
\shortans{0,32}
\end{ex}

% ===== Câu 229 =====
% ID: L11C1B5-59-TN
% Mức: NB
\dangbt{Tính động năng, thế năng và cơ năng tại một li độ hoặc thời điểm xác định}
\begin{ex}
% ID: L11C1B5-59-TN
% Mức: NB
Cơ năng của một chất điểm dao động điều hoà tỉ lệ thuận với
\choice
{\True bình phương biên độ dao động.}
{biên độ dao động.}
{bình phương tần số dao động.}
{tần số dao động.}
\loigiai{
Cơ năng của một chất điểm dao động điều hòa được xác định bởi công thức $E = \frac{1}{2}kA^2$ hoặc $E = \frac{1}{2}m\omega^2A^2$, trong đó $k$ là độ cứng của lò xo (hoặc hằng số tương đương), $m$ là khối lượng của vật, $\omega$ là tần số góc, và $A$ là biên độ dao động. Từ công thức này, ta thấy rằng cơ năng $E$ tỉ lệ thuận với bình phương biên độ dao động $A^2$. Các đại lượng $k$, $m$, $\omega$ là các hằng số đối với một hệ dao động điều hòa cụ thể.
}
\end{ex}

% ===== Câu 230 =====
% ID: L11C1B5-60-DS
% Mức: TH
\begin{ex}
% ID: L11C1B5-60-DS
% Mức: TH
Xét một vật dao động điều hòa có khối lượng $m$, biên độ $A$ và tần số góc $\omega$ không đổi. Đánh giá về mặt lí thuyết các giá trị năng lượng của vật:
\choiceTF
{Động năng của vật tại một thời điểm xác định có thể nhận giá trị âm nếu vật đang chuyển động ngược chiều dương của trục tọa độ.}
{\True Thế năng của vật luôn nhận giá trị dương hoặc bằng không khi ta chọn mốc thế năng tại vị trí cân bằng.}
{\True Cơ năng dao động của vật không thay đổi và bằng thế năng của vật tại vị trí biên.}
{\True Sự chuyển hóa qua lại giữa động năng và thế năng của vật xảy ra liên tục nhưng tổng của chúng tại mọi thời điểm luôn là hằng số.}
\loigiai{
\begin{itemchoice}
\item (Sai) Động năng của vật tại một thời điểm xác định có thể nhận giá trị âm nếu vật đang chuyển động ngược chiều dương của trục tọa độ. (Vì): Mệnh đề này sai vì động năng $W_d = \dfrac{1}{2}mv^2$ phụ thuộc vào bình phương vận tốc nên luôn nhận giá trị không âm, không phụ thuộc vào chiều chuyển động
\item (Đúng) Thế năng của vật luôn nhận giá trị dương hoặc bằng không khi ta chọn mốc thế năng tại vị trí cân bằng. (Vì): Mệnh đề này đúng vì thế năng $W_t = \dfrac{1}{2}m\omega^2 x^2$ luôn không âm và đạt giá trị nhỏ nhất bằng 0 tại vị trí cân bằng
\item (Đúng) Cơ năng dao động của vật không thay đổi và bằng thế năng của vật tại vị trí biên. (Vì): Mệnh đề này đúng vì tại biên động năng bằng không nên thế năng đạt cực đại và bằng đúng cơ năng toàn phần của hệ
\item (Đúng) Sự chuyển hóa qua lại giữa động năng và thế năng của vật xảy ra liên tục nhưng tổng của chúng tại mọi thời điểm luôn là hằng số.
\end{itemchoice}
}
\end{ex}

% ===== Câu 231 =====
% ID: L11C1B5-60-TLN
% Mức: VD
\dangbt{Xác định độ cứng lò xo từ chu kì biến thiên năng lượng}
\begin{ex}
% ID: L11C1B5-60-TLN
% Mức: VD
Động năng của con lắc lò xo có chu kì biến thiên $0,25\,\text{s}$; vật nặng có khối lượng $0,25\,\text{kg}$. Lấy π² = 10. Độ cứng của lò xo theo đơn vị N/m bằng bao nhiêu?
\shortans{40}
\loigiai{
Chu kì dao động $T= 2·0,25 =0,5\,\text{s}$. k = 4π²m/T² = 4·10·0,25/0,5² = $40\,\text{N}$ /m.
}
\end{ex}

% ===== Câu 232 =====
% ID: L11C1B5-60-TN
% Mức: NB
\dangbt{Tính động năng, thế năng và cơ năng tại một li độ hoặc thời điểm xác định}
\begin{ex}
% ID: L11C1B5-60-TN
% Mức: NB
Cơ năng của một chất điểm dao động điều hoà tỉ lệ thuận với
\choice
{\True bình phương biên độ dao động.}
{biên độ dao động.}
{bình phương tần số dao động.}
{tần số dao động.}
\loigiai{
Cơ năng của một chất điểm dao động điều hòa được xác định bởi công thức $E = \frac{1}{2}kA^2$ hoặc $E = \frac{1}{2}m\omega^2A^2$, trong đó $k$ là độ cứng của lò xo (hoặc hằng số tương đương), $m$ là khối lượng của vật, $\omega$ là tần số góc, và $A$ là biên độ dao động. Từ công thức này, ta thấy rằng cơ năng $E$ tỉ lệ thuận với bình phương biên độ dao động $A^2$. Các đại lượng $k$, $m$, $\omega$ là các hằng số đối với một hệ dao động điều hòa cụ thể.
}
\end{ex}

% ===== Câu 233 =====
% ID: L11C1B5-61-DS
% Mức: TH
\begin{ex}
% ID: L11C1B5-61-DS
% Mức: TH
Xét một vật dao động điều hòa có khối lượng $m$, biên độ $A$ và tần số góc $\omega$ không đổi. Đánh giá về mặt lí thuyết các giá trị năng lượng của vật:
\choiceTF
{Động năng của vật tại một thời điểm xác định có thể nhận giá trị âm nếu vật đang chuyển động ngược chiều dương của trục tọa độ.}
{\True Thế năng của vật luôn nhận giá trị dương hoặc bằng không khi ta chọn mốc thế năng tại vị trí cân bằng.}
{\True Cơ năng dao động của vật không thay đổi và bằng thế năng của vật tại vị trí biên.}
{\True Sự chuyển hóa qua lại giữa động năng và thế năng của vật xảy ra liên tục nhưng tổng của chúng tại mọi thời điểm luôn là hằng số.}
\loigiai{
\begin{itemchoice}
\item (Sai) Động năng của vật tại một thời điểm xác định có thể nhận giá trị âm nếu vật đang chuyển động ngược chiều dương của trục tọa độ. (Vì): Mệnh đề này sai vì động năng $W_d = \dfrac{1}{2}mv^2$ phụ thuộc vào bình phương vận tốc nên luôn nhận giá trị không âm, không phụ thuộc vào chiều chuyển động
\item (Đúng) Thế năng của vật luôn nhận giá trị dương hoặc bằng không khi ta chọn mốc thế năng tại vị trí cân bằng. (Vì): Mệnh đề này đúng vì thế năng $W_t = \dfrac{1}{2}m\omega^2 x^2$ luôn không âm và đạt giá trị nhỏ nhất bằng 0 tại vị trí cân bằng
\item (Đúng) Cơ năng dao động của vật không thay đổi và bằng thế năng của vật tại vị trí biên. (Vì): Mệnh đề này đúng vì tại biên động năng bằng không nên thế năng đạt cực đại và bằng đúng cơ năng toàn phần của hệ
\item (Đúng) Sự chuyển hóa qua lại giữa động năng và thế năng của vật xảy ra liên tục nhưng tổng của chúng tại mọi thời điểm luôn là hằng số.
\end{itemchoice}
}
\end{ex}

% ===== Câu 234 =====
% ID: L11C1B5-61-TLN
% Mức: VD
\dangbt{Xác định độ cứng lò xo từ chu kì biến thiên năng lượng}
\begin{ex}
% ID: L11C1B5-61-TLN
% Mức: VD
Động năng của con lắc lò xo có chu kì biến thiên $0,25\,\text{s}$; vật nặng có khối lượng $0,25\,\text{kg}$. Lấy π² = 10. Độ cứng của lò xo theo đơn vị N/m bằng bao nhiêu?
\shortans{40}
\loigiai{
Chu kì dao động $T= 2·0,25 =0,5\,\text{s}$. k = 4π²m/T² = 4·10·0,25/0,5² = $40\,\text{N}$ /m.
}
\end{ex}

% ===== Câu 235 =====
% ID: L11C1B5-61-TN
% Mức: NB
\dangbt{Tính động năng, thế năng và cơ năng tại một li độ hoặc thời điểm xác định}
\begin{ex}
% ID: L11C1B5-61-TN
% Mức: NB
Năng lượng của vật điều hoà:
\choice
{\True Bằng với thế năng của vật khi vật có li độ cực đại}
{Bằng với thế năng của vật khi vật qua vị trí cân bằng}
{Bằng với động năng của vật khi vật có li độ cực đại}
{Tỉ lệ với biên độ dao động}
\loigiai{
Năng lượng toàn phần của vật dao động điều hoà được bảo toàn và có giá trị $E = \frac{1}{2}kA^2$.
  Khi vật ở li độ cực đại ($x = \pm A$), vận tốc bằng 0 ($v=0$). Lúc này, động năng bằng 0 và thế năng đạt giá trị cực đại $U_{max} = \frac{1}{2}kA^2$. Do đó, năng lượng toàn phần của vật bằng thế năng của vật khi vật có li độ cực đại.
  Các phương án khác:
  - Năng lượng tỉ lệ với bình phương biên độ dao động ($A^2$), không phải tỉ lệ với biên độ ($A$).
  - Khi vật có li độ cực đại, vận tốc bằng 0 nên động năng bằng 0. Do đó năng lượng toàn phần không bằng động năng của vật khi vật có li độ cực đại.
  - Khi vật qua vị trí cân bằng ($x=0$), thế năng bằng 0. Do đó năng lượng toàn phần không bằng thế năng của vật khi vật qua vị trí cân bằng.
}
\end{ex}

% ===== Câu 236 =====
% ID: L11C1B5-62-DS
% Mức: TH
\begin{ex}
% ID: L11C1B5-62-DS
% Mức: TH
Xét một vật dao động điều hòa có khối lượng $m$, biên độ $A$ và tần số góc $\omega$ không đổi. Đánh giá về mặt lí thuyết các giá trị năng lượng của vật:
\choiceTF
{Động năng của vật tại một thời điểm xác định có thể nhận giá trị âm nếu vật đang chuyển động ngược chiều dương của trục tọa độ.}
{\True Thế năng của vật luôn nhận giá trị dương hoặc bằng không khi ta chọn mốc thế năng tại vị trí cân bằng.}
{\True Cơ năng dao động của vật không thay đổi và bằng thế năng của vật tại vị trí biên.}
{\True Sự chuyển hóa qua lại giữa động năng và thế năng của vật xảy ra liên tục nhưng tổng của chúng tại mọi thời điểm luôn là hằng số.}
\loigiai{
\begin{itemchoice}
\item (Sai) Động năng của vật tại một thời điểm xác định có thể nhận giá trị âm nếu vật đang chuyển động ngược chiều dương của trục tọa độ. (Vì): Mệnh đề này sai vì động năng $W_d = \dfrac{1}{2}mv^2$ phụ thuộc vào bình phương vận tốc nên luôn nhận giá trị không âm, không phụ thuộc vào chiều chuyển động
\item (Đúng) Thế năng của vật luôn nhận giá trị dương hoặc bằng không khi ta chọn mốc thế năng tại vị trí cân bằng. (Vì): Mệnh đề này đúng vì thế năng $W_t = \dfrac{1}{2}m\omega^2 x^2$ luôn không âm và đạt giá trị nhỏ nhất bằng 0 tại vị trí cân bằng
\item (Đúng) Cơ năng dao động của vật không thay đổi và bằng thế năng của vật tại vị trí biên. (Vì): Mệnh đề này đúng vì tại biên động năng bằng không nên thế năng đạt cực đại và bằng đúng cơ năng toàn phần của hệ
\item (Đúng) Sự chuyển hóa qua lại giữa động năng và thế năng của vật xảy ra liên tục nhưng tổng của chúng tại mọi thời điểm luôn là hằng số.
\end{itemchoice}
}
\end{ex}

% ===== Câu 237 =====
% ID: L11C1B5-62-TLN
% Mức: VD
\dangbt{Xác định độ cứng lò xo từ chu kì biến thiên năng lượng}
\begin{ex}
% ID: L11C1B5-62-TLN
% Mức: VD
Động năng của con lắc lò xo có chu kì biến thiên $0,25\,\text{s}$; vật nặng có khối lượng $0,25\,\text{kg}$. Lấy π² = 10. Độ cứng của lò xo theo đơn vị N/m bằng bao nhiêu?
\shortans{40}
\loigiai{
Chu kì dao động $T= 2·0,25 =0,5\,\text{s}$. k = 4π²m/T² = 4·10·0,25/0,5² = $40\,\text{N}$ /m.
}
\end{ex}

% ===== Câu 238 =====
% ID: L11C1B5-62-TN
% Mức: NB
\dangbt{Tính động năng, thế năng và cơ năng tại một li độ hoặc thời điểm xác định}
\begin{ex}
% ID: L11C1B5-62-TN
% Mức: NB
Năng lượng của vật điều hoà:
\choice
{\True Bằng với thế năng của vật khi vật có li độ cực đại}
{Bằng với thế năng của vật khi vật qua vị trí cân bằng}
{Bằng với động năng của vật khi vật có li độ cực đại}
{Tỉ lệ với biên độ dao động}
\loigiai{
Năng lượng toàn phần của vật dao động điều hoà được bảo toàn và có giá trị $E = \frac{1}{2}kA^2$.
  Khi vật ở li độ cực đại ($x = \pm A$), vận tốc bằng 0 ($v=0$). Lúc này, động năng bằng 0 và thế năng đạt giá trị cực đại $U_{max} = \frac{1}{2}kA^2$. Do đó, năng lượng toàn phần của vật bằng thế năng của vật khi vật có li độ cực đại.
  Các phương án khác:
  - Năng lượng tỉ lệ với bình phương biên độ dao động ($A^2$), không phải tỉ lệ với biên độ ($A$).
  - Khi vật có li độ cực đại, vận tốc bằng 0 nên động năng bằng 0. Do đó năng lượng toàn phần không bằng động năng của vật khi vật có li độ cực đại.
  - Khi vật qua vị trí cân bằng ($x=0$), thế năng bằng 0. Do đó năng lượng toàn phần không bằng thế năng của vật khi vật qua vị trí cân bằng.
}
\end{ex}

% ===== Câu 239 =====
% ID: L11C1B5-63-DS
% Mức: TH
\begin{ex}
% ID: L11C1B5-63-DS
% Mức: TH
Vật dao động theo phương trình $x=4\cos(2\pi t)~\mathrm{cm}$, $m=100\ g$. $\$ Tại thời điểm $t=0,25~\mathrm{s}$, vật đi qua vị trí cân bằng. Chọn Đúng hoặc Sai:
\choiceTF
{\True Tại thời điểm đó, động năng bằng cơ năng}
{\True Thế năng tại thời điểm đó bằng 0}
{Vận tốc tại thời điểm đó bằng 0}
{Tại thời điểm đó, cơ năng bằng 0}
\loigiai{
\begin{itemchoice}
\item (Đúng) Tại thời điểm đó, động năng bằng cơ năng. (Vì): Vị trí cân bằng: $x=0$, nên $W_t=0$, $W_d=W\Rightarrow \textbf{Đúng}
\item (Đúng) Thế năng tại thời điểm đó bằng 0. (Vì): Thế năng phụ thuộc vào$x^2$nên tại$x=0$thì$ W_t=0\Rightarrow \textbf{Đúng}
\item (Sai) Vận tốc tại thời điểm đó bằng 0. (Vì): Tại $x=0$, vật chuyển động nhậnh nhất $\Rightarrow v=v_{\max} \ne 0\Rightarrow \textbf{Sai}
\item (Sai) Tại thời điểm đó, cơ năng bằng 0. (Vì): Cơ năng trong dao động điều hòa lý tưởng luôn không đổi \Rightarrow\textbf{Sai}.$
\end{itemchoice}
}
\end{ex}

% ===== Câu 240 =====
% ID: L11C1B5-63-TLN
% Mức: VD
\dangbt{Xác định độ cứng lò xo từ chu kì biến thiên năng lượng}
\begin{ex}
% ID: L11C1B5-63-TLN
% Mức: VD
Một con lắc lò xo có khối lượng vật nhỏ là $50\,\text{g}$. Con lắc dao động điều hòa theo một trục cố định nằm ngang với phương trình $x=A\cos \left( \omega t \right)$. Cứ sau những khoảng thời gian $0,05\,\text{s}$ thì động năng và thế năng của vật lại bằng nhau. Lấy $\pi^2 = 10$. Lò xo của con lắc có độ cứng bằng
\shortans{50}
\loigiai{
Động năng bằng thế năng khi li độ $x = \pm \dfrac{A}{\sqrt{2}}$. Thời gian giữa hai lần liên tiếp động năng bằng thế năng là $\dfrac{T}{4}$.

Theo đề bài: $\dfrac{T}{4} = 0,05\,\text{s} \Rightarrow T = 0,2\,\text{s}$.

Tần số góc: $\omega = \dfrac{2\pi}{T} = \dfrac{2\pi}{0,2} = 10\pi\,\text{rad/s}$.

Độ cứng lò xo: $k = m\omega^2 = 0,05 \times (10\pi)^2 = 0,05 \times 100\pi^2 = 5\pi^2 = 5 \times 10 = 50\,\text{N/m}$.

Đáp số: $k = 50\,\text{N/m}$.
}
\end{ex}

% ===== Câu 241 =====
% ID: L11C1B5-63-TN
% Mức: NB
\dangbt{Tính động năng, thế năng và cơ năng tại một li độ hoặc thời điểm xác định}
\begin{ex}
% ID: L11C1B5-63-TN
% Mức: NB
Năng lượng của vật điều hoà:
\choice
{\True Bằng với thế năng của vật khi vật có li độ cực đại}
{Bằng với thế năng của vật khi vật qua vị trí cân bằng}
{Bằng với động năng của vật khi vật có li độ cực đại}
{Tỉ lệ với biên độ dao động}
\loigiai{
Năng lượng toàn phần của vật dao động điều hoà được bảo toàn và có giá trị $E = \frac{1}{2}kA^2$.
  Khi vật ở li độ cực đại ($x = \pm A$), vận tốc bằng 0 ($v=0$). Lúc này, động năng bằng 0 và thế năng đạt giá trị cực đại $U_{max} = \frac{1}{2}kA^2$. Do đó, năng lượng toàn phần của vật bằng thế năng của vật khi vật có li độ cực đại.
  Các phương án khác:
  - Năng lượng tỉ lệ với bình phương biên độ dao động ($A^2$), không phải tỉ lệ với biên độ ($A$).
  - Khi vật có li độ cực đại, vận tốc bằng 0 nên động năng bằng 0. Do đó năng lượng toàn phần không bằng động năng của vật khi vật có li độ cực đại.
  - Khi vật qua vị trí cân bằng ($x=0$), thế năng bằng 0. Do đó năng lượng toàn phần không bằng thế năng của vật khi vật qua vị trí cân bằng.
}
\end{ex}

% ===== Câu 242 =====
% ID: L11C1B5-64-DS
% Mức: TH
\begin{ex}
% ID: L11C1B5-64-DS
% Mức: TH
Vật dao động theo phương trình $x=4\cos(2\pi t)~\mathrm{cm}$, $m=100\ g$. $\$ Tại thời điểm $t=0,25~\mathrm{s}$, vật đi qua vị trí cân bằng. Chọn Đúng hoặc Sai:
\choiceTF
{\True Tại thời điểm đó, động năng bằng cơ năng}
{\True Thế năng tại thời điểm đó bằng 0}
{Vận tốc tại thời điểm đó bằng 0}
{Tại thời điểm đó, cơ năng bằng 0}
\loigiai{
\begin{itemchoice}
\item (Đúng) Tại thời điểm đó, động năng bằng cơ năng. (Vì): Vị trí cân bằng: $x=0$, nên $W_t=0$, $W_d=W\Rightarrow \textbf{Đúng}
\item (Đúng) Thế năng tại thời điểm đó bằng 0. (Vì): Thế năng phụ thuộc vào$x^2$nên tại$x=0$thì$ W_t=0\Rightarrow \textbf{Đúng}
\item (Sai) Vận tốc tại thời điểm đó bằng 0. (Vì): Tại $x=0$, vật chuyển động nhậnh nhất $\Rightarrow v=v_{\max} \ne 0\Rightarrow \textbf{Sai}
\item (Sai) Tại thời điểm đó, cơ năng bằng 0. (Vì): Cơ năng trong dao động điều hòa lý tưởng luôn không đổi \Rightarrow\textbf{Sai}.$
\end{itemchoice}
}
\end{ex}

% ===== Câu 243 =====
% ID: L11C1B5-64-TLN
% Mức: VD
\dangbt{Xác định độ cứng lò xo từ cơ năng và thông số dao động}
\begin{ex}
% ID: L11C1B5-64-TLN
% Mức: VD
Một con lắc lò xo dao động với biên độ $5\,\text{cm}$ và cơ năng $0,125\,\text{J}$. Độ cứng của lò xo theo đơn vị N/m bằng bao nhiêu?
\shortans{100}
\loigiai{
$k = 2W/A² = 2·0,125/(0,05)² = 100 N/m.$
}
\end{ex}

% ===== Câu 244 =====
% ID: L11C1B5-64-TN
% Mức: TH
\dangbt{Tính động năng, thế năng và cơ năng tại một li độ hoặc thời điểm xác định}
\begin{ex}
% ID: L11C1B5-64-TN
% Mức: TH
Một vật có thế năng 50 $J$ ở độ cao $5\,\text{m}$. Khối lượng của vật là bao nhiêu? (lấy $g = 10 \text{ m/s}^2$)
\choice
{$1\,\text{kg}$}
{\True $2\,\text{kg}$}
{$3\,\text{kg}$}
{$4\,\text{kg}$}
\loigiai{
Thế năng được tính bằng công thức $W_t = mgh$. Thay $W_t = 50J$, $h = 5$ m, và $g = 10 \text{ m/s}^2$, ta có $m = \frac{W_t}{gh} = \frac{50}{10 \times 5} = 2$ kg.
}
\end{ex}

% ===== Câu 245 =====
% ID: L11C1B5-65-DS
% Mức: TH
\begin{ex}
% ID: L11C1B5-65-DS
% Mức: TH
Vật dao động theo phương trình $x=4\cos(2\pi t)~\mathrm{cm}$, $m=100\ g$. $\$ Tại thời điểm $t=0,25~\mathrm{s}$, vật đi qua vị trí cân bằng. Chọn Đúng hoặc Sai:
\choiceTF
{\True Tại thời điểm đó, động năng bằng cơ năng}
{\True Thế năng tại thời điểm đó bằng 0}
{Vận tốc tại thời điểm đó bằng 0}
{Tại thời điểm đó, cơ năng bằng 0}
\loigiai{
\begin{itemchoice}
\item (Đúng) Tại thời điểm đó, động năng bằng cơ năng. (Vì): Vị trí cân bằng: $x=0$, nên $W_t=0$, $W_d=W\Rightarrow \textbf{Đúng}
\item (Đúng) Thế năng tại thời điểm đó bằng 0. (Vì): Thế năng phụ thuộc vào$x^2$nên tại$x=0$thì$ W_t=0\Rightarrow \textbf{Đúng}
\item (Sai) Vận tốc tại thời điểm đó bằng 0. (Vì): Tại $x=0$, vật chuyển động nhậnh nhất $\Rightarrow v=v_{\max} \ne 0\Rightarrow \textbf{Sai}
\item (Sai) Tại thời điểm đó, cơ năng bằng 0. (Vì): Cơ năng trong dao động điều hòa lý tưởng luôn không đổi \Rightarrow\textbf{Sai}.$
\end{itemchoice}
}
\end{ex}

% ===== Câu 246 =====
% ID: L11C1B5-65-TLN
% Mức: VD
\dangbt{Xác định độ cứng lò xo từ cơ năng và thông số dao động}
\begin{ex}
% ID: L11C1B5-65-TLN
% Mức: VD
Một con lắc lò xo dao động với biên độ $5\,\text{cm}$ và cơ năng $0,125\,\text{J}$. Độ cứng của lò xo theo đơn vị N/m bằng bao nhiêu?
\shortans{100}
\loigiai{
$k = 2W/A² = 2·0,125/(0,05)² = 100 N/m.$
}
\end{ex}

% ===== Câu 247 =====
% ID: L11C1B5-65-TN
% Mức: TH
\dangbt{Tính động năng, thế năng và cơ năng tại một li độ hoặc thời điểm xác định}
\begin{ex}
% ID: L11C1B5-65-TN
% Mức: TH
Điều nào sau đây là đúng khi nói về động năng và thế năng của một vật khối lượng không đổi dao động điều hòa
\choice
{Thế năng tăng khi li độ của vật tăng}
{Trong một chu kỳ luôn có 2 thời điểm mà ở đó động bằng thế năng}
{\True Trong một chu kì luôn có 4 thời điểm mà ở đó động năng bằng 3 thế năng}
{Động năng của một vật tăng chỉ khi vận Gtốc của vật tăng}
\loigiai{
A. Trong một chu kì, động năng bằng 3 thế năng ($K=3U$) khi $E-U=3U \Rightarrow E=4U$. Mà tổng năng lượng $E = \frac{1}{2}kA^2$ và thế năng $U = \frac{1}{2}kx^2$, nên $\frac{1}{2}kA^2 = 4 \cdot \frac{1}{2}kx^2 \Rightarrow A^2 = 4x^2 \Rightarrow x = \pm \frac{A}{2}$. Trong một chu kì, vật đi qua các vị trí $x = \frac{A}{2}$ và $x = -\frac{A}{2}$ tổng cộng 4 lần (mỗi vị trí 2 lần khi đi qua và đi về). Vậy khẳng định $A$ đúng.B. Thế năng $U = \frac{1}{2}kx^2$. Thế năng tăng khi độ lớn của li độ $|x|$ tăng. Nếu li độ $x$ tăng từ một giá trị âm sang dương (ví dụ từ $-A$ đến $0$), thế năng giảm. Vì vậy khẳng định $B$ sai.C. Động năng bằng thế năng ($K=U$) khi $E=2U$. Tương tự, $\frac{1}{2}kA^2 = 2 \cdot \frac{1}{2}kx^2 \Rightarrow A^2 = 2x^2 \Rightarrow x = \pm \frac{A}{\sqrt{2}}$. Trong một chu kì, vật đi qua các vị trí $x = \frac{A}{\sqrt{2}}$ và $x = -\frac{A}{\sqrt{2}}$ tổng cộng 4 lần. Vậy khẳng định $C$ sai.D. Động năng $K = \frac{1}{2}mv^2$. Động năng của một vật tăng khi độ lớn của vận tốc $|v|$ tăng. Nếu vận tốc $v$ tăng từ giá trị âm nhỏ (ví dụ từ $-5$ m/s lên $-2$ m/s), thì độ lớn của vận tốc giảm, và động năng cũng giảm. Vì vậy khẳng định $D$ sai.
}
\end{ex}

% ===== Câu 248 =====
% ID: L11C1B5-66-DS
% Mức: VD
\begin{ex}
% ID: L11C1B5-66-DS
% Mức: VD
Một vật có khối lượng $m = 200$ g thực hiện dao động điều hòa với tần số góc $\omega = 10$ rad/s và biên độ dao động $A = 10$ cm. Chọn mốc thế năng tại vị trí cân bằng của vật.
\choiceTF
{\True Cơ năng toàn phần của vật dao động bằng $0,1$ J.}
{\True Tại vị trí vật có li độ $x = 6$ cm, thế năng của vật nhận giá trị bằng $0,036$ J.}
{\True Tại vị trí vật có li độ $x = 6$ cm, động năng của vật nhận giá trị bằng $0,064$ J.}
{\True Khi vật đi qua vị trí cân bằng, tốc độ của vật đạt giá trị cực đại là $1$ m/s.}
\loigiai{
\begin{itemchoice}
\item (Đúng) Cơ năng toàn phần của vật dao động bằng $0,1$ J. (Vì): Mệnh đề này đúng vì ta đổi $A = 0,1$ m, $m = 0,2$ kg; cơ năng là $W = \dfrac{1}{2}m\omega^2 A^2 = 0,5 \cdot 0,2 \cdot 10^2 \cdot 0,1^2 = 0,1$ J
\item (Đúng) Tại vị trí vật có li độ $x = 6$ cm, thế năng của vật nhận giá trị bằng $0,036$ J. (Vì): Mệnh đề này đúng vì đổi $x = 0,06$ m; thế năng là $W_t = \dfrac{1}{2}m\omega^2 x^2 = 0,5 \cdot 0,2 \cdot 10^2 \cdot 0,06^2 = 0,036$ J
\item (Đúng) Tại vị trí vật có li độ $x = 6$ cm, động năng của vật nhận giá trị bằng $0,064$ J. (Vì): Mệnh đề này đúng vì động năng của vật là $W_d = W - W_t = 0,1 - 0,036 = 0,064$ J
\item (Đúng) Khi vật đi qua vị trí cân bằng, tốc độ của vật đạt giá trị cực đại là $1$ m/s. (Vì): Mệnh đề này đúng vì tốc độ cực đại khi qua vị trí cân bằng là $v_{\text{max}} = \omega A = 10 \cdot 0,1 = 1$ m/s
\end{itemchoice}
}
\end{ex}

% ===== Câu 249 =====
% ID: L11C1B5-66-TLN
% Mức: VD
\dangbt{Xác định độ cứng lò xo từ cơ năng và thông số dao động}
\begin{ex}
% ID: L11C1B5-66-TLN
% Mức: VD
Một con lắc lò xo dao động với biên độ $5\,\text{cm}$ và cơ năng $0,125\,\text{J}$. Độ cứng của lò xo theo đơn vị N/m bằng bao nhiêu?
\shortans{100}
\loigiai{
$k = 2W/A² = 2·0,125/(0,05)² = 100 N/m.$
}
\end{ex}

% ===== Câu 250 =====
% ID: L11C1B5-66-TN
% Mức: TH
\dangbt{Tính động năng, thế năng và cơ năng tại một li độ hoặc thời điểm xác định}
\begin{ex}
% ID: L11C1B5-66-TN
% Mức: TH
Điều nào sau đây là đúng khi nói về động năng và thế năng của một vật khối lượng không đổi dao động điều hòa
\choice
{Thế năng tăng khi li độ của vật tăng}
{Trong một chu kỳ luôn có 2 thời điểm mà ở đó động bằng thế năng}
{\True Trong một chu kì luôn có 4 thời điểm mà ở đó động năng bằng 3 thế năng}
{Động năng của một vật tăng chỉ khi vận Gtốc của vật tăng}
\loigiai{
A. Trong một chu kì, động năng bằng 3 thế năng ($K=3U$) khi $E-U=3U \Rightarrow E=4U$. Mà tổng năng lượng $E = \frac{1}{2}kA^2$ và thế năng $U = \frac{1}{2}kx^2$, nên $\frac{1}{2}kA^2 = 4 \cdot \frac{1}{2}kx^2 \Rightarrow A^2 = 4x^2 \Rightarrow x = \pm \frac{A}{2}$. Trong một chu kì, vật đi qua các vị trí $x = \frac{A}{2}$ và $x = -\frac{A}{2}$ tổng cộng 4 lần (mỗi vị trí 2 lần khi đi qua và đi về). Vậy khẳng định $A$ đúng.B. Thế năng $U = \frac{1}{2}kx^2$. Thế năng tăng khi độ lớn của li độ $|x|$ tăng. Nếu li độ $x$ tăng từ một giá trị âm sang dương (ví dụ từ $-A$ đến $0$), thế năng giảm. Vì vậy khẳng định $B$ sai.C. Động năng bằng thế năng ($K=U$) khi $E=2U$. Tương tự, $\frac{1}{2}kA^2 = 2 \cdot \frac{1}{2}kx^2 \Rightarrow A^2 = 2x^2 \Rightarrow x = \pm \frac{A}{\sqrt{2}}$. Trong một chu kì, vật đi qua các vị trí $x = \frac{A}{\sqrt{2}}$ và $x = -\frac{A}{\sqrt{2}}$ tổng cộng 4 lần. Vậy khẳng định $C$ sai.D. Động năng $K = \frac{1}{2}mv^2$. Động năng của một vật tăng khi độ lớn của vận tốc $|v|$ tăng. Nếu vận tốc $v$ tăng từ giá trị âm nhỏ (ví dụ từ $-5$ m/s lên $-2$ m/s), thì độ lớn của vận tốc giảm, và động năng cũng giảm. Vì vậy khẳng định $D$ sai.
}
\end{ex}

% ===== Câu 251 =====
% ID: L11C1B5-67-DS
% Mức: VD
\begin{ex}
% ID: L11C1B5-67-DS
% Mức: VD
Một vật có khối lượng $m = 200$ g thực hiện dao động điều hòa với tần số góc $\omega = 10$ rad/s và biên độ dao động $A = 10$ cm. Chọn mốc thế năng tại vị trí cân bằng của vật.
\choiceTF
{\True Cơ năng toàn phần của vật dao động bằng $0,1$ J.}
{\True Tại vị trí vật có li độ $x = 6$ cm, thế năng của vật nhận giá trị bằng $0,036$ J.}
{\True Tại vị trí vật có li độ $x = 6$ cm, động năng của vật nhận giá trị bằng $0,064$ J.}
{\True Khi vật đi qua vị trí cân bằng, tốc độ của vật đạt giá trị cực đại là $1$ m/s.}
\loigiai{
\begin{itemchoice}
\item (Đúng) Cơ năng toàn phần của vật dao động bằng $0,1$ J. (Vì): Mệnh đề này đúng vì ta đổi $A = 0,1$ m, $m = 0,2$ kg; cơ năng là $W = \dfrac{1}{2}m\omega^2 A^2 = 0,5 \cdot 0,2 \cdot 10^2 \cdot 0,1^2 = 0,1$ J
\item (Đúng) Tại vị trí vật có li độ $x = 6$ cm, thế năng của vật nhận giá trị bằng $0,036$ J. (Vì): Mệnh đề này đúng vì đổi $x = 0,06$ m; thế năng là $W_t = \dfrac{1}{2}m\omega^2 x^2 = 0,5 \cdot 0,2 \cdot 10^2 \cdot 0,06^2 = 0,036$ J
\item (Đúng) Tại vị trí vật có li độ $x = 6$ cm, động năng của vật nhận giá trị bằng $0,064$ J. (Vì): Mệnh đề này đúng vì động năng của vật là $W_d = W - W_t = 0,1 - 0,036 = 0,064$ J
\item (Đúng) Khi vật đi qua vị trí cân bằng, tốc độ của vật đạt giá trị cực đại là $1$ m/s. (Vì): Mệnh đề này đúng vì tốc độ cực đại khi qua vị trí cân bằng là $v_{\text{max}} = \omega A = 10 \cdot 0,1 = 1$ m/s
\end{itemchoice}
}
\end{ex}

% ===== Câu 252 =====
% ID: L11C1B5-67-TLN
% Mức: VD
\dangbt{Xác định độ cứng lò xo từ cơ năng và thông số dao động}
\begin{ex}
% ID: L11C1B5-67-TLN
% Mức: VD
Một con lắc lò xo có vật nặng $0,2\,\text{kg}$, tốc độ cực đại $1\,\text{m}$ /s và biên độ $10\,\text{cm}$. Độ cứng của lò xo theo đơn vị N/m bằng bao nhiêu?
\shortans{20}
\loigiai{
$W$ = mvmax²/2 = kA²/2 nên k = mvmax²/A² = 0,2·1²/0,1² = $20\,\text{N}$ /m.
}
\end{ex}

% ===== Câu 253 =====
% ID: L11C1B5-67-TN
% Mức: TH
\dangbt{Tính động năng, thế năng và cơ năng tại một li độ hoặc thời điểm xác định}
\begin{ex}
% ID: L11C1B5-67-TN
% Mức: TH
Điều nào sau đây là đúng khi nói về động năng và thế năng của một vật khối lượng không đổi dao động điều hòa
\choice
{Thế năng tăng khi li độ của vật tăng}
{Trong một chu kỳ luôn có 2 thời điểm mà ở đó động bằng thế năng}
{\True Trong một chu kì luôn có 4 thời điểm mà ở đó động năng bằng 3 thế năng}
{Động năng của một vật tăng chỉ khi vận Gtốc của vật tăng}
\loigiai{
A. Trong một chu kì, động năng bằng 3 thế năng ($K=3U$) khi $E-U=3U \Rightarrow E=4U$. Mà tổng năng lượng $E = \frac{1}{2}kA^2$ và thế năng $U = \frac{1}{2}kx^2$, nên $\frac{1}{2}kA^2 = 4 \cdot \frac{1}{2}kx^2 \Rightarrow A^2 = 4x^2 \Rightarrow x = \pm \frac{A}{2}$. Trong một chu kì, vật đi qua các vị trí $x = \frac{A}{2}$ và $x = -\frac{A}{2}$ tổng cộng 4 lần (mỗi vị trí 2 lần khi đi qua và đi về). Vậy khẳng định $A$ đúng.B. Thế năng $U = \frac{1}{2}kx^2$. Thế năng tăng khi độ lớn của li độ $|x|$ tăng. Nếu li độ $x$ tăng từ một giá trị âm sang dương (ví dụ từ $-A$ đến $0$), thế năng giảm. Vì vậy khẳng định $B$ sai.C. Động năng bằng thế năng ($K=U$) khi $E=2U$. Tương tự, $\frac{1}{2}kA^2 = 2 \cdot \frac{1}{2}kx^2 \Rightarrow A^2 = 2x^2 \Rightarrow x = \pm \frac{A}{\sqrt{2}}$. Trong một chu kì, vật đi qua các vị trí $x = \frac{A}{\sqrt{2}}$ và $x = -\frac{A}{\sqrt{2}}$ tổng cộng 4 lần. Vậy khẳng định $C$ sai.D. Động năng $K = \frac{1}{2}mv^2$. Động năng của một vật tăng khi độ lớn của vận tốc $|v|$ tăng. Nếu vận tốc $v$ tăng từ giá trị âm nhỏ (ví dụ từ $-5$ m/s lên $-2$ m/s), thì độ lớn của vận tốc giảm, và động năng cũng giảm. Vì vậy khẳng định $D$ sai.
}
\end{ex}

% ===== Câu 254 =====
% ID: L11C1B5-68-DS
% Mức: VD
\begin{ex}
% ID: L11C1B5-68-DS
% Mức: VD
Một vật có khối lượng $m = 200$ g thực hiện dao động điều hòa với tần số góc $\omega = 10$ rad/s và biên độ dao động $A = 10$ cm. Chọn mốc thế năng tại vị trí cân bằng của vật.
\choiceTF
{\True Cơ năng toàn phần của vật dao động bằng $0,1$ J.}
{\True Tại vị trí vật có li độ $x = 6$ cm, thế năng của vật nhận giá trị bằng $0,036$ J.}
{\True Tại vị trí vật có li độ $x = 6$ cm, động năng của vật nhận giá trị bằng $0,064$ J.}
{\True Khi vật đi qua vị trí cân bằng, tốc độ của vật đạt giá trị cực đại là $1$ m/s.}
\loigiai{
\begin{itemchoice}
\item (Đúng) Cơ năng toàn phần của vật dao động bằng $0,1$ J. (Vì): Mệnh đề này đúng vì ta đổi $A = 0,1$ m, $m = 0,2$ kg; cơ năng là $W = \dfrac{1}{2}m\omega^2 A^2 = 0,5 \cdot 0,2 \cdot 10^2 \cdot 0,1^2 = 0,1$ J
\item (Đúng) Tại vị trí vật có li độ $x = 6$ cm, thế năng của vật nhận giá trị bằng $0,036$ J. (Vì): Mệnh đề này đúng vì đổi $x = 0,06$ m; thế năng là $W_t = \dfrac{1}{2}m\omega^2 x^2 = 0,5 \cdot 0,2 \cdot 10^2 \cdot 0,06^2 = 0,036$ J
\item (Đúng) Tại vị trí vật có li độ $x = 6$ cm, động năng của vật nhận giá trị bằng $0,064$ J. (Vì): Mệnh đề này đúng vì động năng của vật là $W_d = W - W_t = 0,1 - 0,036 = 0,064$ J
\item (Đúng) Khi vật đi qua vị trí cân bằng, tốc độ của vật đạt giá trị cực đại là $1$ m/s. (Vì): Mệnh đề này đúng vì tốc độ cực đại khi qua vị trí cân bằng là $v_{\text{max}} = \omega A = 10 \cdot 0,1 = 1$ m/s
\end{itemchoice}
}
\end{ex}

% ===== Câu 255 =====
% ID: L11C1B5-68-TLN
% Mức: VD
\dangbt{Xác định độ cứng lò xo từ cơ năng và thông số dao động}
\begin{ex}
% ID: L11C1B5-68-TLN
% Mức: VD
Một con lắc lò xo có vật nặng $0,2\,\text{kg}$, tốc độ cực đại $1\,\text{m}$ /s và biên độ $10\,\text{cm}$. Độ cứng của lò xo theo đơn vị N/m bằng bao nhiêu?
\shortans{20}
\loigiai{
$W$ = mvmax²/2 = kA²/2 nên k = mvmax²/A² = 0,2·1²/0,1² = $20\,\text{N}$ /m.
}
\end{ex}

% ===== Câu 256 =====
% ID: L11C1B5-68-TN
% Mức: TH
\dangbt{Tính động năng, thế năng và cơ năng tại một li độ hoặc thời điểm xác định}
\begin{ex}
% ID: L11C1B5-68-TN
% Mức: TH
Một vật dao động điều hòa theo một trục cố định (mốc thế năng ở vị trí cân bằng) thì
\choice
{động năng của vật cực đại khi gia tốc của vật có độ lớn cực đại}
{\True thế năng của vật cực đại khi vật ở vị trí biên}
{khi ở vị trí cân bằng, thế năng của vật bằng cơ năng}
{khi vật đi từ vị trí cân bằng ra biên, vận tốc và gia tốc của vật luôn cùng dấu}
\loigiai{
Thế năng của vật dao động điều hòa được tính bằng công thức $W_t = \frac{1}{2}kx^2$. Tại vị trí biên, li độ $x$ đạt giá trị cực đại ($x = \pm A$), do đó thế năng đạt giá trị cực đại. Các phương án còn lại là sai:
  - Khi vật đi từ vị trí cân bằng ra biên, vận tốc và gia tốc của vật luôn ngược dấu (ví dụ, khi vật đi từ VTCB ra biên dương, $v {\gt } 0$ và $a {\lt } 0$).
  - Khi ở vị trí cân bằng ($x=0$), thế năng của vật bằng 0, động năng của vật đạt cực đại và bằng cơ năng.
  - Động năng của vật cực đại khi vật ở vị trí cân bằng ($x=0$), tại đó gia tốc của vật bằng 0 (độ lớn cực tiểu), không phải cực đại. Gia tốc của vật có độ lớn cực đại tại vị trí biên, nơi động năng bằng 0.
}
\end{ex}

% ===== Câu 257 =====
% ID: L11C1B5-69-DS
% Mức: VD
\begin{ex}
% ID: L11C1B5-69-DS
% Mức: VD
Một con lắc lò xo có độ cứng của lò xo $k = 100$ N/m, khối lượng vật nặng $m = 1$ kg, dao động điều hòa với biên độ $A = 10$ cm dọc theo trục Ox nằm ngang. Chọn mốc thế năng tại vị trí cân bằng.
\choiceTF
{\True Cơ năng dao động của con lắc lò xo bằng $0,5$ J.}
{\True Khi vật nặng có li độ $x = 6$ cm thì thế năng đàn hồi của lò xo bằng $0,18$ J.}
{\True Khi vật nặng có li độ $x = 6$ cm thì động năng của vật nặng bằng $0,32$ J.}
{Khi vật nặng có li độ $x = 6$ cm thì tốc độ chuyển động của vật nặng bằng $80$ m/s.}
\loigiai{
\begin{itemchoice}
\item (Đúng) Cơ năng dao động của con lắc lò xo bằng $0,5$ J. (Vì): Mệnh đề này đúng vì cơ năng của hệ là $W = \dfrac{1}{2}kA^2 = 0,5 \cdot 100 \cdot 0,1^2 = 0,5$ J
\item (Đúng) Khi vật nặng có li độ $x = 6$ cm thì thế năng đàn hồi của lò xo bằng $0,18$ J. (Vì): Mệnh đề này đúng vì đổi $x = 0,06$ m; thế năng đàn hồi là $W_t = \dfrac{1}{2}kx^2 = 0,5 \cdot 100 \cdot 0,06^2 = 0,18$ J
\item (Đúng) Khi vật nặng có li độ $x = 6$ cm thì động năng của vật nặng bằng $0,32$ J. (Vì): Mệnh đề này đúng vì động năng của vật là $W_d = W - W_t = 0,5 - 0,18 = 0,32$ J
\item (Sai) Khi vật nặng có li độ $x = 6$ cm thì tốc độ chuyển động của vật nặng bằng $80$ m/s. (Vì): Mệnh đề này sai vì động năng $W_d = \dfrac{1}{2}mv^2 \Rightarrow 0,32 = 0,5 \cdot 1 \cdot v^2 \Rightarrow v^2 = 0,64 \Rightarrow v = 0,8$ m/s
\end{itemchoice}
}
\end{ex}

% ===== Câu 258 =====
% ID: L11C1B5-69-TLN
% Mức: VD
\dangbt{Xác định độ cứng lò xo từ cơ năng và thông số dao động}
\begin{ex}
% ID: L11C1B5-69-TLN
% Mức: VD
Một con lắc lò xo có vật nặng $0,2\,\text{kg}$, tốc độ cực đại $1\,\text{m}$ /s và biên độ $10\,\text{cm}$. Độ cứng của lò xo theo đơn vị N/m bằng bao nhiêu?
\shortans{20}
\loigiai{
$W$ = mvmax²/2 = kA²/2 nên k = mvmax²/A² = 0,2·1²/0,1² = $20\,\text{N}$ /m.
}
\end{ex}

% ===== Câu 259 =====
% ID: L11C1B5-69-TN
% Mức: TH
\dangbt{Tính động năng, thế năng và cơ năng tại một li độ hoặc thời điểm xác định}
\begin{ex}
% ID: L11C1B5-69-TN
% Mức: TH
Một vật dao động điều hòa theo một trục cố định (mốc thế năng ở vị trí cân bằng) thì
\choice
{động năng của vật cực đại khi gia tốc của vật có độ lớn cực đại}
{\True thế năng của vật cực đại khi vật ở vị trí biên}
{khi ở vị trí cân bằng, thế năng của vật bằng cơ năng}
{khi vật đi từ vị trí cân bằng ra biên, vận tốc và gia tốc của vật luôn cùng dấu}
\loigiai{
Thế năng của vật dao động điều hòa được tính bằng công thức $W_t = \frac{1}{2}kx^2$. Tại vị trí biên, li độ $x$ đạt giá trị cực đại ($x = \pm A$), do đó thế năng đạt giá trị cực đại. Các phương án còn lại là sai:
  - Khi vật đi từ vị trí cân bằng ra biên, vận tốc và gia tốc của vật luôn ngược dấu (ví dụ, khi vật đi từ VTCB ra biên dương, $v {\gt } 0$ và $a {\lt } 0$).
  - Khi ở vị trí cân bằng ($x=0$), thế năng của vật bằng 0, động năng của vật đạt cực đại và bằng cơ năng.
  - Động năng của vật cực đại khi vật ở vị trí cân bằng ($x=0$), tại đó gia tốc của vật bằng 0 (độ lớn cực tiểu), không phải cực đại. Gia tốc của vật có độ lớn cực đại tại vị trí biên, nơi động năng bằng 0.
}
\end{ex}

% ===== Câu 260 =====
% ID: L11C1B5-70-DS
% Mức: VD
\begin{ex}
% ID: L11C1B5-70-DS
% Mức: VD
Một con lắc lò xo có độ cứng của lò xo $k = 100$ N/m, khối lượng vật nặng $m = 1$ kg, dao động điều hòa với biên độ $A = 10$ cm dọc theo trục Ox nằm ngang. Chọn mốc thế năng tại vị trí cân bằng.
\choiceTF
{\True Cơ năng dao động của con lắc lò xo bằng $0,5$ J.}
{\True Khi vật nặng có li độ $x = 6$ cm thì thế năng đàn hồi của lò xo bằng $0,18$ J.}
{\True Khi vật nặng có li độ $x = 6$ cm thì động năng của vật nặng bằng $0,32$ J.}
{Khi vật nặng có li độ $x = 6$ cm thì tốc độ chuyển động của vật nặng bằng $80$ m/s.}
\loigiai{
\begin{itemchoice}
\item (Đúng) Cơ năng dao động của con lắc lò xo bằng $0,5$ J. (Vì): Mệnh đề này đúng vì cơ năng của hệ là $W = \dfrac{1}{2}kA^2 = 0,5 \cdot 100 \cdot 0,1^2 = 0,5$ J
\item (Đúng) Khi vật nặng có li độ $x = 6$ cm thì thế năng đàn hồi của lò xo bằng $0,18$ J. (Vì): Mệnh đề này đúng vì đổi $x = 0,06$ m; thế năng đàn hồi là $W_t = \dfrac{1}{2}kx^2 = 0,5 \cdot 100 \cdot 0,06^2 = 0,18$ J
\item (Đúng) Khi vật nặng có li độ $x = 6$ cm thì động năng của vật nặng bằng $0,32$ J. (Vì): Mệnh đề này đúng vì động năng của vật là $W_d = W - W_t = 0,5 - 0,18 = 0,32$ J
\item (Sai) Khi vật nặng có li độ $x = 6$ cm thì tốc độ chuyển động của vật nặng bằng $80$ m/s. (Vì): Mệnh đề này sai vì động năng $W_d = \dfrac{1}{2}mv^2 \Rightarrow 0,32 = 0,5 \cdot 1 \cdot v^2 \Rightarrow v^2 = 0,64 \Rightarrow v = 0,8$ m/s
\end{itemchoice}
}
\end{ex}

% ===== Câu 261 =====
% ID: L11C1B5-70-TLN
% Mức: VD
\dangbt{Xác định độ cứng lò xo từ năng lượng ban đầu}
\begin{ex}
% ID: L11C1B5-70-TLN
% Mức: VD
Một con lắc lò xo có cơ năng 0,18 $J$ và biên độ $6\,\text{cm}$. Độ cứng của lò xo theo đơn vị N/m bằng bao nhiêu?
\shortans{100}
\loigiai{
$k = 2W/A² = 2·0,18/(0,06)² = 100 N/m.$
}
\end{ex}

% ===== Câu 262 =====
% ID: L11C1B5-70-TN
% Mức: TH
\dangbt{Tính động năng, thế năng và cơ năng tại một li độ hoặc thời điểm xác định}
\begin{ex}
% ID: L11C1B5-70-TN
% Mức: TH
Một vật dao động điều hòa theo một trục cố định (mốc thế năng ở vị trí cân bằng) thì
\choice
{động năng của vật cực đại khi gia tốc của vật có độ lớn cực đại}
{\True thế năng của vật cực đại khi vật ở vị trí biên}
{khi ở vị trí cân bằng, thế năng của vật bằng cơ năng}
{khi vật đi từ vị trí cân bằng ra biên, vận tốc và gia tốc của vật luôn cùng dấu}
\loigiai{
Thế năng của vật dao động điều hòa được tính bằng công thức $W_t = \frac{1}{2}kx^2$. Tại vị trí biên, li độ $x$ đạt giá trị cực đại ($x = \pm A$), do đó thế năng đạt giá trị cực đại. Các phương án còn lại là sai:
  - Khi vật đi từ vị trí cân bằng ra biên, vận tốc và gia tốc của vật luôn ngược dấu (ví dụ, khi vật đi từ VTCB ra biên dương, $v {\gt } 0$ và $a {\lt } 0$).
  - Khi ở vị trí cân bằng ($x=0$), thế năng của vật bằng 0, động năng của vật đạt cực đại và bằng cơ năng.
  - Động năng của vật cực đại khi vật ở vị trí cân bằng ($x=0$), tại đó gia tốc của vật bằng 0 (độ lớn cực tiểu), không phải cực đại. Gia tốc của vật có độ lớn cực đại tại vị trí biên, nơi động năng bằng 0.
}
\end{ex}

% ===== Câu 263 =====
% ID: L11C1B5-71-DS
% Mức: TH
\dangbt{Xác định biểu thức cơ năng trong dao động điều hòa}
\begin{ex}
% ID: L11C1B5-71-DS
% Mức: TH
Một vật dao động điều hòa có cơ năng $W$. Tại li độ $x=A/2$, xác định tính đúng sai.
\choiceTF
{\True Thế năng bằng $W/4$.}
{\True Động năng bằng $3W/4$.}
{\True Động năng gấp ba lần thế năng.}
{Cơ năng giảm còn $W/2$.}
\end{ex}

% ===== Câu 264 =====
% ID: L11C1B5-71-TLN
% Mức: VD
\dangbt{Xác định độ cứng lò xo từ năng lượng ban đầu}
\begin{ex}
% ID: L11C1B5-71-TLN
% Mức: VD
Một con lắc lò xo có cơ năng 0,18 $J$ và biên độ $6\,\text{cm}$. Độ cứng của lò xo theo đơn vị N/m bằng bao nhiêu?
\shortans{100}
\loigiai{
$k = 2W/A² = 2·0,18/(0,06)² = 100 N/m.$
}
\end{ex}

% ===== Câu 265 =====
% ID: L11C1B5-71-TN
% Mức: VD
\dangbt{Tính động năng, thế năng và cơ năng tại một li độ hoặc thời điểm xác định}
\begin{ex}
% ID: L11C1B5-71-TN
% Mức: VD
Một con lắc lò xo có độ cứng $k = 120$ N/m, vật nặng khối lượng $m = 300$ g dao động điều hòa với biên độ $A = 5$ cm. Tại thời điểm vận tốc của vật bằng $\frac{\sqrt{3}}{2}$ lần vận tốc cực đại, thế năng của vật bằng bao nhiêu?
\choice
{\True $3,75 \times 10^{-3}J$}
{$1,5 \times 10^{-2}J$}
{$7,5 \times 10^{-3}J$}
{$1,125 \times 10^{-2}J$}
\loigiai{
Cơ năng toàn phần của con lắc lò xo: $W = \frac{1}{2}kA^2 = \frac{1}{2} \times 120 \times (0,05)^2 = 0,15$ J. Vận tốc cực đại: $v_{max} = A\omega$, nên động năng cực đại $W_{\text{đ max}} = \frac{1}{2}mv_{max}^2 = W = 0,15$ J. Khi $v = \frac{\sqrt{3}}{2}v_{max}$, động năng lúc đó: $W_{\text{đ}} = \frac{1}{2}mv^2 = \frac{1}{2}m \cdot \frac{3}{4}v_{max}^2 = \frac{3}{4} \cdot \frac{1}{2}mv_{max}^2 = \frac{3}{4} \times 0,15 = 0,1125$ J. Thế năng: $W_t = W - W_{\text{đ}} = 0,15 - 0,1125 = 0,0375J= 3,75 \times 10^{-2}$ J. Kiểm tra lại: $W_t = \frac{1}{4} \times 0,15 = 0,0375J= 3,75 \times 10^{-2}$ J. Vậy đáp án đúng là $3,75 \times 10^{-2}$ J. Lưu ý: đáp án $A$ ghi $3,75 \times 10^{-2}J$ (tương đương $37,5$ mJ). Đáp án đúng là A.
}
\end{ex}

% ===== Câu 266 =====
% ID: L11C1B5-72-DS
% Mức: VD
\dangbt{Xác định biểu thức cơ năng trong dao động điều hòa}
\begin{ex}
% ID: L11C1B5-72-DS
% Mức: VD
Một vật dao động điều hòa có cơ năng $W$. Tại li độ $x=A/2$, xác định tính đúng sai.
\choiceTF
{\True Thế năng bằng $W/4$.}
{\True Động năng bằng $3W/4$.}
{\True Động năng gấp ba lần thế năng.}
{Cơ năng giảm còn $W/2$.}
\end{ex}

% ===== Câu 267 =====
% ID: L11C1B5-72-TLN
% Mức: VD
\dangbt{Xác định độ cứng lò xo từ năng lượng ban đầu}
\begin{ex}
% ID: L11C1B5-72-TLN
% Mức: VD
Một con lắc lò xo có cơ năng 0,18 $J$ và biên độ $6\,\text{cm}$. Độ cứng của lò xo theo đơn vị N/m bằng bao nhiêu?
\shortans{100}
\loigiai{
$k = 2W/A² = 2·0,18/(0,06)² = 100 N/m.$
}
\end{ex}

% ===== Câu 268 =====
% ID: L11C1B5-72-TN
% Mức: NB
\dangbt{Xác định biểu thức cơ năng trong dao động điều hòa}
\begin{ex}
% ID: L11C1B5-72-TN
% Mức: NB
Một con lắc lò xo gồm vật nhỏ và lò xo nhẹ có độ cứng k, đang dao động điều hòa. Mốc thế năng tại vị trí cân bằng. Biểu thức thế năng của con lắc ở li độ x là
\choice
{$W_t = 2kx^2$}
{\True $W_t = \dfrac{1}{2}kx^2$}
{$W_t = 2kx$}
{$W_t = \dfrac{1}{2}kx$}
\loigiai{
Thế năng của con lắc lò xo ở li độ x được xác định bởi công thức $W_t = \dfrac{1}{2}kx^2$, với k là độ cứng của lò xo và x là li độ.
}
\end{ex}

% ===== Câu 269 =====
% ID: L11C1B5-73-DS
% Mức: VD
\begin{ex}
% ID: L11C1B5-73-DS
% Mức: VD
Một vật dao động điều hòa có cơ năng $W$. Tại li độ $x=A/2$, xác định tính đúng sai.
\choiceTF
{\True Thế năng bằng $W/4$.}
{\True Động năng bằng $3W/4$.}
{\True Động năng gấp ba lần thế năng.}
{Cơ năng giảm còn $W/2$.}
\end{ex}

% ===== Câu 270 =====
% ID: L11C1B5-73-TLN
% Mức: VD
\dangbt{Xác định độ cứng lò xo từ năng lượng ban đầu}
\begin{ex}
% ID: L11C1B5-73-TLN
% Mức: VD
Một con lắc lò xo có vật nặng $0,5\,\text{kg}$, biên độ $5\,\text{cm}$. Tại li độ $3\,\text{cm}$, tốc độ của vật là $0,4\,\text{m}$ /s. Độ cứng của lò xo theo đơn vị N/m bằng bao nhiêu?
\shortans{50}
\loigiai{
Bảo toàn cơ năng: k(A² - x²) = mv², nên k = 0,5·0,4²/(0,05² - 0,03²) = $50\,\text{N}$ /m.
}
\end{ex}

% ===== Câu 271 =====
% ID: L11C1B5-73-TN
% Mức: VD
\dangbt{Xác định biểu thức cơ năng trong dao động điều hòa}
\begin{ex}
% ID: L11C1B5-73-TN
% Mức: VD
Một chất điểm khối lượng $m=100g$ dao động điều hoà với phương trình $x=4\cos (2t)$ cm. Cơ năng trong dao động điều hoà của chất điểm là
\choice
{0,32 $J$}
{3,2 $J$}
{\True 0,32 mJ}
{3200 $J$}
\loigiai{
Cơ năng: $W = \dfrac{1}{2}m\omega^2 A^2 = \dfrac{1}{2} \cdot 0,1 \cdot 2^2 \cdot (0,04)^2 = 0,32 \cdot 10^{-3} \mathrm{J} = 0,32 \mathrm{mJ}$
}
\end{ex}

% ===== Câu 272 =====
% ID: L11C1B5-74-DS
% Mức: TH
\dangbt{Xác định và tính toán các đại lượng vận tốc, gia tốc tại thời điểm xác định}
\begin{ex}
% ID: L11C1B5-74-DS
% Mức: TH
Vật dao động theo $x=6\cos(3t)$ cm.
\choiceTF
{\True $v=-18\sin(3t)$ cm/s.}
{\True $a=-54\cos(3t)$ cm/s².}
{\True Tốc độ cực đại $18\,\text{cm}$ /s.}
{Gia tốc bằng 0 ở biên.}
\end{ex}

% ===== Câu 273 =====
% ID: L11C1B5-74-TLN
% Mức: VD
\dangbt{Xác định độ cứng lò xo từ năng lượng ban đầu}
\begin{ex}
% ID: L11C1B5-74-TLN
% Mức: VD
Một con lắc lò xo có vật nặng $0,5\,\text{kg}$, biên độ $5\,\text{cm}$. Tại li độ $3\,\text{cm}$, tốc độ của vật là $0,4\,\text{m}$ /s. Độ cứng của lò xo theo đơn vị N/m bằng bao nhiêu?
\shortans{50}
\loigiai{
Bảo toàn cơ năng: k(A² - x²) = mv², nên k = 0,5·0,4²/(0,05² - 0,03²) = $50\,\text{N}$ /m.
}
\end{ex}

% ===== Câu 274 =====
% ID: L11C1B5-74-TN
% Mức: VD
\dangbt{Xác định biểu thức cơ năng trong dao động điều hòa}
\begin{ex}
% ID: L11C1B5-74-TN
% Mức: VD
Một vật nhỏ khối lượng $100\,\text{g}$ dao động điều hoà trên một quỹ đạo thẳng dài $20\,\text{cm}$ với tần số góc 6 rad/s. Cơ năng của vật dao động này là
\choice
{\True 0,018 $J$}
{0,036 $J$}
{0,18 $J$}
{0,36 $J$}
\loigiai{
Biên độ dao động $A = \dfrac{L}{2} = \dfrac{20}{2} = 10\,\text{cm} = 0,1\,\text{m}$. Khối lượng $m = 100\,\text{g} = 0,1\,\text{kg}$. Tần số góc $\omega = 6\,\text{rad/s}$. Cơ năng của vật dao động điều hòa được tính bằng công thức $W = \dfrac{1}{2}m\omega^2A^2$. Thay số vào ta có: $W = \dfrac{1}{2} \times 0,1 \times 6^2 \times (0,1)^2 = \dfrac{1}{2} \times 0,1 \times 36 \times 0,01 = 0,018\,\text{J}$.
}
\end{ex}

% ===== Câu 275 =====
% ID: L11C1B5-75-DS
% Mức: VD
\dangbt{Xác định và tính toán các đại lượng vận tốc, gia tốc tại thời điểm xác định}
\begin{ex}
% ID: L11C1B5-75-DS
% Mức: VD
Vật dao động theo $x=7\cos(4t)$ cm.
\choiceTF
{\True $v=-28\sin(4t)$ cm/s.}
{\True $a=-112\cos(4t)$ cm/s².}
{\True Tốc độ cực đại $28\,\text{cm}$ /s.}
{Gia tốc bằng 0 ở biên.}
\end{ex}

% ===== Câu 276 =====
% ID: L11C1B5-75-TLN
% Mức: VD
\dangbt{Xác định độ cứng lò xo từ năng lượng ban đầu}
\begin{ex}
% ID: L11C1B5-75-TLN
% Mức: VD
Một con lắc lò xo có vật nặng $0,5\,\text{kg}$, biên độ $5\,\text{cm}$. Tại li độ $3\,\text{cm}$, tốc độ của vật là $0,4\,\text{m}$ /s. Độ cứng của lò xo theo đơn vị N/m bằng bao nhiêu?
\shortans{50}
\loigiai{
Bảo toàn cơ năng: k(A² - x²) = mv², nên k = 0,5·0,4²/(0,05² - 0,03²) = $50\,\text{N}$ /m.
}
\end{ex}

% ===== Câu 277 =====
% ID: L11C1B5-75-TN
% Mức: TH
\dangbt{Xác định và tính toán các đại lượng vận tốc, gia tốc tại thời điểm xác định}
\begin{ex}
% ID: L11C1B5-75-TN
% Mức: TH
Vật dao động theo $x=10\cos(5t)$ cm. Tốc độ cực đại bằng
\choice
{\True $50\,\text{cm}$ /s}
{$10\,\text{cm}$ /s}
{$5\,\text{cm}$ /s}
{$250\,\text{cm}$ /s}
\end{ex}

% ===== Câu 278 =====
% ID: L11C1B5-76-DS
% Mức: VD
\begin{ex}
% ID: L11C1B5-76-DS
% Mức: VD
Vật dao động theo $x=8\cos(5t)$ cm.
\choiceTF
{\True $v=-40\sin(5t)$ cm/s.}
{\True $a=-200\cos(5t)$ cm/s².}
{\True Tốc độ cực đại $40\,\text{cm}$ /s.}
{Gia tốc bằng 0 ở biên.}
\end{ex}

% ===== Câu 279 =====
% ID: L11C1B5-76-TN
% Mức: TH
\begin{ex}
% ID: L11C1B5-76-TN
% Mức: TH
Thời gian ngắn nhất vật đi từ vị trí cân bằng đến vị trí có li độ $\dfrac{A}{2}$ là
\choice
{$\dfrac{T}{4}.$}
{$\dfrac{T}{6}.$}
{$\dfrac{T}{8}.$}
{\True $\dfrac{T}{12}.$}
\loigiai{
Thời gian ngắn nhất vật đi từ vị trí cân bằng đến vị trí có li độ $\dfrac{A}{2}$ là $\dfrac{T}{12}.$
}
\end{ex}

% ===== Câu 280 =====
% ID: L11C1B5-77-DS
% Mức: TH
\dangbt{Xác định vị trí dựa vào tỉ số động năng và thế năng}
\begin{ex}
% ID: L11C1B5-77-DS
% Mức: TH
Khi một vật dao động điều hòa quanh vị trí cân bằng với biên độ $A$, ta xét tỉ số giữa động năng và thế năng tại một số vị trí đặc trưng:
\choiceTF
{Tại vị trí có li độ $x = \pm \dfrac{A}{2}$ thì động năng bằng thế năng.}
{\True Tại vị trí có li độ $x = \pm \dfrac{A}{2}$ thì động năng bằng ba lần thế năng của hệ.}
{\True Tại vị trí có li độ $x = \pm \dfrac{A\sqrt{3}}{2}$ thì thế năng bằng ba lần động năng của vật.}
{Vị trí vật có động năng bằng thế năng nằm chính giữa vị trí cân bằng và vị trí biên.}
\loigiai{
\begin{itemchoice}
\item (Sai) Tại vị trí có li độ $x = \pm \dfrac{A}{2}$ thì động năng bằng thế năng. (Vì): Mệnh đề này sai vì khi động năng bằng thế năng ($W_d = W_t$) thì cơ năng $W = 2W_t \Rightarrow x = \pm \dfrac{A}{\sqrt{2}}$
\item (Đúng) Tại vị trí có li độ $x = \pm \dfrac{A}{2}$ thì động năng bằng ba lần thế năng của hệ. (Vì): Mệnh đề này đúng vì khi $x = \pm \dfrac{A}{2}$ thì thế năng $W_t = \dfrac{1}{4}W \Rightarrow W_d = W - W_t = \dfrac{3}{4}W = 3W_t$
\item (Đúng) Tại vị trí có li độ $x = \pm \dfrac{A\sqrt{3}}{2}$ thì thế năng bằng ba lần động năng của vật. (Vì): Mệnh đề này đúng vì khi $x = \pm \dfrac{A\sqrt{3}}{2}$ thì thế năng $W_t = \dfrac{3}{4}W \Rightarrow W_d = W - W_t = \dfrac{1}{4}W \Rightarrow W_t = 3W_d$
\item (Sai) Vị trí vật có động năng bằng thế năng nằm chính giữa vị trí cân bằng và vị trí biên. (Vì): Mệnh đề này sai vì vị trí có động năng bằng thế năng là $x \approx 0,707A$, lệch xa hơn so với trung điểm của đoạn nối vị trí cân bằng và biên ($0,5A$
\end{itemchoice}
}
\end{ex}

% ===== Câu 281 =====
% ID: L11C1B5-77-TN
% Mức: VD
\dangbt{Xác định và tính toán các đại lượng vận tốc, gia tốc tại thời điểm xác định}
\begin{ex}
% ID: L11C1B5-77-TN
% Mức: VD
Vật dao động theo $x=5\cos(2t)$ cm. Tốc độ cực đại bằng
\choice
{\True $10\,\text{cm}$ /s}
{$5\,\text{cm}$ /s}
{$2\,\text{cm}$ /s}
{$20\,\text{cm}$ /s}
\end{ex}

% ===== Câu 282 =====
% ID: L11C1B5-78-DS
% Mức: TH
\dangbt{Xác định vị trí dựa vào tỉ số động năng và thế năng}
\begin{ex}
% ID: L11C1B5-78-DS
% Mức: TH
Khi một vật dao động điều hòa quanh vị trí cân bằng với biên độ $A$, ta xét tỉ số giữa động năng và thế năng tại một số vị trí đặc trưng:
\choiceTF
{Tại vị trí có li độ $x = \pm \dfrac{A}{2}$ thì động năng bằng thế năng.}
{\True Tại vị trí có li độ $x = \pm \dfrac{A}{2}$ thì động năng bằng ba lần thế năng của hệ.}
{\True Tại vị trí có li độ $x = \pm \dfrac{A\sqrt{3}}{2}$ thì thế năng bằng ba lần động năng của vật.}
{Vị trí vật có động năng bằng thế năng nằm chính giữa vị trí cân bằng và vị trí biên.}
\loigiai{
\begin{itemchoice}
\item (Sai) Tại vị trí có li độ $x = \pm \dfrac{A}{2}$ thì động năng bằng thế năng. (Vì): Mệnh đề này sai vì khi động năng bằng thế năng ($W_d = W_t$) thì cơ năng $W = 2W_t \Rightarrow x = \pm \dfrac{A}{\sqrt{2}}$
\item (Đúng) Tại vị trí có li độ $x = \pm \dfrac{A}{2}$ thì động năng bằng ba lần thế năng của hệ. (Vì): Mệnh đề này đúng vì khi $x = \pm \dfrac{A}{2}$ thì thế năng $W_t = \dfrac{1}{4}W \Rightarrow W_d = W - W_t = \dfrac{3}{4}W = 3W_t$
\item (Đúng) Tại vị trí có li độ $x = \pm \dfrac{A\sqrt{3}}{2}$ thì thế năng bằng ba lần động năng của vật. (Vì): Mệnh đề này đúng vì khi $x = \pm \dfrac{A\sqrt{3}}{2}$ thì thế năng $W_t = \dfrac{3}{4}W \Rightarrow W_d = W - W_t = \dfrac{1}{4}W \Rightarrow W_t = 3W_d$
\item (Sai) Vị trí vật có động năng bằng thế năng nằm chính giữa vị trí cân bằng và vị trí biên. (Vì): Mệnh đề này sai vì vị trí có động năng bằng thế năng là $x \approx 0,707A$, lệch xa hơn so với trung điểm của đoạn nối vị trí cân bằng và biên ($0,5A$
\end{itemchoice}
}
\end{ex}

% ===== Câu 283 =====
% ID: L11C1B5-78-TN
% Mức: TH
\begin{ex}
% ID: L11C1B5-78-TN
% Mức: TH
Một vật dao động điều hòa với biên độ $A$. Vị trí mà tại đó thế năng của vật bằng động năng của vật thỏa mãn hệ thức nào sau đây?
\choice
{$x = \pm \dfrac{A}{2}$}
{$x = \pm \dfrac{A\sqrt{3}}{2}$}
{\True $x = \pm \dfrac{A}{\sqrt{2}}$}
{$x = \pm \dfrac{A}{4}$}
\loigiai{
• Điều kiện đề bài: $W_{\text{t}} = W_{\text{đ}}$.
 
• Áp dụng bảo toàn cơ năng:
  \begin{aligned}
 $W$ &= W_{\text{đ}} + W_{ $\text{t}$ } = 2W_{ $\text{t}$ } 
$\Rightarrow\dfrac{1}{2}m\omega^2$ A^2 &= 2 $\cdot\dfrac{1}{2}m\omega^2$ x^2.
 \end{aligned} 
 
• Rút gọn hai vế:
 $A^2 = 2x^2 \Rightarrow x = \pm \dfrac{A}{\sqrt{2}}.$
}
\end{ex}

% ===== Câu 284 =====
% ID: L11C1B5-79-DS
% Mức: TH
\begin{ex}
% ID: L11C1B5-79-DS
% Mức: TH
Khi một vật dao động điều hòa quanh vị trí cân bằng với biên độ $A$, ta xét tỉ số giữa động năng và thế năng tại một số vị trí đặc trưng:
\choiceTF
{Tại vị trí có li độ $x = \pm \dfrac{A}{2}$ thì động năng bằng thế năng.}
{\True Tại vị trí có li độ $x = \pm \dfrac{A}{2}$ thì động năng bằng ba lần thế năng của hệ.}
{\True Tại vị trí có li độ $x = \pm \dfrac{A\sqrt{3}}{2}$ thì thế năng bằng ba lần động năng của vật.}
{Vị trí vật có động năng bằng thế năng nằm chính giữa vị trí cân bằng và vị trí biên.}
\loigiai{
\begin{itemchoice}
\item (Sai) Tại vị trí có li độ $x = \pm \dfrac{A}{2}$ thì động năng bằng thế năng. (Vì): Mệnh đề này sai vì khi động năng bằng thế năng ($W_d = W_t$) thì cơ năng $W = 2W_t \Rightarrow x = \pm \dfrac{A}{\sqrt{2}}$
\item (Đúng) Tại vị trí có li độ $x = \pm \dfrac{A}{2}$ thì động năng bằng ba lần thế năng của hệ. (Vì): Mệnh đề này đúng vì khi $x = \pm \dfrac{A}{2}$ thì thế năng $W_t = \dfrac{1}{4}W \Rightarrow W_d = W - W_t = \dfrac{3}{4}W = 3W_t$
\item (Đúng) Tại vị trí có li độ $x = \pm \dfrac{A\sqrt{3}}{2}$ thì thế năng bằng ba lần động năng của vật. (Vì): Mệnh đề này đúng vì khi $x = \pm \dfrac{A\sqrt{3}}{2}$ thì thế năng $W_t = \dfrac{3}{4}W \Rightarrow W_d = W - W_t = \dfrac{1}{4}W \Rightarrow W_t = 3W_d$
\item (Sai) Vị trí vật có động năng bằng thế năng nằm chính giữa vị trí cân bằng và vị trí biên. (Vì): Mệnh đề này sai vì vị trí có động năng bằng thế năng là $x \approx 0,707A$, lệch xa hơn so với trung điểm của đoạn nối vị trí cân bằng và biên ($0,5A$
\end{itemchoice}
}
\end{ex}

% ===== Câu 285 =====
% ID: L11C1B5-79-TN
% Mức: TH
\begin{ex}
% ID: L11C1B5-79-TN
% Mức: TH
Một vật dao động điều hòa với biên độ $A$. Vị trí mà tại đó thế năng của vật bằng động năng của vật thỏa mãn hệ thức nào sau đây?
\choice
{$x = \pm \dfrac{A}{2}$}
{$x = \pm \dfrac{A\sqrt{3}}{2}$}
{\True $x = \pm \dfrac{A}{\sqrt{2}}$}
{$x = \pm \dfrac{A}{4}$}
\loigiai{
• Điều kiện đề bài: $W_{\text{t}} = W_{\text{đ}}$.
 
• Áp dụng bảo toàn cơ năng:
  \begin{aligned}
 $W$ &= W_{\text{đ}} + W_{ $\text{t}$ } = 2W_{ $\text{t}$ } 
$\Rightarrow\dfrac{1}{2}m\omega^2$ A^2 &= 2 $\cdot\dfrac{1}{2}m\omega^2$ x^2.
 \end{aligned} 
 
• Rút gọn hai vế:
 $A^2 = 2x^2 \Rightarrow x = \pm \dfrac{A}{\sqrt{2}}.$
}
\end{ex}

% ===== Câu 286 =====
% ID: L11C1B5-80-DS
% Mức: TH
\begin{ex}
% ID: L11C1B5-80-DS
% Mức: TH
Một con lắc đơn dao động điều hòa quanh vị trí cân bằng với biên độ góc là $\alpha_0$. Chọn mốc thế năng trọng trường tại vị trí cân bằng của quả cầu.
\choiceTF
{\True Khi đi qua vị trí có động năng bằng thế năng, góc lệch dây treo của con lắc là $\alpha = \pm \dfrac{\alpha_0}{\sqrt{2}}$.}
{Tại vị trí dây treo lệch góc $\alpha = \pm \dfrac{\alpha_0}{2}$ so với phương đứng, thế năng bằng ba lần động năng.}
{\True Khi quả cầu đi qua vị trí cân bằng, động năng đạt cực đại và thế năng bằng không.}
{\True Khi vật nặng ở vị trí biên, thế năng đạt cực đại và tốc độ của quả cầu bằng không.}
\loigiai{
\begin{itemchoice}
\item (Đúng) Khi đi qua vị trí có động năng bằng thế năng, góc lệch dây treo của con lắc là $\alpha = \pm \dfrac{\alpha_0}{\sqrt{2}}$. (Vì): Mệnh đề này đúng vì khi thế năng bằng động năng thì góc lệch li độ thỏa mãn hệ thức: $\alpha = \pm \dfrac{\alpha_0}{\sqrt{2}}$
\item (Sai) Tại vị trí dây treo lệch góc $\alpha = \pm \dfrac{\alpha_0}{2}$ so với phương đứng, thế năng bằng ba lần động năng. (Vì): Mệnh đề này sai vì tại góc lệch $\alpha = \pm \dfrac{\alpha_0}{2}$ thì thế năng bằng một phần tư cơ năng, suy ra động năng bằng ba phần tư cơ năng ($W_d = 3W_t$ hay thế năng bằng một phần ba động năng
\item (Đúng) Khi quả cầu đi qua vị trí cân bằng, động năng đạt cực đại và thế năng bằng không. (Vì): Mệnh đề này đúng vì qua vị trí cân bằng vật có li độ góc bằng không nên thế năng trọng trường bằng không, đồng thời tốc độ đạt cực đại nên động năng cực đại
\item (Đúng) Khi vật nặng ở vị trí biên, thế năng đạt cực đại và tốc độ của quả cầu bằng không. (Vì): Mệnh đề này đúng vì tại biên vật đổi chiều chuyển động nên vận tốc bằng không (động năng bằng không), độ dời đạt cực đại nên thế năng cực đại
\end{itemchoice}
}
\end{ex}

% ===== Câu 287 =====
% ID: L11C1B5-80-TN
% Mức: TH
\begin{ex}
% ID: L11C1B5-80-TN
% Mức: TH
Một vật dao động điều hòa với biên độ $A$. Vị trí mà tại đó thế năng của vật bằng động năng của vật thỏa mãn hệ thức nào sau đây?
\choice
{$x = \pm \dfrac{A}{2}$}
{$x = \pm \dfrac{A\sqrt{3}}{2}$}
{\True $x = \pm \dfrac{A}{\sqrt{2}}$}
{$x = \pm \dfrac{A}{4}$}
\loigiai{
• Điều kiện đề bài: $W_{\text{t}} = W_{\text{đ}}$.
 
• Áp dụng bảo toàn cơ năng:
  \begin{aligned}
 $W$ &= W_{\text{đ}} + W_{ $\text{t}$ } = 2W_{ $\text{t}$ } 
$\Rightarrow\dfrac{1}{2}m\omega^2$ A^2 &= 2 $\cdot\dfrac{1}{2}m\omega^2$ x^2.
 \end{aligned} 
 
• Rút gọn hai vế:
 $A^2 = 2x^2 \Rightarrow x = \pm \dfrac{A}{\sqrt{2}}.$
}
\end{ex}

% ===== Câu 288 =====
% ID: L11C1B5-81-DS
% Mức: TH
\begin{ex}
% ID: L11C1B5-81-DS
% Mức: TH
Một con lắc đơn dao động điều hòa quanh vị trí cân bằng với biên độ góc là $\alpha_0$. Chọn mốc thế năng trọng trường tại vị trí cân bằng của quả cầu.
\choiceTF
{\True Khi đi qua vị trí có động năng bằng thế năng, góc lệch dây treo của con lắc là $\alpha = \pm \dfrac{\alpha_0}{\sqrt{2}}$.}
{Tại vị trí dây treo lệch góc $\alpha = \pm \dfrac{\alpha_0}{2}$ so với phương đứng, thế năng bằng ba lần động năng.}
{\True Khi quả cầu đi qua vị trí cân bằng, động năng đạt cực đại và thế năng bằng không.}
{\True Khi vật nặng ở vị trí biên, thế năng đạt cực đại và tốc độ của quả cầu bằng không.}
\loigiai{
\begin{itemchoice}
\item (Đúng) Khi đi qua vị trí có động năng bằng thế năng, góc lệch dây treo của con lắc là $\alpha = \pm \dfrac{\alpha_0}{\sqrt{2}}$. (Vì): Mệnh đề này đúng vì khi thế năng bằng động năng thì góc lệch li độ thỏa mãn hệ thức: $\alpha = \pm \dfrac{\alpha_0}{\sqrt{2}}$
\item (Sai) Tại vị trí dây treo lệch góc $\alpha = \pm \dfrac{\alpha_0}{2}$ so với phương đứng, thế năng bằng ba lần động năng. (Vì): Mệnh đề này sai vì tại góc lệch $\alpha = \pm \dfrac{\alpha_0}{2}$ thì thế năng bằng một phần tư cơ năng, suy ra động năng bằng ba phần tư cơ năng ($W_d = 3W_t$ hay thế năng bằng một phần ba động năng
\item (Đúng) Khi quả cầu đi qua vị trí cân bằng, động năng đạt cực đại và thế năng bằng không. (Vì): Mệnh đề này đúng vì qua vị trí cân bằng vật có li độ góc bằng không nên thế năng trọng trường bằng không, đồng thời tốc độ đạt cực đại nên động năng cực đại
\item (Đúng) Khi vật nặng ở vị trí biên, thế năng đạt cực đại và tốc độ của quả cầu bằng không. (Vì): Mệnh đề này đúng vì tại biên vật đổi chiều chuyển động nên vận tốc bằng không (động năng bằng không), độ dời đạt cực đại nên thế năng cực đại
\end{itemchoice}
}
\end{ex}

% ===== Câu 289 =====
% ID: L11C1B5-81-TN
% Mức: VD
\begin{ex}
% ID: L11C1B5-81-TN
% Mức: VD
Một con lắc đơn dao động điều hòa với biên độ góc $\alpha_0$. Khi vật đi qua vị trí có li độ góc $\alpha = \pm \dfrac{\alpha_0}{\sqrt{2}}$ thì tỉ số giữa động năng và thế năng của con lắc bằng:
\choice
{$2$}
{$\dfrac{1}{2}$}
{\True $1$}
{$3$}
\loigiai{
• Thế năng của con lắc đơn ở li độ góc $\alpha = \pm \dfrac{\alpha_0}{\sqrt{2}}$:
 $W_{\text{t}} = \dfrac{1}{2}mgl\alpha^2 = \dfrac{1}{2}mgl\left(\dfrac{\alpha_0}{\sqrt{2}}\right)^2 = \dfrac{1}{2} \cdot \left(\dfrac{1}{2}mgl\alpha_0^2\right) = \dfrac{W}{2}.$
 
• Động năng, theo bảo toàn cơ năng:
 $W_{\text{đ}} = W - W_{\text{t}} = W - \dfrac{W}{2} = \dfrac{W}{2}.$
 
• Lập tỉ số:
 $\dfrac{W_{\text{đ}}}{W_{\text{t}}} = \dfrac{\frac{W}{2}}{\frac{W}{2}} = 1.$
}
\end{ex}

% ===== Câu 290 =====
% ID: L11C1B5-82-DS
% Mức: TH
\begin{ex}
% ID: L11C1B5-82-DS
% Mức: TH
Một con lắc đơn dao động điều hòa quanh vị trí cân bằng với biên độ góc là $\alpha_0$. Chọn mốc thế năng trọng trường tại vị trí cân bằng của quả cầu.
\choiceTF
{\True Khi đi qua vị trí có động năng bằng thế năng, góc lệch dây treo của con lắc là $\alpha = \pm \dfrac{\alpha_0}{\sqrt{2}}$.}
{Tại vị trí dây treo lệch góc $\alpha = \pm \dfrac{\alpha_0}{2}$ so với phương đứng, thế năng bằng ba lần động năng.}
{\True Khi quả cầu đi qua vị trí cân bằng, động năng đạt cực đại và thế năng bằng không.}
{\True Khi vật nặng ở vị trí biên, thế năng đạt cực đại và tốc độ của quả cầu bằng không.}
\loigiai{
\begin{itemchoice}
\item (Đúng) Khi đi qua vị trí có động năng bằng thế năng, góc lệch dây treo của con lắc là $\alpha = \pm \dfrac{\alpha_0}{\sqrt{2}}$. (Vì): Mệnh đề này đúng vì khi thế năng bằng động năng thì góc lệch li độ thỏa mãn hệ thức: $\alpha = \pm \dfrac{\alpha_0}{\sqrt{2}}$
\item (Sai) Tại vị trí dây treo lệch góc $\alpha = \pm \dfrac{\alpha_0}{2}$ so với phương đứng, thế năng bằng ba lần động năng. (Vì): Mệnh đề này sai vì tại góc lệch $\alpha = \pm \dfrac{\alpha_0}{2}$ thì thế năng bằng một phần tư cơ năng, suy ra động năng bằng ba phần tư cơ năng ($W_d = 3W_t$ hay thế năng bằng một phần ba động năng
\item (Đúng) Khi quả cầu đi qua vị trí cân bằng, động năng đạt cực đại và thế năng bằng không. (Vì): Mệnh đề này đúng vì qua vị trí cân bằng vật có li độ góc bằng không nên thế năng trọng trường bằng không, đồng thời tốc độ đạt cực đại nên động năng cực đại
\item (Đúng) Khi vật nặng ở vị trí biên, thế năng đạt cực đại và tốc độ của quả cầu bằng không. (Vì): Mệnh đề này đúng vì tại biên vật đổi chiều chuyển động nên vận tốc bằng không (động năng bằng không), độ dời đạt cực đại nên thế năng cực đại
\end{itemchoice}
}
\end{ex}

% ===== Câu 291 =====
% ID: L11C1B5-82-TN
% Mức: VD
\begin{ex}
% ID: L11C1B5-82-TN
% Mức: VD
Một con lắc đơn dao động điều hòa với biên độ góc $\alpha_0$. Khi vật đi qua vị trí có li độ góc $\alpha = \pm \dfrac{\alpha_0}{\sqrt{2}}$ thì tỉ số giữa động năng và thế năng của con lắc bằng:
\choice
{$2$}
{$\dfrac{1}{2}$}
{\True $1$}
{$3$}
\loigiai{
• Thế năng của con lắc đơn ở li độ góc $\alpha = \pm \dfrac{\alpha_0}{\sqrt{2}}$:
 $W_{\text{t}} = \dfrac{1}{2}mgl\alpha^2 = \dfrac{1}{2}mgl\left(\dfrac{\alpha_0}{\sqrt{2}}\right)^2 = \dfrac{1}{2} \cdot \left(\dfrac{1}{2}mgl\alpha_0^2\right) = \dfrac{W}{2}.$
 
• Động năng, theo bảo toàn cơ năng:
 $W_{\text{đ}} = W - W_{\text{t}} = W - \dfrac{W}{2} = \dfrac{W}{2}.$
 
• Lập tỉ số:
 $\dfrac{W_{\text{đ}}}{W_{\text{t}}} = \dfrac{\frac{W}{2}}{\frac{W}{2}} = 1.$
}
\end{ex}

% ===== Câu 292 =====
% ID: L11C1B5-83-DS
% Mức: VD
\begin{ex}
% ID: L11C1B5-83-DS
% Mức: VD
Một con lắc lò xo dao động điều hòa với biên độ $A = 10$ cm và tần số góc $\omega = 20$ rad/s. Khảo sát động học của vật tại các vị trí phân chia năng lượng:
\choiceTF
{\True Tại vị trí có li độ $x = 5$ cm, động năng của vật bằng ba lần thế năng.}
{\True Tại vị trí có li độ $x = 5\sqrt{2}$ cm, động năng của vật bằng thế năng của con lắc.}
{\True Tốc độ của vật tại vị trí có động năng bằng thế năng bằng $100\sqrt{2}$ cm/s.}
{Tốc độ của vật tại vị trí có động năng bằng ba lần thế năng bằng $100$ cm/s.}
\loigiai{
\begin{itemchoice}
\item (Đúng) Tại vị trí có li độ $x = 5$ cm, động năng của vật bằng ba lần thế năng. (Vì): Mệnh đề này đúng vì $x = 5\,\mathrm{cm}= A/2$, tương ứng với vị trí có động năng bằng 3 lần thế năng
\item (Đúng) Tại vị trí có li độ $x = 5\sqrt{2}$ cm, động năng của vật bằng thế năng của con lắc. (Vì): Mệnh đề này đúng vì $x = 5\sqrt{2}\,\mathrm{cm}= A/\sqrt{2}$, tương ứng với vị trí có động năng bằng thế năng
\item (Đúng) Tốc độ của vật tại vị trí có động năng bằng thế năng bằng $100\sqrt{2}$ cm/s. (Vì): Mệnh đề đúng vì khi $W_d = W_t \Rightarrow v = \dfrac{v_{\text{max}}}{\sqrt{2}} = \dfrac{\omega A}{\sqrt{2}} = \dfrac{20 \cdot 10}{\sqrt{2}} = 100\sqrt{2}$ cm/s
\item (Sai) Tốc độ của vật tại vị trí có động năng bằng ba lần thế năng bằng $100$ cm/s. (Vì): Mệnh đề này sai vì khi $W_d = 3W_t \Rightarrow v = \dfrac{v_{\text{max}}\sqrt{3}}{2} = \dfrac{\omega A\sqrt{3}}{2} = 100\sqrt{3}$ cm/s
\end{itemchoice}
}
\end{ex}

% ===== Câu 293 =====
% ID: L11C1B5-83-TN
% Mức: VD
\begin{ex}
% ID: L11C1B5-83-TN
% Mức: VD
Một con lắc đơn dao động điều hòa với biên độ góc $\alpha_0$. Khi vật đi qua vị trí có li độ góc $\alpha = \pm \dfrac{\alpha_0}{\sqrt{2}}$ thì tỉ số giữa động năng và thế năng của con lắc bằng:
\choice
{$2$}
{$\dfrac{1}{2}$}
{\True $1$}
{$3$}
\loigiai{
• Thế năng của con lắc đơn ở li độ góc $\alpha = \pm \dfrac{\alpha_0}{\sqrt{2}}$:
 $W_{\text{t}} = \dfrac{1}{2}mgl\alpha^2 = \dfrac{1}{2}mgl\left(\dfrac{\alpha_0}{\sqrt{2}}\right)^2 = \dfrac{1}{2} \cdot \left(\dfrac{1}{2}mgl\alpha_0^2\right) = \dfrac{W}{2}.$
 
• Động năng, theo bảo toàn cơ năng:
 $W_{\text{đ}} = W - W_{\text{t}} = W - \dfrac{W}{2} = \dfrac{W}{2}.$
 
• Lập tỉ số:
 $\dfrac{W_{\text{đ}}}{W_{\text{t}}} = \dfrac{\frac{W}{2}}{\frac{W}{2}} = 1.$
}
\end{ex}

% ===== Câu 294 =====
% ID: L11C1B5-84-DS
% Mức: VD
\begin{ex}
% ID: L11C1B5-84-DS
% Mức: VD
Một con lắc lò xo dao động điều hòa với biên độ $A = 10$ cm và tần số góc $\omega = 20$ rad/s. Khảo sát động học của vật tại các vị trí phân chia năng lượng:
\choiceTF
{\True Tại vị trí có li độ $x = 5$ cm, động năng của vật bằng ba lần thế năng.}
{\True Tại vị trí có li độ $x = 5\sqrt{2}$ cm, động năng của vật bằng thế năng của con lắc.}
{\True Tốc độ của vật tại vị trí có động năng bằng thế năng bằng $100\sqrt{2}$ cm/s.}
{Tốc độ của vật tại vị trí có động năng bằng ba lần thế năng bằng $100$ cm/s.}
\loigiai{
\begin{itemchoice}
\item (Đúng) Tại vị trí có li độ $x = 5$ cm, động năng của vật bằng ba lần thế năng. (Vì): Mệnh đề này đúng vì $x = 5\,\mathrm{cm}= A/2$, tương ứng với vị trí có động năng bằng 3 lần thế năng
\item (Đúng) Tại vị trí có li độ $x = 5\sqrt{2}$ cm, động năng của vật bằng thế năng của con lắc. (Vì): Mệnh đề này đúng vì $x = 5\sqrt{2}\,\mathrm{cm}= A/\sqrt{2}$, tương ứng với vị trí có động năng bằng thế năng
\item (Đúng) Tốc độ của vật tại vị trí có động năng bằng thế năng bằng $100\sqrt{2}$ cm/s. (Vì): Mệnh đề đúng vì khi $W_d = W_t \Rightarrow v = \dfrac{v_{\text{max}}}{\sqrt{2}} = \dfrac{\omega A}{\sqrt{2}} = \dfrac{20 \cdot 10}{\sqrt{2}} = 100\sqrt{2}$ cm/s
\item (Sai) Tốc độ của vật tại vị trí có động năng bằng ba lần thế năng bằng $100$ cm/s. (Vì): Mệnh đề này sai vì khi $W_d = 3W_t \Rightarrow v = \dfrac{v_{\text{max}}\sqrt{3}}{2} = \dfrac{\omega A\sqrt{3}}{2} = 100\sqrt{3}$ cm/s
\end{itemchoice}
}
\end{ex}

% ===== Câu 295 =====
% ID: L11C1B5-84-TN
% Mức: VD
\begin{ex}
% ID: L11C1B5-84-TN
% Mức: VD
Một chất điểm dao động điều hòa với biên độ $A = 6\ \mathrm{cm}$. Khoảng cách từ chất điểm đến vị trí cân bằng tại nơi có động năng bằng ba lần thế năng bằng:
\choice
{$3\sqrt{2}\ \mathrm{cm}$}
{\True $3\ \mathrm{cm}$}
{$1,5\ \mathrm{cm}$}
{$3\sqrt{3}\ \mathrm{cm}$}
\loigiai{
• Điều kiện đề bài: $W_{\text{đ}} = 3W_{\text{t}}$.
 
• Áp dụng bảo toàn cơ năng:
  \begin{aligned}
 $W$ &= W_{\text{đ}} + W_{ $\text{t}$ } = 4W_{ $\text{t}$ } 
$\Rightarrow\dfrac{1}{2}m\omega^2$ A^2 &= 4 $\cdot\dfrac{1}{2}m\omega^2$ x^2 
\Rightarrow x &= $\pm\dfrac{A}{2}$.
 \end{aligned} 
 
• Thay biên độ $A = 6\ \mathrm{cm}$:
 $d = |x| = \dfrac{6}{2} = 3\ \mathrm{cm}.$
}
\end{ex}

% ===== Câu 296 =====
% ID: L11C1B5-85-DS
% Mức: VD
\begin{ex}
% ID: L11C1B5-85-DS
% Mức: VD
Một con lắc lò xo dao động điều hòa với biên độ $A = 10$ cm và tần số góc $\omega = 20$ rad/s. Khảo sát động học của vật tại các vị trí phân chia năng lượng:
\choiceTF
{\True Tại vị trí có li độ $x = 5$ cm, động năng của vật bằng ba lần thế năng.}
{\True Tại vị trí có li độ $x = 5\sqrt{2}$ cm, động năng của vật bằng thế năng của con lắc.}
{\True Tốc độ của vật tại vị trí có động năng bằng thế năng bằng $100\sqrt{2}$ cm/s.}
{Tốc độ của vật tại vị trí có động năng bằng ba lần thế năng bằng $100$ cm/s.}
\loigiai{
\begin{itemchoice}
\item (Đúng) Tại vị trí có li độ $x = 5$ cm, động năng của vật bằng ba lần thế năng. (Vì): Mệnh đề này đúng vì $x = 5\,\mathrm{cm}= A/2$, tương ứng với vị trí có động năng bằng 3 lần thế năng
\item (Đúng) Tại vị trí có li độ $x = 5\sqrt{2}$ cm, động năng của vật bằng thế năng của con lắc. (Vì): Mệnh đề này đúng vì $x = 5\sqrt{2}\,\mathrm{cm}= A/\sqrt{2}$, tương ứng với vị trí có động năng bằng thế năng
\item (Đúng) Tốc độ của vật tại vị trí có động năng bằng thế năng bằng $100\sqrt{2}$ cm/s. (Vì): Mệnh đề đúng vì khi $W_d = W_t \Rightarrow v = \dfrac{v_{\text{max}}}{\sqrt{2}} = \dfrac{\omega A}{\sqrt{2}} = \dfrac{20 \cdot 10}{\sqrt{2}} = 100\sqrt{2}$ cm/s
\item (Sai) Tốc độ của vật tại vị trí có động năng bằng ba lần thế năng bằng $100$ cm/s. (Vì): Mệnh đề này sai vì khi $W_d = 3W_t \Rightarrow v = \dfrac{v_{\text{max}}\sqrt{3}}{2} = \dfrac{\omega A\sqrt{3}}{2} = 100\sqrt{3}$ cm/s
\end{itemchoice}
}
\end{ex}

% ===== Câu 297 =====
% ID: L11C1B5-85-TN
% Mức: VD
\begin{ex}
% ID: L11C1B5-85-TN
% Mức: VD
Một chất điểm dao động điều hòa với biên độ $A = 6\ \mathrm{cm}$. Khoảng cách từ chất điểm đến vị trí cân bằng tại nơi có động năng bằng ba lần thế năng bằng:
\choice
{$3\sqrt{2}\ \mathrm{cm}$}
{\True $3\ \mathrm{cm}$}
{$1,5\ \mathrm{cm}$}
{$3\sqrt{3}\ \mathrm{cm}$}
\loigiai{
• Điều kiện đề bài: $W_{\text{đ}} = 3W_{\text{t}}$.
 
• Áp dụng bảo toàn cơ năng:
  \begin{aligned}
 $W$ &= W_{\text{đ}} + W_{ $\text{t}$ } = 4W_{ $\text{t}$ } 
$\Rightarrow\dfrac{1}{2}m\omega^2$ A^2 &= 4 $\cdot\dfrac{1}{2}m\omega^2$ x^2 
\Rightarrow x &= $\pm\dfrac{A}{2}$.
 \end{aligned} 
 
• Thay biên độ $A = 6\ \mathrm{cm}$:
 $d = |x| = \dfrac{6}{2} = 3\ \mathrm{cm}.$
}
\end{ex}

% ===== Câu 298 =====
% ID: L11C1B5-86-DS
% Mức: TH
\dangbt{Xác định độ cứng lò xo từ chu kì biến thiên năng lượng}
\begin{ex}
% ID: L11C1B5-86-DS
% Mức: TH
Con lắc lò xo có cơ năng 0.04 $J$ và biên độ $4\,\text{cm}$. Xác định tính đúng sai.
\choiceTF
{\True Cơ năng được tính bởi $W=kA^2/2$.}
{\True Độ cứng lò xo bằng $50\,\text{N}$ /m.}
{\True Tại biên, động năng bằng 0.}
{Tại vị trí cân bằng, thế năng cực đại.}
\end{ex}

% ===== Câu 299 =====
% ID: L11C1B5-86-TN
% Mức: VD
\dangbt{Xác định vị trí dựa vào tỉ số động năng và thế năng}
\begin{ex}
% ID: L11C1B5-86-TN
% Mức: VD
Một chất điểm dao động điều hòa với biên độ $A = 6\ \mathrm{cm}$. Khoảng cách từ chất điểm đến vị trí cân bằng tại nơi có động năng bằng ba lần thế năng bằng:
\choice
{$3\sqrt{2}\ \mathrm{cm}$}
{\True $3\ \mathrm{cm}$}
{$1,5\ \mathrm{cm}$}
{$3\sqrt{3}\ \mathrm{cm}$}
\loigiai{
• Điều kiện đề bài: $W_{\text{đ}} = 3W_{\text{t}}$.
 
• Áp dụng bảo toàn cơ năng:
  \begin{aligned}
 $W$ &= W_{\text{đ}} + W_{ $\text{t}$ } = 4W_{ $\text{t}$ } 
$\Rightarrow\dfrac{1}{2}m\omega^2$ A^2 &= 4 $\cdot\dfrac{1}{2}m\omega^2$ x^2 
\Rightarrow x &= $\pm\dfrac{A}{2}$.
 \end{aligned} 
 
• Thay biên độ $A = 6\ \mathrm{cm}$:
 $d = |x| = \dfrac{6}{2} = 3\ \mathrm{cm}.$
}
\end{ex}

% ===== Câu 300 =====
% ID: L11C1B5-87-DS
% Mức: VD
\dangbt{Xác định độ cứng lò xo từ chu kì biến thiên năng lượng}
\begin{ex}
% ID: L11C1B5-87-DS
% Mức: VD
Con lắc lò xo có cơ năng 0.09 $J$ và biên độ $6\,\text{cm}$. Xác định tính đúng sai.
\choiceTF
{\True Cơ năng được tính bởi $W=kA^2/2$.}
{\True Độ cứng lò xo bằng $50\,\text{N}$ /m.}
{\True Tại biên, động năng bằng 0.}
{Tại vị trí cân bằng, thế năng cực đại.}
\end{ex}

% ===== Câu 301 =====
% ID: L11C1B5-87-TN
% Mức: VD
\dangbt{Xác định vị trí dựa vào tỉ số động năng và thế năng}
\begin{ex}
% ID: L11C1B5-87-TN
% Mức: VD
Một con lắc lò xo dao động điều hòa với biên độ $A = 8\ \mathrm{cm}$. Li độ của vật tại nơi thế năng của con lắc bằng ba lần động năng bằng:
\choice
{$x = \pm 4\ \mathrm{cm}$}
{$x = \pm 4\sqrt{2}\ \mathrm{cm}$}
{\True $x = \pm 4\sqrt{3}\ \mathrm{cm}$}
{$x = \pm 2\ \mathrm{cm}$}
\loigiai{
• Điều kiện đề bài: $W_{\text{t}} = 3W_{\text{đ}} \Rightarrow W_{\text{đ}} = \dfrac{1}{3}W_{\text{t}}$.
 
• Áp dụng bảo toàn cơ năng:
  \begin{aligned}
 $W$ &= W_{\text{đ}} + W_{ $\text{t}$ } = $\dfrac{4}{3}W_{\text{t}}\Rightarrow\dfrac{1}{2}kA^2$ &= $\dfrac{4}{3}\cdot\dfrac{1}{2}kx^2\Rightarrow$ x &= $\pm\dfrac{A\sqrt{3}}{2}$.
 \end{aligned} 
 
• Thay biên độ $A = 8\ \mathrm{cm}$:
 $x = \pm \dfrac{8\sqrt{3}}{2} = \pm 4\sqrt{3}\ \mathrm{cm}.$
}
\end{ex}

% ===== Câu 302 =====
% ID: L11C1B5-88-DS
% Mức: VD
\dangbt{Xác định độ cứng lò xo từ chu kì biến thiên năng lượng}
\begin{ex}
% ID: L11C1B5-88-DS
% Mức: VD
Con lắc lò xo có cơ năng 0.2 $J$ và biên độ $10\,\text{cm}$. Xác định tính đúng sai.
\choiceTF
{\True Cơ năng được tính bởi $W=kA^2/2$.}
{\True Độ cứng lò xo bằng $39.99999999999999\,\text{N}$ /m.}
{\True Tại biên, động năng bằng 0.}
{Tại vị trí cân bằng, thế năng cực đại.}
\end{ex}

% ===== Câu 303 =====
% ID: L11C1B5-88-TN
% Mức: VD
\dangbt{Xác định vị trí dựa vào tỉ số động năng và thế năng}
\begin{ex}
% ID: L11C1B5-88-TN
% Mức: VD
Một con lắc lò xo dao động điều hòa với biên độ $A = 8\ \mathrm{cm}$. Li độ của vật tại nơi thế năng của con lắc bằng ba lần động năng bằng:
\choice
{$x = \pm 4\ \mathrm{cm}$}
{$x = \pm 4\sqrt{2}\ \mathrm{cm}$}
{\True $x = \pm 4\sqrt{3}\ \mathrm{cm}$}
{$x = \pm 2\ \mathrm{cm}$}
\loigiai{
• Điều kiện đề bài: $W_{\text{t}} = 3W_{\text{đ}} \Rightarrow W_{\text{đ}} = \dfrac{1}{3}W_{\text{t}}$.
 
• Áp dụng bảo toàn cơ năng:
  \begin{aligned}
 $W$ &= W_{\text{đ}} + W_{ $\text{t}$ } = $\dfrac{4}{3}W_{\text{t}}\Rightarrow\dfrac{1}{2}kA^2$ &= $\dfrac{4}{3}\cdot\dfrac{1}{2}kx^2\Rightarrow$ x &= $\pm\dfrac{A\sqrt{3}}{2}$.
 \end{aligned} 
 
• Thay biên độ $A = 8\ \mathrm{cm}$:
 $x = \pm \dfrac{8\sqrt{3}}{2} = \pm 4\sqrt{3}\ \mathrm{cm}.$
}
\end{ex}

% ===== Câu 304 =====
% ID: L11C1B5-89-DS
% Mức: TH
\dangbt{Xác định độ cứng lò xo từ cơ năng và thông số dao động}
\begin{ex}
% ID: L11C1B5-89-DS
% Mức: TH
Con lắc lò xo có cơ năng 0.04 $J$ và biên độ $4\,\text{cm}$. Xác định tính đúng sai.
\choiceTF
{\True Cơ năng được tính bởi $W=kA^2/2$.}
{\True Độ cứng lò xo bằng $50\,\text{N}$ /m.}
{\True Tại biên, động năng bằng 0.}
{Tại vị trí cân bằng, thế năng cực đại.}
\end{ex}

% ===== Câu 305 =====
% ID: L11C1B5-89-TN
% Mức: VD
\dangbt{Xác định vị trí dựa vào tỉ số động năng và thế năng}
\begin{ex}
% ID: L11C1B5-89-TN
% Mức: VD
Một con lắc lò xo dao động điều hòa với biên độ $A = 8\ \mathrm{cm}$. Li độ của vật tại nơi thế năng của con lắc bằng ba lần động năng bằng:
\choice
{$x = \pm 4\ \mathrm{cm}$}
{$x = \pm 4\sqrt{2}\ \mathrm{cm}$}
{\True $x = \pm 4\sqrt{3}\ \mathrm{cm}$}
{$x = \pm 2\ \mathrm{cm}$}
\loigiai{
• Điều kiện đề bài: $W_{\text{t}} = 3W_{\text{đ}} \Rightarrow W_{\text{đ}} = \dfrac{1}{3}W_{\text{t}}$.
 
• Áp dụng bảo toàn cơ năng:
  \begin{aligned}
 $W$ &= W_{\text{đ}} + W_{ $\text{t}$ } = $\dfrac{4}{3}W_{\text{t}}\Rightarrow\dfrac{1}{2}kA^2$ &= $\dfrac{4}{3}\cdot\dfrac{1}{2}kx^2\Rightarrow$ x &= $\pm\dfrac{A\sqrt{3}}{2}$.
 \end{aligned} 
 
• Thay biên độ $A = 8\ \mathrm{cm}$:
 $x = \pm \dfrac{8\sqrt{3}}{2} = \pm 4\sqrt{3}\ \mathrm{cm}.$
}
\end{ex}

% ===== Câu 306 =====
% ID: L11C1B5-90-DS
% Mức: VD
\dangbt{Xác định độ cứng lò xo từ cơ năng và thông số dao động}
\begin{ex}
% ID: L11C1B5-90-DS
% Mức: VD
Con lắc lò xo có cơ năng 0.09 $J$ và biên độ $6\,\text{cm}$. Xác định tính đúng sai.
\choiceTF
{\True Cơ năng được tính bởi $W=kA^2/2$.}
{\True Độ cứng lò xo bằng $50\,\text{N}$ /m.}
{\True Tại biên, động năng bằng 0.}
{Tại vị trí cân bằng, thế năng cực đại.}
\end{ex}

% ===== Câu 307 =====
% ID: L11C1B5-90-TN
% Mức: VD
\dangbt{Xác định vị trí dựa vào tỉ số động năng và thế năng}
\begin{ex}
% ID: L11C1B5-90-TN
% Mức: VD
Một vật dao động điều hòa với chu kì $T$ và biên độ $A$. Khi vật đi qua vị trí có li độ $x = \dfrac{A}{2}$ thì tỉ số giữa thế năng và động năng của vật bằng:
\choice
{$3$}
{\True $\dfrac{1}{3}$}
{$\dfrac{1}{2}$}
{$1$}
\loigiai{
• Thế năng tại li độ $x = \dfrac{A}{2}$:
 $W_{\text{t}} = \dfrac{1}{2}m\omega^2 \left(\dfrac{A}{2}\right)^2 = \dfrac{1}{4}W.$
 
• Động năng tương ứng, theo bảo toàn cơ năng:
 $W_{\text{đ}} = W - W_{\text{t}} = W - \dfrac{1}{4}W = \dfrac{3}{4}W.$
 
• Lập tỉ số:
 $\dfrac{W_{\text{t}}}{W_{\text{đ}}} = \dfrac{\frac{1}{4}W}{\frac{3}{4}W} = \dfrac{1}{3}.$
}
\end{ex}

% ===== Câu 308 =====
% ID: L11C1B5-91-DS
% Mức: VD
\dangbt{Xác định độ cứng lò xo từ cơ năng và thông số dao động}
\begin{ex}
% ID: L11C1B5-91-DS
% Mức: VD
Con lắc lò xo có cơ năng 0.2 $J$ và biên độ $10\,\text{cm}$. Xác định tính đúng sai.
\choiceTF
{\True Cơ năng được tính bởi $W=kA^2/2$.}
{\True Độ cứng lò xo bằng $39.99999999999999\,\text{N}$ /m.}
{\True Tại biên, động năng bằng 0.}
{Tại vị trí cân bằng, thế năng cực đại.}
\end{ex}

% ===== Câu 309 =====
% ID: L11C1B5-91-TN
% Mức: VD
\dangbt{Xác định vị trí dựa vào tỉ số động năng và thế năng}
\begin{ex}
% ID: L11C1B5-91-TN
% Mức: VD
Một vật dao động điều hòa với chu kì $T$ và biên độ $A$. Khi vật đi qua vị trí có li độ $x = \dfrac{A}{2}$ thì tỉ số giữa thế năng và động năng của vật bằng:
\choice
{$3$}
{\True $\dfrac{1}{3}$}
{$\dfrac{1}{2}$}
{$1$}
\loigiai{
• Thế năng tại li độ $x = \dfrac{A}{2}$:
 $W_{\text{t}} = \dfrac{1}{2}m\omega^2 \left(\dfrac{A}{2}\right)^2 = \dfrac{1}{4}W.$
 
• Động năng tương ứng, theo bảo toàn cơ năng:
 $W_{\text{đ}} = W - W_{\text{t}} = W - \dfrac{1}{4}W = \dfrac{3}{4}W.$
 
• Lập tỉ số:
 $\dfrac{W_{\text{t}}}{W_{\text{đ}}} = \dfrac{\frac{1}{4}W}{\frac{3}{4}W} = \dfrac{1}{3}.$
}
\end{ex}

% ===== Câu 310 =====
% ID: L11C1B5-92-DS
% Mức: TH
\dangbt{Xác định độ cứng lò xo từ năng lượng ban đầu}
\begin{ex}
% ID: L11C1B5-92-DS
% Mức: TH
Cho đồ thị biểu diễn mối quan hệ giữa động năng, thế năng và cơ năng trong dao động điều hòa. Chọn phát biểu đúng.
\choiceTF
{Động năng của vật dao động điều hòa luôn luôn tăng khi vật chuyển động từ vị trí cân bằng ra biên}
{Thế năng của vật dao động điều hòa luôn luôn giảm khi vật chuyển động từ vị trí cân bằng ra biên}
{\True Cơ năng của vật dao động điều hòa được bảo toàn và luôn bằng tổng động năng và thế năng tại mọi vị trí}
{Tại vị trí biên, động năng của vật đạt giá trị cực đại và bằng cơ năng.}
\loigiai{
\begin{itemchoice}
\item (Sai) Động năng của vật dao động điều hòa luôn luôn tăng khi vật chuyển động từ vị trí cân bằng ra biên. (Vì): Đáp án 1: Sai. Khi vật đi từ vị trí cân bằng ra biên thì vận tốc giảm động năng giảm
\item (Sai) Thế năng của vật dao động điều hòa luôn luôn giảm khi vật chuyển động từ vị trí cân bằng ra biên. (Vì): Đáp án 2: Sai. Thế năng tăng dần khi ra biên do $x$ tăng $W_t \propto x^2$
\item (Đúng) Cơ năng của vật dao động điều hòa được bảo toàn và luôn bằng tổng động năng và thế năng tại mọi vị trí. (Vì): áp án 3: Đúng. Cơ năng luôn là tổng hai năng lượng và không đổi
\item (Sai) Tại vị trí biên, động năng của vật đạt giá trị cực đại và bằng cơ năng. (Vì): Đáp án 4: Sai. Tại biên, vận tốc bằng 0 động năng bằng 0, thế năng cực đại = cơ năng
\end{itemchoice}
}
\end{ex}

% ===== Câu 311 =====
% ID: L11C1B5-92-TN
% Mức: VD
\dangbt{Xác định vị trí dựa vào tỉ số động năng và thế năng}
\begin{ex}
% ID: L11C1B5-92-TN
% Mức: VD
Một vật dao động điều hòa với chu kì $T$ và biên độ $A$. Khi vật đi qua vị trí có li độ $x = \dfrac{A}{2}$ thì tỉ số giữa thế năng và động năng của vật bằng:
\choice
{$3$}
{\True $\dfrac{1}{3}$}
{$\dfrac{1}{2}$}
{$1$}
\loigiai{
• Thế năng tại li độ $x = \dfrac{A}{2}$:
 $W_{\text{t}} = \dfrac{1}{2}m\omega^2 \left(\dfrac{A}{2}\right)^2 = \dfrac{1}{4}W.$
 
• Động năng tương ứng, theo bảo toàn cơ năng:
 $W_{\text{đ}} = W - W_{\text{t}} = W - \dfrac{1}{4}W = \dfrac{3}{4}W.$
 
• Lập tỉ số:
 $\dfrac{W_{\text{t}}}{W_{\text{đ}}} = \dfrac{\frac{1}{4}W}{\frac{3}{4}W} = \dfrac{1}{3}.$
}
\end{ex}

% ===== Câu 312 =====
% ID: L11C1B5-93-DS
% Mức: TH
\dangbt{Xác định độ cứng lò xo từ năng lượng ban đầu}
\begin{ex}
% ID: L11C1B5-93-DS
% Mức: TH
Cho đồ thị như hình vẽ mô tả biến thiên năng lượng của vật dao động điều hòa theo thời gian. Chọn phát biểu đúng.
\choiceTF
{\True Cơ năng luôn không đổi và là đường song song với trục thời gian $t$}
{Động năng và thế năng luôn biến thiên cùng chu kỳ và cùng pha}
{Trong 1 chu kỳ luôn có 2 lần động năng bằng thế năng}
{Khi động năng đạt cực đại thì thế năng sẽ bằng cơ năng và ngược lại}
\loigiai{
\begin{itemchoice}
\item (Đúng) Cơ năng luôn không đổi và là đường song song với trục thời gian $t$. (Vì): áp án 1: Đúng. Vì $W = W_d + W_t$ luôn không đổi nên là đường thẳng song song với trục $t$
\item (Đúng) Động năng và thế năng luôn biến thiên cùng chu kỳ và cùng pha. (Vì): áp án 2: Sai. Động năng và thế năng có cùng chu kỳ nhưng ngược pha nhau
\item (Đúng) Trong 1 chu kỳ luôn có 2 lần động năng bằng thế năng. (Vì): áp án 3: Sai. Mỗi chu kỳ có 4 lần động năng bằng thế năng, không phải 2
\item (Đúng) Khi động năng đạt cực đại thì thế năng sẽ bằng cơ năng và ngược lại. (Vì): áp án 4: Sai. Khi $W_d$ cực đại thì $W_t = 0$, không thể bằng cơ năng
\end{itemchoice}
}
\end{ex}

% ===== Câu 313 =====
% ID: L11C1B5-93-TN
% Mức: TH
\dangbt{Xác định độ cứng lò xo từ chu kì biến thiên năng lượng}
\begin{ex}
% ID: L11C1B5-93-TN
% Mức: TH
Một vật dao động điều hoà có li độ x có chu kì $T$, động năng thời gian là hàm tuần hoàn có chu kì là
\choice
{$T$}
{\True $\dfrac{T}{2}$}
{2 $T$}
{$\dfrac{T}{4}$}
\loigiai{
$({\text{W)}_{d}}=\dfrac{1}{2}mv^{2}=\dfrac{1}{2}m{\omega}^{2}A^{2}{\sin}^{2}(\omega t+\varphi )=\dfrac{1}{4}m{\omega}^{2}A^{2}-\dfrac{1}{4}m{\omega}^{2}A^{2}\cos$ 2 $\omega t+\varphi )$ Động năng là hàm điều hoà có tần số góc $\omega$ '=2 $\omega\Rightarrow$ T'= $\dfrac{T}{2}$
}
\end{ex}

% ===== Câu 314 =====
% ID: L11C1B5-94-DS
% Mức: TH
\dangbt{Xác định độ cứng lò xo từ năng lượng ban đầu}
\begin{ex}
% ID: L11C1B5-94-DS
% Mức: TH
Cho đồ thị sau biểu diễn động năng $W_d$ và thế năng $W_t$ của một vật dao động điều hòa:
\choiceTF
{\True Tổng $W_t+W_d$ luôn không đổi theo thời gian}
{Động năng cực đại và thế năng cực đại xảy ra cùng lúc}
{\True Thế năng có chu kỳ bằng $\dfrac{T}{2}$}
{Tại $t=\frac{T}{6}$, động năng bằng 0}
\loigiai{
\begin{itemchoice}
\item (Đúng) Tổng $W_t+W_d$ luôn không đổi theo thời gian. (Vì): Trong dao động điều hòa lý tưởng, cơ năng $W=W_d+W_t$ luôn bảo toàn \Rightarrow \textbf{Đúng}
\item (Sai) Động năng cực đại và thế năng cực đại xảy ra cùng lúc. (Vì): $W_t$ và $W_d$ biến thiên theo quy luật ngược pha (lệch pha $\pi$), nên không thể cực đại cùng lúc $\Rightarrow\textbf{Sai}$
\item (Đúng) Thế năng có chu kỳ bằng $\dfrac{T}{2}$. (Vì): Vì $W_t \sim \cos^2(\omega t)$ nên chu kỳ của $W_t$ là $T/2\Rightarrow \textbf{Đúng}
\item (Sai) Tại $t=\frac{T}{6}$, động năng bằng 0. (Vì): Tại$t=$\frac{T}{6}$, theo đồ thị thì động năng không đạt cực đại, không bằng 0$\Rightarrow\textbf{Sai}$
\end{itemchoice}
}
\end{ex}

% ===== Câu 315 =====
% ID: L11C1B5-94-TN
% Mức: TH
\dangbt{Xác định độ cứng lò xo từ chu kì biến thiên năng lượng}
\begin{ex}
% ID: L11C1B5-94-TN
% Mức: TH
Một vật dao động điều hoà có li độ x có tần số f, thế năng theo thời gian là hàm tuần hoàn có tần số là
\choice
{$\dfrac{f}{4}$}
{$\dfrac{f}{2}$}
{f}
{\True 2f}
\loigiai{
$({\text{W)}_{t}}=\dfrac{1}{2}m{\omega}^{2}x^{2}=\dfrac{1}{2}m{\omega}^{2}A^{2}{\cos}^{2}(\omega t+\varphi )=\dfrac{1}{4}m{\omega}^{2}A^{2}+\dfrac{1}{4}m{\omega}^{2}A^{2}\cos$ 2 $\omega t+\varphi )$ Thế năng là hàm điều hoà có tần số góc $\omega$ '=2 $\omega\Rightarrow$ f'=2f
}
\end{ex}

% ===== Câu 316 =====
% ID: L11C1B5-95-TN
% Mức: TH
\begin{ex}
% ID: L11C1B5-95-TN
% Mức: TH
Đối với một chất điểm dao động điều hòa với chu kì $T$ thì
\choice
{động năng và thế năng đều biến thiên tuần hoàn theo thời gian với chu kì $T$}
{động năng và thế năng đều biến thiên tuần hoàn theo thời gian nhưng không điều hòa}
{\True động năng và thế năng đều biến thiên tuần hoàn theo thời gian với chu kì T/2}
{động năng và thế năng đều biến thiên tuần hoàn theo thời gian với chu kì 2 $T$}
\loigiai{
Trong dao động cơ điều hòa, động năng và thế năng biến thiên tuần hoàn với chu kỳ bằng nửa chu kỳ dao động.
}
\end{ex}

% ===== Câu 317 =====
% ID: L11C1B5-96-TN
% Mức: VD
\begin{ex}
% ID: L11C1B5-96-TN
% Mức: VD
Một vật nhỏ thực hiện dao động điều hòa theo phương trình $x=10\cos(4\pi t)$ cm. Động năng của vật đó biến thiên với chu kì bằng:
\choice
{$0,5s$}
{$1s$}
{\True $0,25s$}
{$2s$}
\loigiai{
Phương trình dao động điều hòa của vật là $x=10\cos(4\pi t)$ cm.Từ phương trình, ta có tần số góc của dao động là $\omega = 4\pi$ rad/s.Chu kì dao động của vật là $T = \frac{2\pi}{\omega} = \frac{2\pi}{4\pi} = 0,5$ s.Trong dao động điều hòa, động năng và thế năng biến thiên tuần hoàn với tần số góc $2\omega$ và chu kì $T' = \frac{T}{2}$.Vậy chu kì biến thiên của động năng là $T' = \frac{T}{2} = \frac{0,5}{2} = 0,25$ s.
}
\end{ex}

% ===== Câu 318 =====
% ID: L11C1B5-97-TN
% Mức: TH
\dangbt{Xác định độ cứng lò xo từ cơ năng và thông số dao động}
\begin{ex}
% ID: L11C1B5-97-TN
% Mức: TH
Con lắc lò xo có cơ năng 0.04 $J$ và biên độ $4\,\text{cm}$. Độ cứng lò xo bằng
\choice
{\True $50\,\text{N}$ /m}
{$25\,\text{N}$ /m}
{$100\,\text{N}$ /m}
{$12.5\,\text{N}$ /m}
\end{ex}

% ===== Câu 319 =====
% ID: L11C1B5-98-TN
% Mức: VD
\begin{ex}
% ID: L11C1B5-98-TN
% Mức: VD
Con lắc lò xo có cơ năng 0.09 $J$ và biên độ $6\,\text{cm}$. Độ cứng lò xo bằng
\choice
{\True $50\,\text{N}$ /m}
{$25\,\text{N}$ /m}
{$100\,\text{N}$ /m}
{$12.5\,\text{N}$ /m}
\end{ex}

% ===== Câu 320 =====
% ID: L11C1B5-99-TN
% Mức: VD
\begin{ex}
% ID: L11C1B5-99-TN
% Mức: VD
Con lắc lò xo có cơ năng 0.2 $J$ và biên độ $10\,\text{cm}$. Độ cứng lò xo bằng
\choice
{\True $39.99999999999999\,\text{N}$ /m}
{$19.999999999999996\,\text{N}$ /m}
{$79.99999999999999\,\text{N}$ /m}
{$9.999999999999998\,\text{N}$ /m}
\end{ex}
